% Leçon 38 — Vocabulaire et compréhension de texte : rédaction d’un CV — Chinois Tle A — Cours signé OnBuch+
% Programme officiel MINESEC. Compiler avec : tectonic ZH38-vocabulaire-redaction-dun-cv.tex
\newcommand{\DOCMATIERE}{Chinois}
\newcommand{\DOCNIVEAU}{Tle A}
\newcommand{\DOCMODULE}{Module 5 : Mass médias et télécommunication}
\newcommand{\DOCLECON}{ZH38}
\newcommand{\DOCTITRE}{Leçon 38 — Vocabulaire et compréhension de texte : rédaction d’un CV}
% OnBuch+ — préambule « cours signés » ENRICHI (compilé avec tectonic / XeLaTeX)
% Système visuel « Pop » d'OnBuch+ : fond crème (jamais blanc), bordures encre
% épaisses, ombres « dures » décalées, titres Archivo Black en capitales.
% Version MATHÉMATIQUES : graphes (pgfplots/tikz), boîtes pédagogiques
% (définition, propriété, méthode, attention, le savais-tu, exercice résolu,
% exercice, TP, bilan), page de garde et signature OnBuch+ ancrée en bas.
%
% Macros attendues dans main.tex AVANT \input{preamble} :
%   \DOCMATIERE  \DOCNIVEAU  \DOCMODULE  \DOCLECON (numéro)  \DOCTITRE
\documentclass[11pt]{article}

\usepackage{fontspec}
\usepackage[french]{babel} % français = langue principale ; le chinois via \zh{..} / textebox (xeCJK)
\usepackage{xeCJK}
\usepackage{pifont} % \ding{51} \ding{55} utilisés par les modèles
\usepackage[a4paper,margin=2cm,top=3.4cm,bottom=2.8cm,headheight=1.4cm,footskip=1.3cm]{geometry}
\usepackage{amsmath,amssymb,amsfonts}
\usepackage{mathtools} % \xmapsto, \coloneqq... souvent utilisés par les modèles
\usepackage{cancel} % simplifications barrées (bilans de forces)
% \abs{z} (module, valeur absolue) et \norm{u} (norme), utilisés par les modèles
\providecommand{\abs}{}\let\abs\relax\DeclarePairedDelimiter{\abs}{\lvert}{\rvert}
\providecommand{\norm}{}\let\norm\relax\DeclarePairedDelimiter{\norm}{\lVert}{\rVert}
\usepackage[table]{xcolor} % [table] : fournit aussi \rowcolors (alternance de lignes)
\usepackage{pagecolor}
\usepackage{tikz}
\usetikzlibrary{arrows.meta,positioning,calc,shapes.geometric,shapes.misc,shapes.arrows,decorations.pathmorphing,decorations.pathreplacing,decorations.markings,patterns,shadows,mindmap,trees,fit,backgrounds,matrix,intersections,angles,quotes}
\usepackage{pgfplots}
\usepackage[version=4]{mhchem} % équations nucléaires \ce{^{235}_{92}U}
\usepackage[european,siunitx]{circuitikz} % schémas électriques
\pgfplotsset{compat=1.18}
\usepgfplotslibrary{fillbetween,groupplots} % aires entre courbes (calcul intégral)
\usepackage{enumitem}
\usepackage{fancyhdr}
% [most] charge toutes les bibliothèques tcolorbox (listings, minted, raster,
% documentation...) et gonfle inutilement la mémoire TeX sur un document long
% (~300 boîtes) : on ne charge que ce qui est réellement utilisé.
\usepackage[skins,breakable]{tcolorbox}
\usepackage{graphicx}
\usepackage{colortbl}
\usepackage{booktabs}
\usepackage{tabularx}
\usepackage{diagbox} % case d'en-tête diagonale des tableaux à double entrée
% Colonnes extensibles centrée / gauche / droite (C, L, R), souvent utilisées par les modèles
\newcolumntype{C}{>{\centering\arraybackslash}X}
\newcolumntype{L}{>{\raggedright\arraybackslash}X}
\newcolumntype{R}{>{\raggedleft\arraybackslash}X}
\usepackage{array}
\usepackage{multicol}
\usepackage{multirow}
\usepackage{siunitx}
% parse-numbers=false : accepte aussi \SI{2,5 \times 10^{-3}}{...}
\sisetup{locale=FR,output-decimal-marker={,},group-minimum-digits=4,parse-numbers=false}
\DeclareSIUnit\ohm{\ensuremath{\symup{\Omega}}} % U+2126 absent de la police
\DeclareSIUnit\division{div} % réglages d'oscilloscope (ms/div, V/div)
\DeclareSIUnit\tr{tr} % tours (vitesse de rotation, ex. \SI{1500}{\tr\per\minute})
\DeclareSIUnit\molar{M} % concentration molaire (ex. \SI{3}{\molar}, \SI{1}{\micro\molar})
\DeclareSIUnit\an{an} % année (le modèle confond parfois avec \year, un compteur TeX) — ex. \SI{5}{\kilo\meter\per\an}
\DeclareSIUnit\FCFA{FCFA} % franc CFA (ex. \SI{500}{\FCFA\per\kilo\gram})
\DeclareSIUnit\degreeBrix{\textdegree Brix} % teneur en sucre d'un sirop/jus (réfractométrie)
\DeclareSIUnit\month{mois} % le modèle utilise parfois \month, un compteur TeX, comme unité de durée
\usepackage{qrcode}
% \degree : commande fréquemment utilisée par les modèles pour un angle ou une
% température en dehors de \SI{}{} (non fournie par les paquets chargés).
\providecommand{\degree}{^{\circ}}
% \qed (fin de démonstration), utilisé par les modèles sans amsthm
\providecommand{\qed}{\hfill$\square$}
% \barre : double barre des valeurs interdites dans un tableau de variations (tkz-tab)
\providecommand{\barre}{\ensuremath{\|}}
% environnement proof (amsthm non chargé) utilisé par les modèles
\ifdefined\proof\else\newenvironment{proof}[1][Démonstration]{\par\noindent\textit{#1.}\ }{\qed\par}\fi
\usepackage{lastpage}
\usepackage{hyperref}
\hypersetup{hidelinks}

% ============================================================
% Palette « Pop » OnBuch+ — mêmes hex que la classe P (plus_ui.dart)
% ============================================================
\definecolor{poporange}{HTML}{F59321}
\definecolor{popdark}{HTML}{E07A0C}
\definecolor{popcream}{HTML}{FFF6E8}
\definecolor{popink}{HTML}{1C1714}
\definecolor{poppurple}{HTML}{7A5AE0}
\definecolor{popgreen}{HTML}{2E9E6A}
\definecolor{poppink}{HTML}{E0457B}
\definecolor{popblue}{HTML}{2F7DE1}
\definecolor{popgold}{HTML}{C98A12}
\definecolor{popmuted}{HTML}{8A7F76}
\definecolor{popline}{HTML}{EADFD2}
\definecolor{popgray}{HTML}{8A7F76}
\colorlet{popgrey}{popgray}
% Alias tolérants (le modèle nomme parfois les couleurs différemment)
\colorlet{popwhite}{white}
\colorlet{popblack}{popink}
\colorlet{popred}{poppink}
\colorlet{popyellow}{popgold}
% Teintes claires (fonds de tableaux/encadrés) — le modèle nomme parfois
% les couleurs avec un suffixe L (ex. popblueL) sans les avoir définies.
\colorlet{poporangeL}{poporange!20}
\colorlet{popdarkL}{popdark!20}
\colorlet{poppurpleL}{poppurple!20}
\colorlet{popgreenL}{popgreen!20}
\colorlet{poppinkL}{poppink!20}
\colorlet{popblueL}{popblue!20}
\colorlet{popgoldL}{popgold!20}
\colorlet{popredL}{popred!20}
\colorlet{popgrayL}{popgray!20}
% Teintes claires pour fonds de tableaux / zones
\colorlet{popblueL}{popblue!10!white}
\colorlet{poporangeL}{poporange!14!white}
\colorlet{popgreenL}{popgreen!12!white}
\colorlet{poppurpleL}{poppurple!12!white}
\colorlet{poppinkL}{poppink!12!white}
\colorlet{popgoldL}{popgold!14!white}

\pagecolor{popcream}

% ============================================================
% Typographie
% ============================================================
\defaultfontfeatures{Ligatures=TeX}
\setmainfont{PlusJakartaSans-Medium.ttf}[Path=../../../fonts/,BoldFont=PlusJakartaSans-Bold.ttf]
% Symboles, API et emojis absents de la police principale : repli sur DejaVu Sans (caractères actifs)
\newfontface\symfont{DejaVuSans.ttf}[Path=../../../fonts/]
\makeatletter
\newcommand\symactive[1]{\catcode#1=13 \begingroup\lccode`\~=#1 \lowercase{\endgroup\def~}{{\symfont\char#1}}}
\newcommand\symblank[1]{\catcode#1=13 \begingroup\lccode`\~=#1 \lowercase{\endgroup\def~}{}}
\newcommand\symhyph[1]{\catcode#1=13 \begingroup\lccode`\~=#1 \lowercase{\endgroup\def~}{\mbox{-}}}
\newcount\symc
\newcommand\symrange[2]{\symc=#1 \@whilenum\symc<#2 \do{\expandafter\symactive\expandafter{\the\symc}\advance\symc by 1 }}
\newcommand\symblankrange[2]{\symc=#1 \@whilenum\symc<#2 \do{\expandafter\symblank\expandafter{\the\symc}\advance\symc by 1 }}
\makeatother
\symrange{"0250}{"02FF}\symactive{"2713}\symactive{"2714}\symactive{"2717}\symactive{"2718}\symactive{"2610}\symactive{"2611}\symactive{"2612}
\symrange{"2600}{"27BF}
\catcode"2705=13 \begingroup\lccode`\~="2705 \lowercase{\endgroup\def~}{{\symfont\char"2713}}\symblank{"26BD}\symblank{"26A1}
\symrange{"2460}{"257F}\symactive{"223C}\symactive{"2248}\symactive{"2260}
\symactive{"01F8}\symactive{"01F9}\symactive{"1E3E}\symactive{"1E3F}
\symhyph{"2011}\symhyph{"2E1A}\symblankrange{"1F300}{"1FAFF}\symblank{"FE0F}
% Caractères chinois : WenQuanYi Zen Hei (sous-ensemble GB2312 + pinyin), gras/italique simulés
\setCJKmainfont{WenQuanYiZenHei-subset.ttf}[Path=../../../fonts/,AutoFakeBold=3,AutoFakeSlant=0.2]
\xeCJKsetup{CJKspace=false}
% Pinyin : voyelles à caron (ǎ ǐ ǒ ǔ ǖ ǘ ǚ ǜ) absentes de la police principale -> caractères actifs
\newfontface\pyfont{WenQuanYiZenHei-subset.ttf}[Path=../../../fonts/]
\newcommand\pyactive[1]{\catcode#1=13 \begingroup\lccode`\~=#1 \lowercase{\endgroup\def~}{{\pyfont\char#1}}}
\pyactive{"01CD}\pyactive{"01CE}\pyactive{"01CF}\pyactive{"01D0}\pyactive{"01D1}\pyactive{"01D2}\pyactive{"01D3}\pyactive{"01D4}\pyactive{"01D5}\pyactive{"01D6}\pyactive{"01D7}\pyactive{"01D8}\pyactive{"01D9}\pyactive{"01DA}\pyactive{"01DB}\pyactive{"01DC}
% Maths en sans-serif (Fira Math) avec chiffres et lettres droites de Plus
% Jakarta, pour que les formules se fondent dans le texte.
\usepackage[math-style=ISO,bold-style=ISO]{unicode-math}
\setmathfont{FiraMath-Regular.otf}[Path=../../../fonts/]
\setmathfont{PlusJakartaSans-Medium.ttf}[Path=../../../fonts/,range={up/num,up/{Latin,latin},\mathup}]
\usepackage{icomma} % virgule décimale sans espace en mode math
% Caractères mathématiques Unicode absents des polices de texte (Plus Jakarta,
% Archivo Black) : rendus par leur équivalent mathématique, en texte comme en
% formule — sinon ils s'affichent en carré vide (ex. « Arithmétique dans ℤ »).
\usepackage{newunicodechar}
\newunicodechar{ℝ}{\ensuremath{\mathbb{R}}}
\newunicodechar{ℤ}{\ensuremath{\mathbb{Z}}}
\newunicodechar{ℂ}{\ensuremath{\mathbb{C}}}
\newunicodechar{ℕ}{\ensuremath{\mathbb{N}}}
\newunicodechar{ℚ}{\ensuremath{\mathbb{Q}}}
\newunicodechar{𝔻}{\ensuremath{\mathbb{D}}}
\newunicodechar{α}{\ensuremath{\alpha}}
\newunicodechar{β}{\ensuremath{\beta}}
\newunicodechar{γ}{\ensuremath{\gamma}}
\newunicodechar{δ}{\ensuremath{\delta}}
\newunicodechar{ε}{\ensuremath{\varepsilon}}
\newunicodechar{θ}{\ensuremath{\theta}}
\newunicodechar{λ}{\ensuremath{\lambda}}
\newunicodechar{μ}{\ensuremath{\mu}}
\newunicodechar{σ}{\ensuremath{\sigma}}
\newunicodechar{τ}{\ensuremath{\tau}}
\newunicodechar{φ}{\ensuremath{\varphi}}
\newunicodechar{ψ}{\ensuremath{\psi}}
\newunicodechar{ω}{\ensuremath{\omega}}
\newunicodechar{Σ}{\ensuremath{\Sigma}}
% majuscules produites par \MakeUppercase dans les titres (α-aminés -> Α)
\newunicodechar{Α}{\ensuremath{\alpha}}
\newunicodechar{Β}{\ensuremath{\beta}}
\newunicodechar{Γ}{\ensuremath{\Gamma}}
\newunicodechar{Δ}{\ensuremath{\Delta}}
\newunicodechar{Ε}{\ensuremath{\varepsilon}}
\newunicodechar{Θ}{\ensuremath{\Theta}}
\newunicodechar{Λ}{\ensuremath{\Lambda}}
\newunicodechar{Μ}{\ensuremath{\mu}}
\newunicodechar{Τ}{\ensuremath{\tau}}
\newunicodechar{Φ}{\ensuremath{\Phi}}
\newunicodechar{Ψ}{\ensuremath{\Psi}}
\newunicodechar{Ω}{\ensuremath{\Omega}}
\newunicodechar{↦}{\ensuremath{\mapsto}}
\newunicodechar{∈}{\ensuremath{\in}}
\newunicodechar{∉}{\ensuremath{\notin}}
\newunicodechar{∧}{\ensuremath{\wedge}}
\newunicodechar{⇔}{\ensuremath{\Leftrightarrow}}
\newunicodechar{⟺}{\ensuremath{\Longleftrightarrow}}
\newunicodechar{⇒}{\ensuremath{\Rightarrow}}
\newunicodechar{⟹}{\ensuremath{\Longrightarrow}}
\newunicodechar{∘}{\ensuremath{\circ}}
\newunicodechar{∪}{\ensuremath{\cup}}
\newunicodechar{∩}{\ensuremath{\cap}}
\newunicodechar{≡}{\ensuremath{\equiv}}
\newunicodechar{⊂}{\ensuremath{\subset}}
\newunicodechar{⊥}{\ensuremath{\perp}}
\newunicodechar{∃}{\ensuremath{\exists}}
\newunicodechar{∀}{\ensuremath{\forall}}
\newunicodechar{∛}{\ensuremath{\sqrt[3]{\,}}}
\newunicodechar{∜}{\ensuremath{\sqrt[4]{\,}}}
\newunicodechar{⃗}{\ensuremath{{}^{\to}}}
\newunicodechar{ⁿ}{\ifmmode^{n}\else\textsuperscript{\ensuremath{n}}\fi}
\newunicodechar{ˣ}{\ifmmode^{x}\else\textsuperscript{\ensuremath{x}}\fi}
\newunicodechar{ᵇ}{\ifmmode^{b}\else\textsuperscript{\ensuremath{b}}\fi}
\newunicodechar{ʳ}{\ifmmode^{r}\else\textsuperscript{\ensuremath{r}}\fi}
\newunicodechar{ᵘ}{\ifmmode^{u}\else\textsuperscript{\ensuremath{u}}\fi}
\newunicodechar{ʸ}{\ifmmode^{y}\else\textsuperscript{\ensuremath{y}}\fi}
\newunicodechar{ᵃ}{\ifmmode^{a}\else\textsuperscript{\ensuremath{a}}\fi}
\newunicodechar{ᵉ}{\ifmmode^{e}\else\textsuperscript{\ensuremath{e}}\fi}
\newunicodechar{ᵏ}{\ifmmode^{k}\else\textsuperscript{\ensuremath{k}}\fi}
\newunicodechar{ᵖ}{\ifmmode^{p}\else\textsuperscript{\ensuremath{p}}\fi}
\newunicodechar{ᴺ}{\ifmmode^{N}\else\textsuperscript{\ensuremath{N}}\fi}
\newunicodechar{ᶜ}{\ifmmode^{c}\else\textsuperscript{\ensuremath{c}}\fi}
\newunicodechar{ʰ}{\ifmmode^{h}\else\textsuperscript{\ensuremath{h}}\fi}
\newunicodechar{ᵠ}{\ifmmode^{q}\else\textsuperscript{\ensuremath{q}}\fi}
\newunicodechar{ₙ}{\ifmmode_{n}\else\textsubscript{\ensuremath{n}}\fi}
\newunicodechar{ₐ}{\ifmmode_{a}\else\textsubscript{\ensuremath{a}}\fi}
\newunicodechar{ᵢ}{\ifmmode_{i}\else\textsubscript{\ensuremath{i}}\fi}
\newunicodechar{ᵣ}{\ifmmode_{r}\else\textsubscript{\ensuremath{r}}\fi}
\newunicodechar{ₖ}{\ifmmode_{k}\else\textsubscript{\ensuremath{k}}\fi}
\newunicodechar{ᵦ}{\ifmmode_{\beta}\else\textsubscript{\ensuremath{\beta}}\fi}
\newunicodechar{ₓ}{\ifmmode_{x}\else\textsubscript{\ensuremath{x}}\fi}
\newfontfamily\popdisplay{ArchivoBlack-Regular.ttf}[Path=../../../fonts/]
\newfontfamily\poplogo{SpaceGrotesk-Bold.ttf}[Path=../../../fonts/]
\newfontfamily\popsemi{PlusJakartaSans-SemiBold.ttf}[Path=../../../fonts/]
\newfontfamily\popxbold{PlusJakartaSans-ExtraBold.ttf}[Path=../../../fonts/]
\newfontfamily\popxboldls{PlusJakartaSans-ExtraBold.ttf}[Path=../../../fonts/,LetterSpace=3.0]

\setlist[enumerate]{leftmargin=1.7em,itemsep=5pt,topsep=4pt}
\setlist[itemize]{leftmargin=1.7em,itemsep=4pt,topsep=4pt,label={\color{poporange}\textbullet}}
\setlength{\parindent}{0pt}
\setlength{\parskip}{5pt}
\renewcommand{\arraystretch}{1.25}

% pgfplots : style Pop par défaut
\pgfplotsset{
  every axis/.append style={axis line style={popink,line width=1pt},tick style={popink},
    grid style={popline,line width=0.5pt},label style={font=\small},tick label style={font=\footnotesize}},
  popcurve/.style={poporange,line width=1.8pt,no markers,smooth},
}
% Même style disponible dans un \draw TikZ simple (hors pgfplots)
\tikzset{popcurve/.style={poporange,line width=1.8pt}}
\tikzset{popgrid/.style={popline,very thin}}
\colorlet{popcurve}{poporange} % le nom du style est parfois utilisé comme couleur

% ============================================================
% Pill / squiggle
% ============================================================
\newcommand{\poppill}[3]{%
  \tikz[baseline=-0.55ex]\node[draw=popink,line width=1.2pt,rounded corners=2.6pt,%
    fill=#2,inner xsep=6.5pt,inner ysep=3pt]{%
    \popxboldls\fontsize{7}{7}\selectfont\color{#3}\MakeUppercase{#1}};%
}
\newcommand{\popsquiggle}[1]{%
  \tikz[baseline=-0.4ex]{
    \draw[#1,line width=2.6pt,line cap=round]
      plot [smooth] coordinates {(0,0.09) (0.34,0.01) (0.68,0.21) (1.02,0.01) (1.36,0.10)};
  }%
}
% Mot-clé surligné (marqueur orange sous le texte)
\newcommand{\cle}[1]{{\popxbold\color{popdark}#1}}

% ============================================================
% En-tête / pied
% ============================================================
\pagestyle{fancy}
\fancyhf{}
\renewcommand{\headrulewidth}{0pt}
\renewcommand{\footrulewidth}{0pt}
\lhead{\raisebox{-0.15ex}{\poplogo\fontsize{13}{13}\selectfont\color{popink}OnBuch\color{poporange}+}}
\rhead{\poppill{\DOCMATIERE\ \textbullet\ \DOCNIVEAU\ \textbullet\ Leçon \DOCLECON}{white}{popink}}
\lfoot{\popxbold\fontsize{7.6}{7.6}\selectfont\color{popmuted}\MakeUppercase{Cours signé OnBuch+}\ \textbullet\ {\popsemi\DOCTITRE}}
\rfoot{\tikz[baseline=-0.6ex]\node[rounded corners=8.5pt,fill=popink,minimum height=17pt,minimum width=17pt,inner xsep=5pt,inner sep=0]{\popxbold\fontsize{8}{8}\selectfont\color{popcream}\thepage\,/\,\pageref*{LastPage}};}

% ============================================================
% Titres
% ============================================================
\newcommand{\chaptitre}[2]{%
  \begin{tcolorbox}[enhanced,colback=poporange,colframe=popink,boxrule=2.6pt,arc=11pt,
      drop shadow=popink,left=15pt,right=15pt,top=12pt,bottom=11pt,before skip=3pt,after skip=18pt]
    \poppill{#2}{popink}{popcream}\par\vspace{9pt}
    {\raggedright\hyphenpenalty=10000\exhyphenpenalty=10000\popdisplay\fontsize{22}{24}\selectfont\color{popcream}\MakeUppercase{#1}\par}
  \end{tcolorbox}
}
% Section numérotée : pastille orange avec le numéro + titre Archivo
\newcounter{popsec}
\newcommand{\coursec}[1]{%
  \stepcounter{popsec}\setcounter{popsub}{0}%
  \par\vspace{16pt}\noindent
  \begin{minipage}{\linewidth}
  \tikz[baseline=-0.7ex]\node[draw=popink,line width=1.6pt,fill=poporange,rounded corners=4pt,minimum size=22pt,inner sep=0]{\popdisplay\fontsize{12}{12}\selectfont\color{popcream}\thepopsec};%
  \hspace{8pt}{\popdisplay\fontsize{15}{17}\selectfont\color{popink}\MakeUppercase{#1}}\par
  \vspace{4pt}\noindent{\color{poporange}\rule{36pt}{3.6pt}}
  \end{minipage}\par\nopagebreak\vspace{8pt}
}
\newcounter{popsub}
\newcommand{\courssub}[1]{%
  \stepcounter{popsub}%
  \par\vspace{9pt}\noindent{\popxbold\fontsize{11.5}{13}\selectfont\color{popink}{\color{poporange}\thepopsec.\thepopsub}\hspace{6pt}#1}\par\nopagebreak\vspace{4pt}
}

% ============================================================
% Boîtes pédagogiques — même recette : fond blanc, bordure épaisse colorée,
% ombre dure encre, badge plein. Titre optionnel [#1].
% ============================================================
\tcbset{popbase/.style={breakable,enhanced,colback=white,boxrule=2.3pt,arc=9pt,
  shadow={3pt}{-3pt}{0pt}{popink},left=14pt,right=14pt,top=11pt,bottom=11pt,before skip=11pt,after skip=15pt,
  fontupper=\popsemi\fontsize{10.6}{14.5}\selectfont\color{popink},
  fontlower=\popsemi\fontsize{10.6}{14.5}\selectfont\color{popink}}}
% Coche dessinée (le glyphe ✓ n'existe pas dans les polices du document)
\newcommand{\popcheck}[1]{\tikz[baseline=-0.5ex]\draw[#1,line width=1.6pt,line cap=round,line join=round](-0.13,0.0)--(-0.04,-0.09)--(0.13,0.1);}
\providecommand{\coche}{\popcheck{popgreen}}
\providecommand{\resultat}[1]{\par\smallskip\noindent\fcolorbox{popgreen}{popgreenL}{\parbox{\dimexpr\linewidth-2\fboxsep-2\fboxrule\relax}{\textbf{Résultat.} #1}}\par\smallskip}
\newcommand{\popboxhead}[3]{\poppill{#1}{#2}{white}\ifx\relax#3\relax\else\hspace{6pt}{\popxbold\fontsize{10.6}{12}\selectfont\color{popink}#3}\fi\par\vspace{7pt}}

\newtcolorbox{aretenir}[1][]{popbase,colframe=popgreen,before upper={\popboxhead{À retenir}{popgreen}{#1}}}
% Texte en chinois (document de lecture, dialogue, extrait) : cadre
% d'étude. Un titre optionnel (source, auteur) s'affiche dans l'en-tête.
\newtcolorbox{textebox}[1][]{popbase,colframe=popblue,before upper={\popboxhead{文本 · Texte}{popblue}{#1}},after upper={}}
% Marques ✓ / ✗ (forme correcte / fautive) souvent utilisées par les modèles
\providecommand{\cross}{\textcolor{poppink}{\ensuremath{\times}}}
\providecommand{\xmark}{\cross}\providecommand{\crossmark}{\cross}\providecommand{\cmark}{\ensuremath{\checkmark}}
% Ligne de réponse / justification à compléter (inventée par les modèles)
\providecommand{\justification}[1]{\par\noindent\textit{Justification~:} \dotfill\par}
% \textipa (package tipa non chargé) : la police ne couvre pas l'API -> simple passage
\providecommand{\textipa}[1]{#1}
% Soulignement (ulem non chargé) : \uline, \ul
\providecommand{\uline}[1]{\underline{#1}}\providecommand{\ul}[1]{\underline{#1}}
% \glossaire{\item ...} inventée par les modèles : liste de glossaire
\providecommand{\glossaire}[1]{\begin{itemize}#1\end{itemize}}
% \alert (beamer) inventée par les modèles : texte mis en évidence
\providecommand{\alert}[1]{\textbf{\textcolor{poppink}{#1}}}
% variantes ASCII d'ordinaux écrites en macro
\providecommand{\iere}{\textsuperscript{re}}\providecommand{\ieme}{\textsuperscript{e}}\providecommand{\ere}{\textsuperscript{re}}\providecommand{\eme}{\textsuperscript{e}}\providecommand{\er}{\textsuperscript{er}}\providecommand{\ie}{\textsuperscript{e}}
% Ordinaux français écrits en macro par les modèles : 1\ière, 3\ième
\newcommand{\ière}{\textsuperscript{re}}\newcommand{\ième}{\textsuperscript{e}}
% \exemple (inventée par les modèles) : étiquette « Exemple : »
\providecommand{\exemple}{\textbf{Exemple~:}\ }
% \square utilisable aussi hors mode math (cases à cocher)
\AtBeginDocument{\let\squareMath\square\renewcommand{\square}{\ensuremath{\squareMath}}}
% Mot ou phrase en chinois à l'intérieur d'un paragraphe français
\newcommand{\zh}[1]{\ifmmode\text{#1}\else{#1}\fi}
\let\eng\zh \let\esp\zh \let\all\zh % alias tolérés
% flèches d'intonation inventées par les modèles
\providecommand{\textsearrow}{}\renewcommand{\textsearrow}{\ensuremath{\searrow}}\providecommand{\textnearrow}{}\renewcommand{\textnearrow}{\ensuremath{\nearrow}}
% \fr{..} : mot français (inventé par les modèles)
\providecommand{\fr}[1]{\textit{#1}}
% zhbox : encadré de texte chinois inventé par les modèles
\newenvironment{zhbox}{\begin{quote}}{\end{quote}}
% \courssubsub : titre de niveau 3 inventé par les modèles
\providecommand{\courssubsub}[1]{\par\medskip\noindent\textbf{#1}\par\smallskip}
% \zhbf : texte chinois en gras inventé par les modèles
\providecommand{\zhbf}[1]{\textbf{#1}}
% \vs : « versus » inventé par les modèles
\providecommand{\vs}{\textit{vs}\ }
% \blank : case à compléter inventée par les modèles
\providecommand{\blank}[1][]{\underline{\hspace{1.6em}}}
\newtcolorbox{exemplebox}[1][]{popbase,colframe=poporange,before upper={\popboxhead{Exemple}{poporange}{#1}}}
\newtcolorbox{definition}[1][]{popbase,colframe=popblue,before upper={\popboxhead{Définition}{popblue}{#1}}}
\newtcolorbox{propriete}[1][]{popbase,colframe=poppurple,before upper={\popboxhead{Propriété}{poppurple}{#1}}}
\newtcolorbox{methode}[1][]{popbase,colframe=popgold,before upper={\popboxhead{Méthode}{popgold}{#1}}}
\newtcolorbox{attention}[1][]{popbase,colframe=poppink,before upper={\popboxhead{Attention}{poppink}{#1}}}
\newtcolorbox{experience}[1][]{popbase,colframe=popblue,colback=popblueL,before upper={\popboxhead{Expérience}{popblue}{#1}}}
% « Le savais-tu ? » : carte violette pleine (contexte, culture, Cameroun)
\newtcolorbox{savaistu}[1][]{popbase,colback=poppurple,colframe=popink,
  fontupper=\popsemi\fontsize{10.4}{14.2}\selectfont\color{white},
  before upper={\poppill{Le savais-tu\,?}{white}{poppurple}\ifx\relax#1\relax\else\hspace{6pt}{\popxbold\color{white}#1}\fi\par\vspace{7pt}}}
% Exercice résolu : énoncé en haut, \tcblower puis la solution sur fond crème
\newcounter{popexo}\newcounter{popexr}
\newtcolorbox{exoresolu}[1][]{popbase,colframe=poporange,colbacklower=popcream,
  segmentation style={popink,line width=1.2pt,dashed},
  before upper={\stepcounter{popexr}\popboxhead{Exercice résolu \thepopexr}{poporange}{#1}},
  before lower={\poppill{Solution}{popink}{popcream}\par\vspace{6pt}}}
% Exercice (non résolu en place) — la correction est dans la partie Corrigés
\newtcolorbox{exercice}[1][]{popbase,colframe=popink,
  before upper={\stepcounter{popexo}\popboxhead{Exercice \thepopexo}{popink}{#1}}}
% Corrigé
\newtcolorbox{corrige}[1][]{popbase,colframe=popgreen,colback=popgreenL,
  before upper={\popboxhead{Corrigé}{popgreen}{#1}}}
% Objectifs / prérequis / bilan
\newtcolorbox{objectifs}{popbase,colframe=popink,colback=poporangeL,
  before upper={\popboxhead{Objectifs de la leçon}{popink}{}}}
\newtcolorbox{prerequis}{popbase,colframe=popink,colback=popblueL,
  before upper={\popboxhead{Prérequis}{popblue}{}}}
\newtcolorbox{bilan}{popbase,colframe=popink,colback=white,boxrule=2.8pt,
  before upper={\popboxhead{Fiche bilan}{poporange}{L'essentiel à connaître pour l'examen}}}
% Figure encadrée avec légende
\newtcolorbox{popfigure}[1][]{enhanced,breakable=false,colback=white,colframe=popline,boxrule=1.4pt,arc=8pt,
  left=10pt,right=10pt,top=10pt,bottom=8pt,before skip=10pt,after skip=12pt,halign=center,
  lower separated=false,halign lower=center,
  fontlower=\popsemi\fontsize{9}{11.5}\selectfont\color{popmuted},
  #1}
\newcommand{\legende}[1]{\par\vspace{6pt}{\popsemi\fontsize{9}{11.5}\selectfont\color{popmuted}#1}}

% ============================================================
% Signature OnBuch+
% ============================================================
\newcommand{\poplogoinline}[1]{{\poplogo\fontsize{#1}{#1}\selectfont\color{popink}OnBuch\color{poporange}+}}
\newcommand{\signatureonbuch}{%
  \begin{tcolorbox}[enhanced,colback=popink,colframe=popink,boxrule=2.2pt,arc=10pt,
      drop shadow=poporange,left=14pt,right=14pt,top=10pt,bottom=10pt,before skip=0pt,after skip=0pt]
    \tikz[baseline=-0.7ex]\node[circle,fill=popgreen,draw=popcream,line width=1.3pt,minimum size=22pt,inner sep=0]{\popcheck{white}};%
    \hspace{9pt}\begin{minipage}[c]{0.62\linewidth}
      {\popxbold\fontsize{12}{13}\selectfont\color{popcream}Signé\ \textbullet\ {\poplogo OnBuch\color{poporange}+}}\par\vspace{2pt}
      {\raggedright\popsemi\fontsize{8}{10.5}\selectfont\color{popline}\MakeUppercase{Équipe pédagogique OnBuch+}\\\MakeUppercase{Conforme au programme officiel MINESEC}\par}
    \end{minipage}\hfill
    {\poplogo\fontsize{22}{22}\selectfont\color{popcream}OnBuch\color{poporange}+}
  \end{tcolorbox}%
}

% ============================================================
% Page de garde
% #1 titre · #2 matière · #3 niveau · #4 module · #5 n° leçon · #6 durée ·
% #7 description / objectifs (texte ou itemize)
% ============================================================
\newcommand{\pagegarde}[7]{%
  \thispagestyle{empty}
  \newgeometry{a4paper,margin=1.7cm,top=1.6cm,bottom=1.4cm}%
  \begin{tikzpicture}[remember picture,overlay]
    \fill[poporange!18] ([xshift=-3.2cm,yshift=-3.6cm]current page.north east) circle (4.6cm);
    \fill[poppurple!14] ([xshift=2.4cm,yshift=7.6cm]current page.south west) circle (3.8cm);
  \end{tikzpicture}%
  {\poplogo\fontsize{30}{30}\selectfont\color{popink}OnBuch\color{poporange}+}\hfill
  \raisebox{0.6ex}{\poppill{Cours signé \textbullet\ Premium}{popink}{popcream}}\par
  \vspace{1pt}\popsquiggle{poppurple}\par
  \vspace{8pt}
  \begin{tcolorbox}[enhanced,colback=poporange,colframe=popink,boxrule=2.8pt,arc=13pt,
      drop shadow=popink,left=18pt,right=18pt,top=12pt,bottom=13pt,after skip=10pt]
    \poppill{#2 \textbullet\ #3}{popink}{popcream}\hspace{5pt}\poppill{#4}{white}{popink}\par\vspace{8pt}
    {\popdisplay\fontsize{14}{14}\selectfont\color{popink}LEÇON #5}\par\vspace{4pt}
    {\raggedright\hyphenpenalty=10000\exhyphenpenalty=10000\popdisplay\fontsize{25}{27}\selectfont\color{popcream}\MakeUppercase{#1}\par}
    \vspace{8pt}
    {\popxbold\fontsize{9.5}{9.5}\selectfont\color{popink}\MakeUppercase{Durée conseillée : #6}}
  \end{tcolorbox}
  \begin{tcolorbox}[enhanced,colback=white,colframe=popink,boxrule=2.2pt,arc=10pt,
      drop shadow=popink,left=16pt,right=16pt,top=10pt,bottom=10pt,after skip=8pt]
    \poppill{Dans cette leçon}{poppurple}{white}\par\vspace{5pt}
    \popsemi\fontsize{9.3}{12.6}\selectfont\color{popink}#7
  \end{tcolorbox}
  \noindent\begin{minipage}[t]{0.487\linewidth}
    \begin{tcolorbox}[enhanced,colback=popgreen,colframe=popink,boxrule=2.2pt,arc=10pt,
        drop shadow=popink,left=11pt,right=11pt,top=10pt,bottom=10pt,height=3.1cm,valign=center]
      \tcbox[colback=white,colframe=popink,boxrule=1.3pt,arc=5pt,boxsep=0pt,nobeforeafter,
        left=3pt,right=3pt,top=3pt,bottom=3pt,valign=center]{\qrcode[height=1.9cm]{https://play.google.com/store/apps/details?id=com.luvvix.onbuch&hl=fr}}%
      \hspace{8pt}\begin{minipage}[c]{0.5\linewidth}
        {\popdisplay\fontsize{11}{12.5}\selectfont\color{white}\MakeUppercase{Télécharge OnBuch}}\par\vspace{4pt}
        {\raggedright\popsemi\fontsize{8.4}{11}\selectfont\color{white}Gratuit \textbullet\ Google Play\\Tuteur IA, annales, concours, résultats\par}
      \end{minipage}
    \end{tcolorbox}
  \end{minipage}\hfill
  \begin{minipage}[t]{0.487\linewidth}
    \begin{tcolorbox}[enhanced,colback=poppurple,colframe=popink,boxrule=2.2pt,arc=10pt,
        drop shadow=popink,left=11pt,right=11pt,top=10pt,bottom=10pt,height=3.1cm,valign=center]
      \tikz[baseline=-0.7ex]\node[circle,fill=white,draw=popink,line width=1.3pt,minimum size=1.6cm,inner sep=0]{\popdisplay\fontsize{24}{24}\selectfont\color{poppurple}+};%
      \hspace{8pt}\begin{minipage}[c]{0.6\linewidth}
        {\popdisplay\fontsize{11}{12.5}\selectfont\color{white}\MakeUppercase{Passe à OnBuch+}}\par\vspace{4pt}
        {\raggedright\popsemi\fontsize{8.4}{11}\selectfont\color{white}Tous les cours signés, Tuteur IA illimité, fiches PDF — onglet OnBuch+ de l'app\par}
      \end{minipage}
    \end{tcolorbox}
  \end{minipage}
  \vspace{8pt}
  \popcontenu
  \vfill
  \signatureonbuch
  \restoregeometry
  \newpage
}

% Rangée « Ce cours contient » (page de garde)
\newcommand{\poptile}[3]{% #1 couleur #2 gros texte #3 légende
  \begin{tcolorbox}[enhanced,colback=#1,colframe=popink,boxrule=2pt,arc=9pt,shadow={2.5pt}{-2.5pt}{0pt}{popink},
      left=6pt,right=6pt,top=9pt,bottom=9pt,halign=center,valign=center,height=1.75cm,width=\linewidth,nobeforeafter]
    {\popdisplay\fontsize{12.5}{13}\selectfont\color{white}\MakeUppercase{#2}}\par\vspace{3pt}
    {\centering\popsemi\fontsize{7.8}{9.5}\selectfont\color{white}#3\par}
  \end{tcolorbox}}
\newcommand{\popcontenu}{%
  \par\noindent{\popxboldls\fontsize{7.5}{7.5}\selectfont\color{popmuted}CE COURS SIGNÉ CONTIENT}\par\vspace{7pt}
  \noindent\begin{minipage}[t]{0.235\linewidth}\poptile{popblue}{Cours}{complet, illustré, pas à pas}\end{minipage}\hfill
  \begin{minipage}[t]{0.235\linewidth}\poptile{poporange}{Méthodes}{et exercices résolus}\end{minipage}\hfill
  \begin{minipage}[t]{0.235\linewidth}\poptile{poppink}{Exercices}{type examen + corrigés}\end{minipage}\hfill
  \begin{minipage}[t]{0.235\linewidth}\poptile{popgreen}{Fiche bilan}{l'essentiel à retenir}\end{minipage}\par}

% Sommaire
\newtcolorbox{sommaire}{enhanced,colback=white,colframe=popink,boxrule=2.2pt,arc=10pt,
  drop shadow=popink,left=18pt,right=18pt,top=13pt,bottom=13pt,before skip=6pt,after skip=20pt,
  before upper={\poppill{Sommaire}{poppurple}{white}\par\vspace{9pt}\popsemi\fontsize{10.6}{14.5}\selectfont\color{popink}}}

% Signature de fin de cours — poussée tout en bas de la dernière page
\newcommand{\findecours}{%
  \par\vfill\nopagebreak
  \begin{center}\popsquiggle{poporange}\end{center}\vspace{4pt}
  \signatureonbuch
}
\AtBeginDocument{\def\llbracket{\mathopen{[\mkern-3.5mu[}}\def\rrbracket{\mathclose{]\mkern-3.5mu]}}} % intervalles d'entiers (glyphe ⟦ absent de la police)
% Glyphes absents de Fira Math : redéfinis à partir de glyphes disponibles
\AtBeginDocument{%
  \def\setminus{\mathbin{\backslash}}\let\smallsetminus\setminus
  \def\vdots{\mathord{\vbox{\baselineskip3.2pt\lineskiplimit0pt\kern4pt\hbox{.}\hbox{.}\hbox{.}}}}%
  \def\ddots{\mathinner{\mkern1mu\raise6.5pt\hbox{.}\mkern2mu\raise3.6pt\hbox{.}\mkern2mu\raise0.7pt\hbox{.}\mkern1mu}}%
  \def\triangle{\mathord{\scalebox{0.9}{$\symup{\Delta}$}}}%
  \def\nabla{\mathord{\raisebox{\depth}{\scalebox{1}[-1]{$\symup{\Delta}$}}}}%
  \def\checkmark{\popcheck{popdark}}%
  \def\star{\mathbin{\ast}}\def\bigstar{\mathord{\ast}}%
  \def\diamond{\mathbin{\scalebox{0.6}{$\Diamond$}}}%
  \def\bigcup{\mathop{\mathchoice{\raisebox{-0.3em}{\scalebox{1.7}{$\cup$}}}{\scalebox{1.25}{$\cup$}}{\cup}{\cup}}\displaylimits}%
  \def\bigcap{\mathop{\mathchoice{\raisebox{-0.3em}{\scalebox{1.7}{$\cap$}}}{\scalebox{1.25}{$\cap$}}{\cap}{\cap}}\displaylimits}%
  \def\lesseqqgtr{\mathrel{\lessgtr}}\def\bigoplus{\mathop{\oplus}\displaylimits}%
}
\providecommand{\ssi}{si et seulement si} % abréviation fréquente
\AtBeginDocument{\providecommand{\wideparen}{\overparen}} % arc de cercle

%% FIGFIX-BEGIN v3 — correctifs globaux des figures (docs/onbuch-generation/tools/figfix)
\usepackage{adjustbox}\usepackage{etoolbox}
% toute figure TikZ/pgfplots plus large que la ligne est réduite à la largeur disponible
\BeforeBeginEnvironment{tikzpicture}{\begin{adjustbox}{max width=\linewidth}}
\AfterEndEnvironment{tikzpicture}{\end{adjustbox}}
% légendes pgfplots lisibles par-dessus les courbes
\pgfplotsset{every axis/.append style={legend style={fill=white,fill opacity=0.92,text opacity=1,draw=black!15,inner sep=3pt}}}
% étiquettes d'arêtes (syntaxe edge["..."]) avec fond blanc
\tikzset{every edge quotes/.append style={fill=white,inner sep=1.2pt}}
% bulles de cartes mentales : texte à la ligne dans la bulle, bulle un peu plus grande
\tikzset{every concept/.append style={text width=2.0cm,align=center,minimum size=2.3cm,inner sep=2pt}}
%% FIGFIX-END
\begin{document}

\pagegarde{Leçon 38 — Vocabulaire et compréhension de texte : rédaction d’un CV}{Chinois}{Tle A}{Module 5 : Mass médias et télécommunication}{ZH38}{2 h}{Cette leçon de vocabulaire et compréhension de texte initie les élèves de Terminale A à la lecture et à la rédaction d'un curriculum vitae en chinois. À travers un texte support authentique, un glossaire structuré, l'étude des radicaux et des questions de compréhension, les élèves acquièrent le lexique professionnel indispensable au monde du travail sino-camerounais.
\vspace{4pt}
\begin{multicols}{2}\small
\begin{itemize}
\item À la fin de cette leçon, je sais lire et comprendre un court texte chinois présentant un CV.
\item À la fin de cette leçon, je sais reconnaître, prononcer (pinyin et tons) et employer le vocabulaire du CV et du parcours professionnel.
\item À la fin de cette leçon, je sais identifier les rubriques d'un CV chinois (informations personnelles, formation, expérience, compétences).
\item À la fin de cette leçon, je sais écrire correctement quelques caractères clés du thème en respectant les radicaux et l'ordre des traits.
\item À la fin de cette leçon, je sais répondre à des questions de compréhension sur un texte de type CV.
\item À la fin de cette leçon, je sais utiliser les collocations essentielles (找工作, 工作经验, 毕业, 申请) dans des phrases complètes.
\end{itemize}\end{multicols}}

\begin{sommaire}
\begin{multicols}{2}
\begin{enumerate}
\item Mise en route et découverte du thème
\item Texte support : le CV de Li Ming
\item Glossaire du thème
\item Clés (radicaux) et ordre des traits
\item Compréhension du texte et collocations
\item Production guidée : mon mini-CV
\item Activité d'intégration
\item Méthodes et exercices résolus
\item Exercices
\item Corrigés des exercices
\item Fiche bilan
\end{enumerate}
\end{multicols}
\end{sommaire}

\begin{objectifs}
\begin{itemize}
\item À la fin de cette leçon, je sais lire et comprendre un court texte chinois présentant un CV.
\item À la fin de cette leçon, je sais reconnaître, prononcer (pinyin et tons) et employer le vocabulaire du CV et du parcours professionnel.
\item À la fin de cette leçon, je sais identifier les rubriques d'un CV chinois (informations personnelles, formation, expérience, compétences).
\item À la fin de cette leçon, je sais écrire correctement quelques caractères clés du thème en respectant les radicaux et l'ordre des traits.
\item À la fin de cette leçon, je sais répondre à des questions de compréhension sur un texte de type CV.
\item À la fin de cette leçon, je sais utiliser les collocations essentielles (找工作, 工作经验, 毕业, 申请) dans des phrases complètes.
\item À la fin de cette leçon, je sais rédiger un mini-CV guidé en chinois avec mes propres informations.
\item À la fin de cette leçon, je sais comparer les conventions du CV chinois et du CV camerounais/français.
\end{itemize}
\end{objectifs}

\begin{prerequis}
\begin{itemize}
\item Vocabulaire de la famille et de l'identité (姓名, 年龄, 国籍) vu en Seconde et Première
\item Vocabulaire de l'école et des études (学校, 学生, 学习, 专业) vu en Première
\item Expression de la date et des années (年, 月, 日) et des durées
\item Structures de base : 是, 有, 在 + lieu, verbes modaux 会, 能
\item Notions sur les médias et la communication (Module 5, leçons précédentes)
\end{itemize}
\end{prerequis}

\newpage

\coursec{Mise en route et découverte du thème}

\courssub{Situation d'accroche : Marie et l'offre d'emploi de Douala}

Imaginons la scène. À Douala, capitale économique du Cameroun, une grande entreprise chinoise de télécommunications affiche une annonce : elle recrute des jeunes diplômés \textbf{bilingues français-chinois} pour un stage rémunéré, avec possibilité d'embauche. La consigne est claire : \emph{envoyez votre CV en chinois}.

Marie, élève de Terminale A, apprend le chinois depuis quatre ans. Elle sait se présenter à l'oral : \zh{我叫玛丽，我是喀麦隆人，我学中文学了四年。} (\emph{Wǒ jiào Mǎlì, wǒ shì Kāmàilóng rén, wǒ xué Zhōngwén xué le sì nián.} — « Je m'appelle Marie, je suis Camerounaise, j'apprends le chinois depuis quatre ans. »). Mais rédiger un document officiel, c'est autre chose ! Elle découvre que le CV chinois, \zh{简历} (\emph{jiǎnlì}), a ses propres codes : une photo d'identité, des rubriques précises, un ordre bien défini, un style sobre et direct.

\begin{textebox}[La phrase clé de la leçon]
找工作之前，你需要一份简历。\\
\emph{Zhǎo gōngzuò zhīqián, nǐ xūyào yí fèn jiǎnlì.}\\
« Avant de chercher un emploi, tu as besoin d'un CV. »
\end{textebox}

Observons cette phrase de très près, car elle contient déjà trois trésors :

\begin{itemize}
\item \zh{找工作} (\emph{zhǎo gōngzuò}) : « chercher un emploi ». Le verbe \zh{找} (\emph{zhǎo}, chercher) se combine avec \zh{工作} (\emph{gōngzuò}, travail, emploi).
\item \zh{之前} (\emph{zhīqián}) : « avant ». Attention à l'ordre chinois : on dit littéralement « chercher un emploi \emph{avant} », c'est-à-dire que \zh{之前} se place \textbf{après} le groupe verbal, contrairement au français « avant de ». Structure : \cle{verbe / groupe verbal + 之前}.
\item \zh{一份简历} (\emph{yí fèn jiǎnlì}) : « un CV ». Le mot \zh{份} (\emph{fèn}) est le \cle{classificateur} des documents, journaux et portions : on ne peut pas dire \zh{一个简历}, il faut \zh{一份简历}.
\end{itemize}

\textbf{Problématique de la leçon :} comment présenter son parcours en chinois — son identité, ses études, ses compétences — pour décrocher un stage ou un emploi, au Cameroun ou en Chine ? Cette leçon nous donne toutes les clés : à la fin, chacun rédigera son propre mini-CV en chinois.

\courssub{Brainstorming : qu'est-ce qu'un CV et à quoi sert-il au Cameroun ?}

Avant d'ouvrir le dictionnaire chinois, réfléchissons en français. Un \cle{CV} (curriculum vitae, « déroulement de la vie » en latin) est un document qui résume votre parcours : qui vous êtes, ce que vous avez étudié, ce que vous savez faire, ce que vous avez déjà accompli.

Au Cameroun, le CV est indispensable dans de nombreuses situations de la vie scolaire et professionnelle :

\begin{itemize}
\item les \textbf{concours} administratifs et les grandes écoles (ENAM, ENS, écoles de génie, de médecine) ;
\item la recherche d'un \textbf{emploi} dans une entreprise, une banque, une ONG, une administration ;
\item les demandes de \textbf{stage} académique ou professionnel, souvent exigées en fin de cycle ;
\item les candidatures à des \textbf{bourses} d'études, notamment les bourses chinoises (gouvernement chinois, Instituts Confucius), pour lesquelles un CV en chinois est un vrai atout ;
\item les dossiers pour des programmes d'échange, du volontariat ou des concours de jeunesse.
\end{itemize}

Un bon CV est \textbf{bref} (une page, deux au maximum), \textbf{clair} (des rubriques bien séparées) et \textbf{honnête} (on n'invente rien : un employeur peut toujours vérifier). Ces trois qualités, nous allons les retrouver dans le mot chinois lui-même.

\courssub{Le mot 简历 (jiǎnlì) : un mot qui se raconte}

\begin{definition}[简历 jiǎnlì]
\zh{简历} (\emph{jiǎnlì}), nom, signifie « CV, résumé du parcours ». Il est formé de deux caractères : \zh{简} (\emph{jiǎn}) = « simple, bref » et \zh{历} (\emph{lì}) = « parcours, expérience, histoire ». Un CV chinois est donc littéralement un « parcours simplifié ». On dit aussi \zh{个人简历} (\emph{gèrén jiǎnlì}, « CV personnel ») en titre de document.
\end{definition}

Décortiquons les deux caractères, car comprendre leur composition aide à les mémoriser et à ne plus les confondre.

\textbf{1. 简 (jiǎn), « simple, bref ».} Ce caractère est composé de la clé du \cle{bambou} \zh{⺮} (en haut) et de \zh{间} (\emph{jiān}, « entre, intervalle ») en dessous. Le bambou renvoie aux lamelles de bambou sur lesquelles on écrivait autrefois en Chine : écrire « simple », c'était écrire court sur des lamelles ! On retrouve \zh{简} dans \zh{简单} (\emph{jiǎndān}, simple, facile) et dans \zh{简体字} (\emph{jiǎntǐzì}, caractères simplifiés).

\textbf{2. 历 (lì), « parcours, expérience ».} Ce caractère contient \zh{厂} (la falaise, clé n\textdegree 27) et \zh{力} (\emph{lì}, la force). On le retrouve dans \zh{历史} (\emph{lìshǐ}, l'histoire, la matière scolaire) et \zh{经历} (\emph{jīnglì}, l'expérience vécue). Un CV, c'est bien la « force » de votre parcours mise par écrit !

\begin{exemplebox}[Le mot 简历 en contexte]
\zh{这是我的简历。} \emph{Zhè shì wǒ de jiǎnlì.} « Voici mon CV. »\\
\zh{请把你的简历发给我。} \emph{Qǐng bǎ nǐ de jiǎnlì fā gěi wǒ.} « Envoie-moi ton CV, s'il te plaît. »\\
\zh{她的简历写得很好。} \emph{Tā de jiǎnlì xiě de hěn hǎo.} « Son CV est très bien rédigé. »\\
\zh{我想申请这份工作，这是我的简历。} \emph{Wǒ xiǎng shēnqǐng zhè fèn gōngzuò, zhè shì wǒ de jiǎnlì.} « Je voudrais postuler à cet emploi, voici mon CV. »
\end{exemplebox}

Remarquez au passage deux structures utiles : \zh{把 + objet + 发给我} (la construction en \cle{把}, qui place l'objet avant le verbe) et \zh{写得很好} (le complément de degré introduit par \cle{得} : « rédigé comment ? — très bien »).

\begin{attention}[Piège : le classificateur de 简历]
Les francophones disent souvent \zh{一个简历} par calque de « un CV ». C'est une erreur ! Le classificateur des documents est \zh{份} (\emph{fèn}) : \zh{一份简历} (\emph{yí fèn jiǎnlì}). Retenez aussi la \textbf{mutation de ton} : \zh{一} (\emph{yī}) devient \emph{yí} (2\textsuperscript{e} ton) devant un 4\textsuperscript{e} ton : \emph{yí fèn}, et non \emph{yī fèn}. De même : \zh{一份工作} (\emph{yí fèn gōngzuò}, un emploi).
\end{attention}

\courssub{Carte mentale : les cinq rubriques du CV chinois}

Un CV chinois standard s'organise autour de cinq grandes rubriques. Mémorisons-les dès maintenant avec une carte mentale : nous les retrouverons dans le texte support de la section 2 et dans notre production finale.

\begin{popfigure}
\begin{tikzpicture}[mindmap, grow cyclic, every node/.style={concept, align=center, font=\small}, root concept/.append style={concept color=popblue, font=\large\bfseries}, level 1/.append style={level distance=3.4cm, sibling angle=72}, level 1 concept/.append style={concept color=poporangeL}]
\node[root concept, align=center]{简历\\ jiǎnlì\\ \footnotesize le CV}
  child { node[align=center]{个人信息\\ gèrén xìnxī\\ \footnotesize infos personnelles} }
  child { node[align=center]{教育经历\\ jiàoyù jīnglì\\ \footnotesize formation} }
  child { node[align=center]{工作经验\\ gōngzuò jīngyàn\\ \footnotesize expérience} }
  child { node[align=center]{技能\\ jìnéng\\ \footnotesize compétences} }
  child { node[align=center]{爱好\\ àihào\\ \footnotesize loisirs} };
\end{tikzpicture}
\legende{Carte mentale du CV chinois : cinq rubriques à connaître par cœur.}
\end{popfigure}

Lisons chaque branche :

\begin{itemize}
\item \zh{个人信息} (\emph{gèrén xìnxī}) : les informations personnelles — nom, date de naissance, adresse, téléphone, courriel ;
\item \zh{教育经历} (\emph{jiàoyù jīnglì}) : le parcours scolaire et universitaire — écoles, diplômes, dates ;
\item \zh{工作经验} (\emph{gōngzuò jīngyàn}) : l'expérience professionnelle — stages, petits boulots, emplois ;
\item \zh{技能} (\emph{jìnéng}) : les compétences — langues, informatique, permis de conduire ;
\item \zh{爱好} (\emph{àihào}) : les loisirs et centres d'intérêt — football, musique, lecture.
\end{itemize}

\begin{aretenir}[Deux mots pour « expérience » : 经历 et 经验]
\zh{经历} (\emph{jīnglì}) désigne le \textbf{parcours vécu} (ce qu'on a traversé : études, séjours) ; \zh{经验} (\emph{jīngyàn}) désigne le \textbf{savoir-faire acquis} (l'expérience professionnelle). Dans un CV : \zh{教育经历} (parcours de formation) mais \zh{工作经验} (expérience de travail). Ne les inversez pas !
\end{aretenir}

\courssub{CV français et CV chinois : petites différences, grandes conséquences}

Marie connaît déjà le CV français. Comparons les deux traditions : cela évitera de « traduire » mécaniquement un CV français en chinois.

\begin{center}
\begin{tabularx}{\linewidth}{|X|X|X|}
\hline
\rowcolor{popblueL} \textbf{Rubrique} & \textbf{CV français (Cameroun)} & \textbf{CV chinois} \\
\hline
Titre du document & Curriculum Vitae / CV & \zh{简历} (\emph{jiǎnlì}) ou \zh{个人简历} \\
\hline
Photo d'identité & facultative, parfois déconseillée & presque toujours exigée (\zh{照片} \emph{zhàopiàn}) \\
\hline
Date de naissance & souvent indiquée au Cameroun & très courante (\zh{出生日期} \emph{chūshēng rìqī}) \\
\hline
État civil & parfois mentionné & souvent mentionné (\zh{婚姻状况} \emph{hūnyīn zhuàngkuàng}) \\
\hline
Formation & du plus récent au plus ancien & du plus récent au plus ancien (\zh{教育经历}) \\
\hline
Expérience & stages, emplois & idem, avec dates précises (\zh{工作经验}) \\
\hline
Langues & rubrique « Langues » & dans \zh{技能} : \zh{法语、英语、汉语} \\
\hline
Loisirs & rubrique souvent négligée & \zh{爱好}, appréciée des recruteurs chinois \\
\hline
Longueur & 1 page & 1 page, style très sobre \\
\hline
\end{tabularx}
\end{center}

\begin{savaistu}[La photo sur le CV : deux cultures, deux habitudes]
En Chine, le CV comporte presque toujours une \textbf{photo d'identité} (\zh{照片}, \emph{zhàopiàn}) en haut à droite, et souvent la \textbf{date de naissance}, parfois même la situation de famille. À l'inverse, dans les pays anglophones comme le Royaume-Uni ou les États-Unis, la photo est de moins en moins courante, car on veut éviter toute discrimination à l'embauche. Au Cameroun, les deux pratiques coexistent : la photo est fréquente dans les dossiers de concours. Quand vous rédigez un CV en chinois, suivez le code chinois : photo nette, sourire discret, fond clair !
\end{savaistu}

\courssub{Méthode : comment aborder un texte chinois inconnu}

Dans la section 2, nous lirons le CV complet de Li Ming. Un texte chinois nouveau peut impressionner : voici la méthode professionnelle en trois étapes, à réutiliser toute l'année et à l'examen.

\begin{methode}[Lire un texte chinois en trois passages]
\begin{enumerate}
\item \textbf{Lire d'abord les titres de rubriques.} Dans un CV, un article ou une notice, les titres (\zh{个人信息}, \zh{教育经历}…) donnent le plan du texte avant même de lire les phrases. C'est la carte du territoire.
\item \textbf{Repérer les mots connus.} Soulignez mentalement les mots déjà appris (noms propres, chiffres, dates, mots du programme). Ils forment des « îlots de sens » qui sécurisent la lecture.
\item \textbf{Procéder à la lecture globale.} Lisez le texte en entier, sans dictionnaire, pour saisir le sens général. Ne traduisez pas mot à mot : devinez le sens des mots inconnus grâce au contexte, aux clés (radicaux) et aux îlots repérés. Le dictionnaire ne vient qu'en dernier recours, pour vérifier.
\end{enumerate}
\end{methode}

\courssub{Révision éclair : le vocabulaire de l'identité et des études}

Pour comprendre le CV de Li Ming, réactivons le vocabulaire appris depuis la Seconde. Ce tableau est votre boîte à outils de la leçon.

\begin{center}
\begin{tabularx}{\linewidth}{|X|X|X|X|}
\hline
\rowcolor{popblueL} Caractères & Pinyin & Nature & Traduction \\
\hline
姓名 & xìngmíng & nom & nom et prénom \\
\hline
出生 & chūshēng & verbe & naître \\
\hline
年龄 & niánlíng & nom & âge \\
\hline
学生 & xuésheng & nom & élève, étudiant \\
\hline
学校 & xuéxiào & nom & école \\
\hline
高中 & gāozhōng & nom & lycée \\
\hline
大学 & dàxué & nom & université \\
\hline
学习 & xuéxí & verbe & étudier, apprendre \\
\hline
毕业 & bìyè & verbe & obtenir son diplôme \\
\hline
会 & huì & verbe modal & savoir (faire) \\
\hline
说 & shuō & verbe & parler \\
\hline
汉语 & Hànyǔ & nom & langue chinoise \\
\hline
法语 & Fǎyǔ & nom & langue française \\
\hline
电话 & diànhuà & nom & téléphone \\
\hline
邮箱 & yóuxiāng & nom & boîte courriel \\
\hline
\end{tabularx}
\end{center}

Testons immédiatement cette remise en mémoire avec un exercice de présentation, la phrase la plus utile d'un CV et d'un entretien.

\begin{exoresolu}[Se présenter en une phrase complète]
Énoncé : complétez la phrase avec les mots suivants : \zh{学生}、\zh{高中}、\zh{汉语}、\zh{岁}.\\
\zh{我叫玛丽，今年十八＿＿，是雅温得一所＿＿的＿＿，我会说法语和＿＿。}
\tcblower
Solution détaillée :\\
\zh{我叫玛丽，今年十八岁，是雅温得一所高中的学生，我会说法语和汉语。}\\
\emph{Wǒ jiào Mǎlì, jīnnián shíbā suì, shì Yǎwēndé yì suǒ gāozhōng de xuésheng, wǒ huì shuō Fǎyǔ hé Hànyǔ.}\\
« Je m'appelle Marie, j'ai dix-huit ans, je suis élève d'un lycée de Yaoundé, je sais parler français et chinois. »\\
Points de méthode : (1) l'âge se dit \zh{nombre + 岁} sans verbe « avoir » ; (2) le classificateur des écoles est \zh{所} (\emph{suǒ}) : \zh{一所高中} (\emph{yì suǒ gāozhōng} — \zh{一} devient \emph{yì} devant un 1\textsuperscript{er}, 2\textsuperscript{e} ou 3\textsuperscript{e} ton) ; (3) « savoir parler » une langue se dit \zh{会说 + langue} ; (4) les noms de langues prennent une majuscule en pinyin : \emph{Fǎyǔ}, \emph{Hànyǔ}.
\end{exoresolu}

\begin{attention}[Erreurs fréquentes des francophones]
\begin{itemize}
\item Ne dites pas \zh{我有十八岁} : le chinois n'utilise pas \zh{有} pour l'âge. Dites \zh{我今年十八岁} (littéralement « moi, cette année, dix-huit ans »).
\item Ne dites pas \zh{我会汉语} pour « je parle chinois » : \zh{会} seul signifie « savoir faire ». La forme correcte est \zh{我会说汉语}.
\item Ordre des mots : en chinois, le complément de lieu se place \textbf{avant} le verbe : \zh{我在雅温得学习} (\emph{Wǒ zài Yǎwēndé xuéxí}, « j'étudie à Yaoundé »), jamais \zh{我学习在雅温得}.
\item Attention aux tons de \zh{岁} (\emph{suì}, 4\textsuperscript{e} ton) et de \zh{会} (\emph{huì}, 4\textsuperscript{e} ton) : un ton faux peut changer le sens ou rendre la phrase incompréhensible.
\end{itemize}
\end{attention}

\courssub{Objectifs de la leçon et tâche finale}

À l'issue de cette leçon, vous serez capables de :

\begin{itemize}
\item lire et comprendre un CV rédigé en chinois (texte support de la section 2) ;
\item reconnaître, prononcer avec les bons tons et employer le vocabulaire du parcours personnel et professionnel ;
\item identifier les clés (radicaux) et l'ordre des traits des caractères du thème ;
\item répondre à des questions de compréhension sur un document authentique.
\end{itemize}

\textbf{Tâche finale (section 6) :} à l'image de Marie, vous rédigerez votre propre \textbf{mini-CV en chinois}, avec les cinq rubriques de la carte mentale : \zh{个人信息}、\zh{教育经历}、\zh{工作经验}、\zh{技能}、\zh{爱好}. Ce document pourra réellement servir un jour, pour un stage à Douala, une bourse pour la Chine ou un premier emploi bilingue. Gardez donc précieusement la phrase d'ouverture en mémoire : \zh{找工作之前，你需要一份简历。} — avant de chercher un emploi, il vous faut un CV. Et ce CV, c'est vous qui allez l'écrire !

\coursec{Mise en route et découverte du thème}

\begin{experience}[Activation des connaissances : qu'est-ce qu'un CV ?]
\begin{enumerate}
\item \textbf{Brainstorming collectif} : Au tableau, notez tous les mots (en français ou en chinois) qui vous viennent à l'esprit à l'évocation de « CV » ou « 简历 ». Le professeur regroupe : identité, formation, expérience, compétences, langues, informatique, loisirs, photo, références.
\item \textbf{Question déclencheur} : \emph{« Dans quel ordre mettez-vous ces rubriques sur un CV français ? Et sur un CV chinois ? »} Constat : l'ordre est quasi identique (identité $\rightarrow$ formation $\rightarrow$ expérience $\rightarrow$ compétences $\rightarrow$ centres d'intérêt).
\item \textbf{Objectif de la leçon} : À l'issue de ces 2 h, vous serez capable de lire, comprendre et rédiger un CV simple en chinois (niveau HSK 3-4), en respectant les conventions chinoises (ordre des dates, vocabulaire spécifique, structure \zh{毕业于}).
\end{enumerate}
\end{experience}

\begin{aretenir}[Le mot \cle{简历}]
\zh{简历} \emph{jiǎnlì} = \zh{简} (simple, concis) + \zh{历} (parcours, histoire, calendrier). Littéralement « parcours concis ». C'est le terme standard pour « CV » / « curriculum vitae ». On dit aussi \zh{履历} \emph{lǚlì} (plus formel, « parcours履行é »).
\end{aretenir}

\coursec{Texte support : le CV de Li Ming}

\courssub{Lecture globale du texte}

\begin{methode}[Avant de lire : la bonne démarche]
Pour lire un CV chinois, procédez en quatre étapes :
\begin{enumerate}
\item \textbf{Repérer les rubriques} : cherchez d'abord les titres de parties (ils se terminent souvent par un deux-points \zh{：}).
\item \textbf{Lire globalement} : parcourez le texte sans dictionnaire pour saisir le sens général (qui ? quoi ? où ? quand ?).
\item \textbf{Lire en détail} : ligne par ligne, en vous aidant du pinyin et du glossaire ci-dessous.
\item \textbf{Vérifier} : répondez mentalement aux questions « Qui ? Quel âge ? Quelle formation ? Quelle expérience ? Quelles langues ? ».
\end{enumerate}
\cle{Astuce} : Dans un CV chinois, l'ordre est immuable : \zh{个人信息} $\rightarrow$ \zh{教育经历} $\rightarrow$ \zh{工作经验} $\rightarrow$ \zh{技能} $\rightarrow$ \zh{爱好} (du plus général au plus spécifique).
\end{methode}

Lisons maintenant le document. Lis d'abord silencieusement, puis nous ferons une lecture à voix haute collective, en insistant sur les tons.

\begin{textebox}[个人简历 — Le CV de Li Ming]
个人简历\\
Gèrén jiǎnlì\\
Curriculum vitae\\[0.4em]
姓名：李明。性别：男。\\
Xìngmíng : Lǐ Míng. Xìngbié : nán.\\
Nom : Li Ming. Sexe : masculin.\\[0.4em]
出生日期：2005年3月12日。国籍：中国。\\
Chūshēng rìqī : 2005 nián 3 yuè 12 rì. Guójí : Zhōngguó.\\
Date de naissance : le 12 mars 2005. Nationalité : chinoise.\\[0.4em]
电话：138-0000-1234。邮箱：liming@mail.com。\\
Diànhuà : yāo sān bā - líng líng líng líng - yāo èr sān sì. Yóuxiāng : liming@mail.com.\\
Téléphone : 138-0000-1234. Courriel : liming@mail.com.\\[0.4em]
教育经历：2023年毕业于北京大学，专业是汉语国际教育。\\
Jiàoyù jīnglì : 2023 nián bìyè yú Běijīng Dàxué, zhuānyè shì Hànyǔ guójì jiàoyù.\\
Formation : diplômé de l'Université de Pékin en 2023, spécialité : enseignement international du chinois.\\[0.4em]
工作经验：2023年到2024年在一家网络公司工作，负责翻译和客户服务。\\
Gōngzuò jīngyàn : 2023 nián dào 2024 nián zài yì jiā wǎngluò gōngsī gōngzuò, fùzé fānyì hé kèhù fúwù.\\
Expérience professionnelle : de 2023 à 2024, travail dans une entreprise Internet, responsable de la traduction et du service client.\\[0.4em]
技能：会说汉语、英语和法语，会用电脑。\\
Jìnéng : huì shuō Hànyǔ, Yīngyǔ hé Fǎyǔ, huì yòng diànnǎo.\\
Compétences : sait parler chinois, anglais et français, sait utiliser l'ordinateur.\\[0.4em]
爱好：运动、音乐和旅行。\\
Àihào : yùndòng, yīnyuè hé lǚxíng.\\
Loisirs : le sport, la musique et les voyages.
\end{textebox}

\begin{attention}[Prononciation des numéros de téléphone \zh{幺} \emph{yāo}]
En chinois, le chiffre 1 se prononce \zh{幺} \emph{yāo} (1\up{e} ton) et non \zh{一} \emph{yī} quand on épelle un numéro de téléphone, un numéro de pièce d'identité ou un code, pour éviter la confusion auditive avec \zh{七} \emph{qī} (7). Le 0 se lit \zh{零} \emph{líng}.\\
\emph{Exercice} : Dictez à voix haute votre numéro camerounais (ex: \zh{6 77 12 34 56} $\rightarrow$ \zh{六 七 七 幺 二 三 四 五 六} \emph{liù qī qī yāo èr sān sì wǔ liù}).
\end{attention}

\courssub{Traduction et analyse linguistique}

\textbf{Traduction littérale (mot à mot) pour comprendre la logique chinoise.}\par
« Personne-simple-histoire. Nom : Li Ming. Sexe : homme. Naître-date : année 2005, mois 3, jour 12. Pays-appartenance : Chine. Électricité-parole : 138-0000-1234. Poste-boîte : liming@mail.com. Éducation-parcours : année 2023 terminer-études-à Université de Pékin, spécialité est chinois international enseignement. Travail-expérience : de l'année 2023 à l'année 2024 dans une entreprise Internet travailler, porter-responsabilité traduction et client-service. Capacités : savoir parler chinois, anglais et français, savoir utiliser ordinateur. Aimer-bien : sport, musique et voyage. »

\textbf{Traduction naturelle (français courant).}\par
« Curriculum vitae. Nom : Li Ming. Sexe : masculin. Date de naissance : le 12 mars 2005. Nationalité : chinoise. Téléphone : 138-0000-1234. Courriel : liming@mail.com. Formation : diplômé de l'Université de Pékin en 2023, spécialité enseignement international du chinois. Expérience professionnelle : de 2023 à 2024, a travaillé dans une entreprise Internet, chargé de la traduction et du service client. Compétences : parle chinois, anglais et français ; maîtrise l'outil informatique. Centres d'intérêt : sport, musique, voyages. »

\begin{propriete}[Points de grammaire clés du texte]
\begin{enumerate}
\item \textbf{Ordre de la date} : \cle{Année + 月 + Jour} (\zh{2005年3月12日}). Du plus grand au plus petit. En français : inverser (Jour + Mois + Année).
\item \textbf{Structure \zh{毕业于} \emph{bìyè yú}} : « Être diplômé de / Sortir diplômé de ». \zh{于} \emph{yú} (préposition formelle « en, à, de ») est \textbf{obligatoire}. Schéma : \textbf{Sujet + 时间 + 毕业于 + Établissement}. Ex: \zh{我2023年毕业于北京大学。}
\item \textbf{Structure \zh{负责} \emph{fùzé}} : « Être responsable de, avoir la charge de ». Se construit directement + Nom ou Verbe. Schéma : \textbf{Sujet + 负责 + Nom/Verbe}. Ex: \zh{负责翻译} (traduction), \zh{负责客户服务} (service client).
\item \textbf{\zh{会} \emph{huì} + Verbe d'action} : Exprime une \textbf{compétence acquise, un savoir-faire} (« savoir faire »). Pas un futur, pas une simple possibilité. \zh{会说汉语} = « savoir parler chinois ». \zh{会用电脑} = « savoir utiliser l'ordinateur / maîtriser l'informatique ».
\item \textbf{Énumération} : Utilisation de la virgule chinoise \zh{、} pour les listes nominales (\zh{汉语、英语和法语}) ; \zh{和} \emph{hé} avant le dernier élément.
\end{enumerate}
\end{propriete}

\begin{exemplebox}[Collocations professionnelles à mémoriser (HSK 4)]
\begin{tabular}{lll}
\zh{毕业于北京大学} & \emph{bìyè yú Běijīng Dàxué} & être diplômé de l'Univ. de Pékin \\
\zh{专业是汉语国际教育} & \emph{zhuānyè shì Hànyǔ guójì jiàoyù} & spécialité : enseignement chinois international \\
\zh{网络公司} & \emph{wǎngluò gōngsī} & entreprise Internet / société du numérique \\
\zh{负责翻译} & \emph{fùzé fānyì} & être responsable de la traduction \\
\zh{客户服务} & \emph{kèhù fúwù} & service client (client + service) \\
\zh{会说汉语/英语/法语} & \emph{huì shuō Hànyǔ/Yīngyǔ/Fǎyǔ} & parler chinois/anglais/français (compétence) \\
\zh{会用电脑} & \emph{huì yòng diànnǎo} & maîtriser l'informatique / savoir utiliser un ordi
\end{tabular}
\end{exemplebox}

\begin{attention}[Pièges fréquents pour francophones]
\begin{itemize}
\item \textbf{\zh{毕业} sans \zh{于}} : \textcolor{red}{$\times$ 我毕业北京大学。} $\rightarrow$ \textcolor{popgreen}{$\checkmark$ \zh{我毕业于北京大学。}} (La préposition \zh{于} ne se traduit pas en français mais est obligatoire en chinois).
\item \textbf{Ordre de la date en traduction} : \textcolor{red}{$\times$ « 2005, 3, 12 »} $\rightarrow$ \textcolor{popgreen}{$\checkmark$ « le 12 mars 2005 »}.
\item \textbf{\zh{工作经验} $\neq$ « expérience de travail »} : En français professionnel, on écrit \textbf{« Expérience professionnelle »}.
\item \textbf{\zh{会} vs \zh{能}} : \zh{会} = compétence apprise (savoir faire) ; \zh{能} = capacité physique/autorisation/circonstance (pouvoir). Pour les langues : toujours \zh{会说}。
\end{itemize}
\end{attention}

\courssub{Lecture à voix haute et travail des tons}

\begin{experience}[Lecture en trois temps (15 min)]
\begin{enumerate}
\item \textbf{Lecture collective modélisée} : Le professeur lit, la classe répète en chœur. Points d'attention : \zh{李明} \emph{Lǐ Míng} (3+2), \zh{毕业于} \emph{bìyè yú} (4+4+2, \zh{于} ton 2), \zh{爱好} \emph{àihào} (4+4, tons descendants secs), \zh{幺} \emph{yāo} (1).
\item \textbf{Lecture individuelle par rubrique} : Chaque élève tire au sort une rubrique (Identité, Formation, Expérience, Compétences, Loisirs) et la lit. La classe note les erreurs de tons sur fiche.
\item \textbf{Dictée de numéros} : Par paires, l'élève A dicte son numéro de téléphone camerounais en chinois (\zh{幺}, \zh{零}), l'élève B l'écrit en chiffres. Inversion des rôles.
\end{enumerate}
\end{experience}

\begin{exoresolu}[Repérage dans le texte : compréhension fine]
Énoncé. Répondez en français en citant la ligne chinoise exacte (caractères + pinyin).
\begin{enumerate}
\item Quelle est la nationalité de Li Ming ?
\item En quelle année a-t-il obtenu son diplôme ? Quelle structure grammaticale marque le lien avec l'université ?
\item Citez les trois langues qu'il parle. Quel verbe exprime la compétence ?
\item Quelles étaient ses fonctions dans l'entreprise Internet ? (Deux noms).
\item Donnez l'ordre chronologique de sa formation et de son expérience.
\end{enumerate}
\tcblower
Solution détaillée.
\begin{enumerate}
\item Nationalité chinoise : \zh{国籍：中国。} \emph{Guójí : Zhōngguó.}
\item Diplômé en 2023. Structure : \zh{2023年毕业于北京大学。} \emph{2023 nián bìyè yú Běijīng Dàxué.} (\zh{毕业于} + lieu).
\item Chinois, anglais, français : \zh{会说汉语、英语和法语。} \emph{Huì shuō Hànyǔ, Yīngyǔ hé Fǎyǔ.} Verbe \zh{会} (compétence acquise).
\item Traduction et service client : \zh{负责翻译和客户服务。} \emph{Fùzé fānyì hé kèhù fúwù.}
\item Formation (2023) $\rightarrow$ Expérience (2023-2024). L'ordre du CV respecte la chronologie inverse (la plus récente en premier) dans chaque rubrique, mais la rubrique Formation précède Expérience.
\end{enumerate}
\end{exoresolu}

\coursec{Glossaire du thème : vocabulaire complet (HSK 3-4)}

\begin{center}
\begin{tabularx}{\linewidth}{|X|X|X|X|}
\hline
\rowcolor{popblueL} \textbf{Caractères} & \textbf{Pinyin} & \textbf{Nature} & \textbf{Traduction} \\
\hline
个人简历 & gèrén jiǎnlì & nom & CV, curriculum vitae \\
\hline
个人信息 & gèrén xìnxī & nom & informations personnelles \\
\hline
姓名 & xìngmíng & nom & nom et prénom \\
\hline
性别 & xìngbié & nom & sexe, genre \\
\hline
出生日期 & chūshēng rìqī & nom & date de naissance \\
\hline
国籍 & guójí & nom & nationalité \\
\hline
电话 & diànhuà & nom & téléphone \\
\hline
邮箱 & yóuxiāng & nom & courriel, boîte mail \\
\hline
教育经历 & jiàoyù jīnglì & nom & formation, parcours scolaire \\
\hline
毕业于 & bìyè yú & verbe + part. & être diplômé de (université) \\
\hline
专业 & zhuānyè & nom & spécialité, filière (univ.) \\
\hline
汉语国际教育 & Hànyǔ guójì jiàoyù & nom propre & enseignement international du chinois (TCSOL) \\
\hline
工作经验 & gōngzuò jīngyàn & nom & expérience professionnelle \\
\hline
网络公司 & wǎngluò gōngsī & nom & entreprise Internet, boîte tech \\
\hline
负责 & fùzé & verbe & être responsable de, gérer \\
\hline
翻译 & fānyì & verbe/nom & traduire / traduction \\
\hline
客户服务 & kèhù fúwù & nom & service client, relation client \\
\hline
技能 & jìnéng & nom & compétences, savoir-faire \\
\hline
会说 & huì shuō & loc. verbale & savoir parler (langue) \\
\hline
会用 & huì yòng & loc. verbale & savoir utiliser, maîtriser \\
\hline
电脑 & diànnǎo & nom & ordinateur \\
\hline
爱好 & àihào & nom & loisirs, centres d'intérêt \\
\hline
运动 & yùndòng & nom/verbe & sport / faire du sport \\
\hline
音乐 & yīnyuè & nom & musique \\
\hline
旅行 & lǚxíng & nom/verbe & voyage / voyager \\
\hline
\multicolumn{4}{|l|}{\textbf{Vocabulaire complémentaire utile pour votre CV}} \\
\hline
地址 & dìzhǐ & nom & adresse \\
\hline
求职意向 & qiúzhí yìxiàng & nom & poste visé, objectif professionnel \\
\hline
实习经历 & shíxí jīnglì & nom & stage, expérience de stage \\
\hline
项目经验 & xiàngmù jīngyàn & nom & expérience de projet \\
\hline
证书 & zhèngshū & nom & certificat, attestation \\
\hline
奖项 & jiǎngxiàng & nom & prix, distinction \\
\hline
自我评价 & zìwǒ píngjià & nom & auto-évaluation, profil \\
\hline
\end{tabularx}
\end{center}

\coursec{Clés (radicaux) et ordre des traits}

\begin{propriete}[Rappel méthodologique : Clé (radical) et ordre des traits]
\begin{itemize}
\item La \cle{clé} (部首 \emph{bùshǒu}) classe le caractère au dictionnaire et donne souvent un indice sémantique.
\item L'\cle{ordre des traits} (笔顺 \emph{bǐshùn}) suit des règles fixes : haut $\rightarrow$ bas, gauche $\rightarrow$ droite, horizontal avant vertical, trait central avant les ailes, cadre avant contenu, fermer le cadre en dernier.
\item Respecter l'ordre assure une écriture lisible, rapide et une reconnaissance OCR/écran tactile fiable.
\end{itemize}
\end{propriete}

\begin{center}
\begin{tabularx}{\linewidth}{|X|X|X|X|X|}
\hline
\rowcolor{poppurpleL} \textbf{Caractère} & \textbf{Clé (Radical)} & \textbf{Sens de la clé} & \textbf{Nb traits} & \textbf{Ordre des traits (description)} \\
\hline
\zh{简} & \zh{竹} (zhú) & bambou (simplifié \zh{⺮}) & 13 & 1-2 bambou (haut), puis \zh{间} (porte + soleil) : haut $\rightarrow$ bas \\
\hline
\zh{历} & \zh{止} (zhǐ) & pied/arrêt (ou \zh{厂} cliff) & 8 & \zh{厂} (cadre haut-gauche), puis \zh{止} à l'intérieur : cadre $\rightarrow$ contenu \\
\hline
\zh{毕} & \zh{十} (shí) & dix / achever & 11 & \zh{十} (haut), puis \zh{比} (bas) : haut $\rightarrow$ bas \\
\hline
\zh{业} & \zh{木} (mù) & arbre / travail & 5 & Horizontal, vertical, gauche, droite, horizontal bas (comme \zh{木} sans le trait du milieu) \\
\hline
\zh{专} & \zh{寸} (cùn) & pouce / mesure & 4 & Point (haut), horizontal, vertical, point final (bas) \\
\hline
\zh{经} & \zh{糸} (mì) & fil de soie (simplifié \zh{纟}) & 8 & \zh{纟} (gauche : 3 traits), puis \zh{巳} (droite : haut $\rightarrow$ bas) \\
\hline
\zh{验} & \zh{马} (mǎ) & cheval & 10 & \zh{马} (gauche : 3 traits), puis \zh{ㄝ} (droite : haut $\rightarrow$ bas) \\
\hline
\zh{技} & \zh{扌} (shǒu) & main (action manuelle) & 7 & \zh{扌} (gauche : 3 traits), puis \zh{枝} (droite : haut $\rightarrow$ bas) \\
\hline
\zh{能} & \zh{月} (yuè) & chair/viande (simplifié de \zh{肉}) & 10 & \zh{月} (gauche), puis \zh{匕} + \zh{匕} (droite : haut $\rightarrow$ bas) \\
\hline
\zh{爱} & \zh{爫} (zhǎo) / \zh{心} (xīn) & griffe / cœur (simplifié : cœur au milieu) & 10 & \zh{爫} (haut), \zh{友} (milieu), \zh{心} (bas) : haut $\rightarrow$ bas \\
\hline
\zh{好} & \zh{女} (nǚ) & femme & 6 & \zh{女} (gauche), \zh{子} (droite) : gauche $\rightarrow$ droite \\
\hline
\end{tabularx}
\end{center}

\begin{experience}[Atelier écriture : « Course aux traits » (10 min)]
\begin{enumerate}
\item Le professeur écrit un caractère au tableau \textbf{trait par trait} dans le désordre.
\item Les élèves lèvent la main pour dire le numéro du trait manquant ou mal placé.
\item Ensuite, chaque élève écrit 3 fois les 5 caractères \zh{简 历 专 业 爱} sur papier quadrillé (格子纸), en respectant l'ordre. Vérification par le pair (correction croisée).
\end{enumerate}
\end{experience}

\coursec{Compréhension du texte et collocations}

\courssub{Exercices de compréhension (niveau examen)}

\begin{exercice}[Questions Vrai / Faux / Non mentionné \emph{(Type Bac)}]
Lisez le CV de Li Ming et cochez V (Vrai), F (Faux) ou NM (Non mentionné).
\begin{enumerate}
\item Li Ming est né en 2005. \hfill $\square$ V \quad $\square$ F \quad $\square$ NM
\item Il a étudié à l'Université de Tsinghua. \hfill $\square$ V \quad $\square$ F \quad $\square$ NM
\item Sa spécialité est l'enseignement du chinois aux étrangers. \hfill $\square$ V \quad $\square$ F \quad $\square$ NM
\item Il a travaillé deux ans dans une entreprise Internet. \hfill $\square$ V \quad $\square$ F \quad $\square$ NM
\item Il était développeur informatique. \hfill $\square$ V \quad $\square$ F \quad $\square$ NM
\item Il parle trois langues étrangères. \hfill $\square$ V \quad $\square$ F \quad $\square$ NM
\item Il sait programmer en Python. \hfill $\square$ V \quad $\square$ F \quad $\square$ NM
\item Ses loisirs sont la lecture et le cinéma. \hfill $\square$ V \quad $\square$ F \quad $\square$ NM
\end{enumerate}
\end{exercice}

\begin{corrige}[Exercice V/F/NM]
\begin{enumerate}
\item \textbf{V} — \zh{出生日期：2005年3月12日。}
\item \textbf{F} — \zh{毕业于北京大学} (Pékin), pas Tsinghua.
\item \textbf{V} — \zh{专业是汉语国际教育} = enseignement international du chinois (TCSOL).
\item \textbf{F} — \zh{2023年到2024年} = 1 an (ou 2 années civiles, mais « deux ans » est faux).
\item \textbf{F} — \zh{负责翻译和客户服务} (traduction + service client), pas développeur.
\item \textbf{F} — Il parle \zh{汉语} (chinois, langue maternelle), \zh{英语} (anglais), \zh{法语} (français). Une seule langue \emph{étrangère} au sens strict (si chinois = L1) ou deux (anglais, français). La phrase dit « trois langues ».
\item \textbf{NM} — \zh{会用电脑} (maîtrise bureautique), pas de mention de Python.
\item \textbf{F} — \zh{运动、音乐和旅行} (sport, musique, voyages).
\end{enumerate}
\end{corrige}

\begin{exercice}[Questions ouvertes : expression écrite guidée]
Répondez en phrases chinoises complètes (caractères + pinyin).
\begin{enumerate}
\item 李明的国籍是什么？(Quelle est la nationalité de Li Ming ?)
\item 他是什么时候毕业的？(Quand a-t-il été diplômé ?)
\item 他在网络公司做什么工作？(Quel travail a-t-il fait dans l'entreprise Internet ?)
\item 他会说哪三种语言？(Quelles sont les trois langues qu'il parle ?)
\item 你觉得李明的简历怎么样？为什么？(Ton avis sur son CV ? Justifie.) \emph{(Ouvert)}
\end{enumerate}
\end{exercice}

\begin{corrige}[Exercice questions ouvertes]
\begin{enumerate}
\item 他的国籍是中国。 \emph{Tā de guójí shì Zhōngguó.}
\item 他2023年毕业于北京大学。 \emph{Tā 2023 nián bìyè yú Běijīng Dàxué.}
\item 他负责翻译和客户服务。 \emph{Tā fùzé fānyì hé kèhù fúwù.}
\item 他会说法语、英语和汉语。 \emph{Tā huì shuō Fǎyǔ, Yīngyǔ hé Hànyǔ.} (Ordre libre)
\item \emph{Exemple de réponse :} 我觉得他的简历很好，因为他会说三种语言，还有工作经验。 \emph{Wǒ juéde tā de jiǎnlì hěn hǎo, yīnwèi tā huì shuō sān zhǒng yǔyán, hái yǒu gōngzuò jīngyàn.}
\end{enumerate}
\end{corrige}

\begin{exercice}[Collocations : Reliez les colonnes]
\begin{center}
\begin{tabular}{lll}
\zh{1. 毕业} & \zh{A. 翻译} & \zh{a. service client} \\
\zh{2. 专业} & \zh{B. 于 北京大学} & \zh{b. enseignement chinois international} \\
\zh{3. 负责} & \zh{C. 是 汉语国际教育} & \zh{c. être diplômé de} \\
\zh{4. 客户服务} & \zh{D. 说 汉语、英语} & \zh{d. responsable de} \\
\zh{5. 会说} & \zh{E. 用 电脑} & \zh{e. savoir parler} \\
\zh{6. 会用} & & \zh{f. savoir utiliser} \\
\end{tabular}
\end{center}
\emph{Consigne : Faites les paires (chiffre + lettre) et donnez le sens français.}
\end{exercice}

\begin{corrige}[Exercice Collocations]
1-B-c (\zh{毕业于北京大学} : être diplômé de l'Univ. de Pékin) \\
2-C-b (\zh{专业是汉语国际教育} : spécialité = enseignement chinois international) \\
3-D-d (\zh{负责翻译} : responsable de la traduction) \\
4-A-a (\zh{客户服务} : service client) \\
5-D-e (\zh{会说法语} : savoir parler français) \\
6-E-f (\zh{会用电脑} : savoir utiliser l'ordi)
\end{corrige}

\coursec{Production guidée : mon mini-CV}

\begin{methode}[Étapes pour rédiger votre CV en chinois (HSK 3-4)]
\begin{enumerate}
\item \textbf{Remplissez le tableau d'identité} : Nom chinois (choisi ou transcrit), sexe, date de naissance (année-mois-jour), nationalité (camerounaise \zh{喀麦隆} \emph{Kāmàilóng}), téléphone (avec \zh{幺}), email.
\item \textbf{Rédigez la rubrique Formation} : \zh{教育经历：202X年毕业于[Lycée/Univ], 专业是[Filière].} Si encore lycéen : \zh{就读于[Lycée], 年级：高三。} (\zh{就读于} \emph{jiùdú yú} = « étudier à »).
\item \textbf{Rédigez Expérience} : \zh{工作经验：202X年到202X年在[Entreprise/Stage], 负责[Tâches].} Ou \zh{实习经历} pour stage. Si aucune : \zh{暂无工作经验。} (\zh{暂无} « pour l'instant pas de »).
\item \textbf{Listez Compétences} : \zh{技能：会说[Langues]，会用[Logiciels/Outils].} Au minimum : \zh{法语、英语、汉语} (si appris), \zh{电脑} (Word, Excel, WeChat).
\item \textbf{Ajoutez Loisirs} : \zh{爱好：[Sport], [Musique], [Voyage/Lecture].} Utilisez \zh{、} et \zh{和}.
\item \textbf{Relisez-vous} : Vérifiez \zh{毕业于}, \zh{负责}, \zh{会}, ordre des dates, ponctuation \zh{： 、 。}.
\end{enumerate}
\end{methode}

\begin{textebox}[Modèle de CV vierge à compléter (Photocopie élève)]
\begin{tabularx}{\linewidth}{|X|}
\hline
\rowcolor{popblueL} \textbf{个人简历} \\
\hline
姓名：\underline{\hspace{4cm}} \quad 性别：\underline{\hspace{2cm}} \\
\hline
出生日期：\underline{\hspace{3cm}} 年 \underline{\hspace{1cm}} 月 \underline{\hspace{1cm}} 日 \quad 国籍：\underline{\hspace{2cm}} \\
\hline
电话：\underline{\hspace{5cm}} \quad 邮箱：\underline{\hspace{5cm}} \\
\hline
\rowcolor{poppurpleL} \textbf{教育经历} \\
\hline
\underline{\hspace{16cm}} \\
\hline
\rowcolor{poppurpleL} \textbf{工作经验 / 实习经历} \\
\hline
\underline{\hspace{16cm}} \\
\hline
\rowcolor{poppurpleL} \textbf{技能} \\
\hline
\underline{\hspace{16cm}} \\
\hline
\rowcolor{poppurpleL} \textbf{爱好} \\
\hline
\underline{\hspace{16cm}} \\
\hline
\end{tabularx}
\end{textebox}

\begin{exemplebox}[Exemple de CV élève (Contexte Cameroun) — Modèle corrigé]
\begin{tabularx}{\linewidth}{|X|}
\hline
\rowcolor{popblueL} \textbf{个人简历} \\
\hline
姓名：\zh{张伟} \emph{Zhāng Wěi} (Jean Dupont) \quad 性别：\zh{男} \emph{nán} \\
\hline
出生日期：\zh{2006年5月20日} \quad 国籍：\zh{喀麦隆} \emph{Kāmàilóng} \\
\hline
电话：\zh{6 77 12 34 56} (\emph{liù qī qī yāo èr sān sì wǔ liù}) \quad 邮箱：\zh{jean.dupont@email.cm} \\
\hline
\rowcolor{poppurpleL} \textbf{教育经历} \\
\hline
\zh{2024年毕业于雅温得第一中学，专业是文学系列。} \\
\emph{2024 nián bìyè yú Yǎwēndé Dìyī Zhōngxué, zhuānyè shì wénxué xìliè.} \\
\hline
\rowcolor{poppurpleL} \textbf{实习经历} \\
\hline
\zh{2023年7月到8月在雅温得一家旅行社实习，负责接待客户。} \\
\emph{2023 nián 7 yuè dào 8 yuè zài Yǎwēndé yì jiā lǚxíngshè shíxí, fùzé jiēdài kèhù.} \\
\hline
\rowcolor{poppurpleL} \textbf{技能} \\
\hline
\zh{会说法语、英语和汉语（HSK 3），会用电脑（Word, Excel, 微信）。} \\
\emph{Huì shuō Fǎyǔ, Yīngyǔ hé Hànyǔ (HSK 3), huì yòng diànnǎo (Word, Excel, Wēixìn).} \\
\hline
\rowcolor{poppurpleL} \textbf{爱好} \\
\hline
\zh{足球、音乐和阅读。} \emph{Zúqiú, yīnyuè hé yuèdú.} \\
\hline
\end{tabularx}
\end{exemplebox}

\begin{exercice}[Production finale : Rédigez votre CV (30 min)]
Remplissez le modèle vierge avec vos vraies informations (ou inventées si vous préférez).
\begin{itemize}
\item \textbf{Contrainte} : Minimum 5 lignes de chinois (Identité + Formation + Expérience/Stage + Compétences + Loisirs).
\item \textbf{Outils autorisés} : Glossaire de la leçon, dictionnaire, professeur.
\item \textbf{Rendu} : Remettez la feuille au professeur pour correction individuelle (code de correction : \zh{GT} grammaire, \zh{VC} vocabulaire, \zh{OS} ordre des traits, \zh{PT} ponctuation).
\end{itemize}
\end{exercice}

\begin{savaistu}[Le CV en Chine vs au Cameroun : points culturels]
\begin{itemize}
\item \textbf{Photo} : En Chine, la photo d'identité (fond bleu/rouge, costume/cravate) est \textbf{quasi obligatoire} en haut à droite. Au Cameroun, elle est recommandée mais de plus en plus optionnelle (lutte contre discrimination).
\item \textbf{Âge / Date de naissance} : En Chine, on met toujours la date de naissance (donc l'âge se calcule). En France/Cameroun, l'âge seul suffit parfois, la date est facultative (loi anti-âge).
\item \textbf{État civil} : \zh{婚姻状况} \emph{hūnyīn zhuàngkuàng} (célibataire/marié) figure souvent sur les CV chinois standards. Au Cameroun, c'est illégal de le demander (non-discrimination).
\item \textbf{Objectif de carrière} : \zh{求职意向} \emph{qiúzhí yìxiàng} (poste visé, salaire attendu, ville) est une rubrique standard en Chine. Au Cameroun, on l'intègre dans la lettre de motivation.
\item \textbf{Références} : En Chine, on note souvent \zh{推荐人} \emph{tuījiànrén} (personne de référence + tel) en bas. Au Cameroun, « Références sur demande ».
\end{itemize}
\end{savaistu}

\begin{aretenir}[Check-list finale avant de rendre votre CV]
\begin{enumerate}
\item \zh{毕业于} + Université (pas \zh{毕业} seul).
\item Dates : \zh{年 月 日} (chiffres) + \zh{到} pour les périodes.
\item \zh{负责} + Nom/Verbe (pas \zh{负责做}... si possible, direct \zh{负责翻译}).
\item Langues : \zh{会说} + Langue.
\item Informatique : \zh{会用} + Outil.
\item Ponctuation : \zh{：} après rubrique, \zh{、} dans listes, \zh{。} fin de phrase.
\item Écriture : Caractères lisibles, ordre des traits respecté (surtout \zh{简 历 专 业 爱 好}).
\end{enumerate}
\end{aretenir}

\coursec{Glossaire du thème}

Dans la section précédente, nous avons lu attentivement le CV de Li Ming : son identité, sa formation, son expérience professionnelle et ses loisirs. Ce texte contenait de nombreux mots nouveaux. Avant de rédiger votre propre mini-CV, il faut maintenant \cle{maîtriser ce vocabulaire} : savoir le reconnaître, le prononcer avec les bons tons, l'écrire et l'employer dans des phrases. C'est l'objet de ce glossaire, organisé par familles de sens, comme les rubriques d'un vrai CV.

\courssub{Le mot central : 简历 (le CV)}

\begin{definition}[简历 jiǎnlì — le CV]
\zh{简历} (jiǎnlì) est un nom qui désigne le \cle{curriculum vitae} : un document bref qui présente l'\textbf{identité} (姓名), la \textbf{formation} (教育) et l'\textbf{expérience} (经验) d'un candidat à un emploi. Le caractère \zh{简} signifie « simple, bref » (clé du bambou ⺮) et \zh{历} signifie « parcours, expérience » : le CV est donc littéralement un « parcours résumé ». On dit aussi \zh{个人简历} (gèrén jiǎnlì), « CV personnel ». Le classificateur des documents comme le CV est \zh{份} (fèn) : \zh{一份简历} (yí fèn jiǎnlì), « un CV ».
\end{definition}

\begin{exemplebox}[Observer avant de mémoriser]
\zh{这是我的简历。} \emph{Zhè shì wǒ de jiǎnlì.} « Voici mon CV. »\\
\zh{请看一下我的简历。} \emph{Qǐng kàn yíxià wǒ de jiǎnlì.} « Veuillez jeter un coup d'œil à mon CV. »
\end{exemplebox}

\courssub{Tableau de vocabulaire : l'identité et les coordonnées}

Voici la première famille de mots : ceux de la rubrique « état civil et coordonnées » du CV.

\begin{center}
\begin{tabularx}{\linewidth}{|X|X|X|X|}
\hline
\rowcolor{popblueL} Caractères & Pinyin & Nature & Traduction \\
\hline
简历 & jiǎnlì & n. & CV, curriculum vitae \\
\hline
姓名 & xìngmíng & n. & nom complet (nom + prénom) \\
\hline
性别 & xìngbié & n. & sexe (masculin / féminin) \\
\hline
出生 & chūshēng & v. & naître \\
\hline
国籍 & guójí & n. & nationalité \\
\hline
电话 & diànhuà & n. & téléphone \\
\hline
邮箱 & yóuxiāng & n. & adresse e-mail, boîte mail \\
\hline
地址 & dìzhǐ & n. & adresse (postale) \\
\hline
\end{tabularx}
\end{center}

Quelques remarques de précision :

\begin{itemize}
\item \zh{姓名} (xìngmíng) désigne le \cle{nom complet}. Pour demander seulement le nom de famille, on utilise \zh{姓} (xìng) : \zh{你姓什么？} \emph{Nǐ xìng shénme ?} « Comment t'appelles-tu (nom de famille) ? »
\item Pour la rubrique \zh{性别} (xìngbié), on écrit \zh{男} (nán), « masculin », ou \zh{女} (nǚ), « féminin ».
\item \zh{出生} (chūshēng) s'emploie avec \zh{于} + le lieu ou la date : \zh{我出生于雅温得。} \emph{Wǒ chūshēng yú Yǎwēndé.} « Je suis né(e) à Yaoundé. »
\item \zh{国籍} (guójí) se construit avec le nom du \textbf{pays}, jamais avec le gentilé : \zh{我的国籍是喀麦隆。} \emph{Wǒ de guójí shì Kāmàilóng.} « Ma nationalité est camerounaise » (litt. « mon pays de nationalité est le Cameroun »). Pour utiliser le gentilé, dites : \zh{我是喀麦隆人。} « Je suis Camerounais(e). »
\item \zh{邮箱} (yóuxiāng) signifie littéralement « boîte postale » ; dans le contexte moderne, c'est l'\textbf{adresse e-mail}. On précise parfois \zh{电子邮箱} (diànzǐ yóuxiāng), « boîte mail électronique ».
\item \zh{电话} (diànhuà) : littéralement « parole électrique » (电 = électricité, 话 = parole). Le numéro de téléphone se dit \zh{电话号码} (diànhuà hàomǎ).
\end{itemize}

\courssub{Tableau de vocabulaire : formation, expérience et compétences}

Deuxième famille : les rubriques « formation », « expérience professionnelle » et « compétences ».

\begin{center}
\begin{tabularx}{\linewidth}{|X|X|X|X|}
\hline
\rowcolor{popblueL} Caractères & Pinyin & Nature & Traduction \\
\hline
毕业 & bìyè & v. & obtenir son diplôme, être diplômé \\
\hline
大学 & dàxué & n. & université \\
\hline
专业 & zhuānyè & n. & spécialité, filière \\
\hline
教育 & jiàoyù & n. & éducation, formation \\
\hline
经验 & jīngyàn & n. & expérience \\
\hline
工作 & gōngzuò & n. / v. & travail ; travailler \\
\hline
公司 & gōngsī & n. & entreprise, société \\
\hline
负责 & fùzé & v. & être responsable de \\
\hline
技能 & jìnéng & n. & compétence \\
\hline
翻译 & fānyì & n. / v. & traduction ; traduire \\
\hline
申请 & shēnqǐng & v. & postuler, faire une demande \\
\hline
爱好 & àihào & n. & loisir, passe-temps \\
\hline
\end{tabularx}
\end{center}

\begin{exemplebox}[Exemples traduits à étudier]
\zh{我想申请这个工作。} \emph{Wǒ xiǎng shēnqǐng zhège gōngzuò.} « Je voudrais postuler à cet emploi. »\\
\zh{她有很多工作经验。} \emph{Tā yǒu hěn duō gōngzuò jīngyàn.} « Elle a beaucoup d'expérience professionnelle. »\\
\zh{他毕业于雅温得第一大学。} \emph{Tā bìyè yú Yǎwēndé dì-yī dàxué.} « Il est diplômé de l'Université de Yaoundé I. »\\
\zh{我在一家公司负责翻译。} \emph{Wǒ zài yì jiā gōngsī fùzé fānyì.} « Je suis responsable de la traduction dans une entreprise. »\\
\zh{我的爱好是足球和音乐。} \emph{Wǒ de àihào shì zúqiú hé yīnyuè.} « Mes loisirs sont le football et la musique. »
\end{exemplebox}

\begin{attention}[工作 : nom ET verbe — le piège classique du francophone]
\zh{工作} (gōngzuò) est à la fois un \textbf{nom} (« le travail ») et un \textbf{verbe} (« travailler »). Mais attention : en chinois, on ne peut pas dire « je travaille professeur ». Pour indiquer sa profession, il faut utiliser \zh{是} (être) ou \zh{当} (exercer le métier de).\\
Ne dites jamais : \zh{我工作老师。} (incorrect)\\
Dites : \zh{我是老师。} \emph{Wǒ shì lǎoshī.} « Je suis professeur. »\\
ou : \zh{我当老师。} \emph{Wǒ dāng lǎoshī.} « Je travaille comme professeur. »\\
Autre piège : \zh{毕业} se construit avec \zh{于} + l'établissement (\zh{毕业于……大学}), jamais avec \zh{在}. Enfin, rappelez-vous : \zh{国籍} + nom de pays (\zh{我的国籍是喀麦隆。}), mais \zh{是} + gentilé (\zh{我是喀麦隆人。}).
\end{attention}

\courssub{Mots de liaison utiles : 从……到…… et 会 + verbe}

Pour rédiger les rubriques d'un CV, deux outils grammaticaux sont indispensables.

\begin{propriete}[从……到…… : « de... à... »]
Structure : \textbf{从 + point de départ + 到 + point d'arrivée}. Elle exprime une durée (de telle date à telle date) ou un parcours spatial. Dans un CV, elle sert à indiquer les périodes de formation ou d'emploi.\\
\zh{从2018年到2021年，我在杜阿拉大学学习。} \emph{Cóng 2018 nián dào 2021 nián, wǒ zài Dù'ālā dàxué xuéxí.} « De 2018 à 2021, j'ai étudié à l'Université de Douala. »\\
Attention à l'ordre : en chinois, la date commence par l'année, puis le mois, puis le jour (du plus grand au plus petit), contrairement au français.
\end{propriete}

\begin{propriete}[会 + verbe : « savoir faire »]
\zh{会} (huì) placé devant un verbe exprime une \cle{compétence acquise} : « savoir faire quelque chose ». C'est l'outil idéal pour la rubrique « compétences » (技能) du CV. Négation : \zh{不会}.\\
\zh{我会说法语和汉语。} \emph{Wǒ huì shuō Fǎyǔ hé Hànyǔ.} « Je sais parler français et chinois. »\\
\zh{他会用电脑。} \emph{Tā huì yòng diànnǎo.} « Il sait utiliser l'ordinateur. »\\
Ne pas confondre avec \zh{能} (pouvoir, être capable dans une situation donnée) ni avec \zh{可以} (avoir la permission).
\end{propriete}

\begin{popfigure}
\begin{tikzpicture}[mindmap, grow cyclic, every node/.style={concept, align=center, font=\small}, concept color=popblue!40, level 1 concept/.append style={level distance=3.4cm, sibling angle=72}, text=white]
\node[concept color=popblue, text=white, align=center]{\zh{简历}\\ jiǎnlì\\ \emph{CV}}
  child[concept color=poporange!80] { node[align=center]{\zh{身份}\\ identité\\ \zh{姓名 性别}\\ \zh{出生 国籍}} }
  child[concept color=popgreen!80] { node[align=center]{\zh{联系}\\ coordonnées\\ \zh{电话 邮箱}\\ \zh{地址}} }
  child[concept color=poppurple!80] { node[align=center]{\zh{教育}\\ formation\\ \zh{毕业 大学}\\ \zh{专业}} }
  child[concept color=poppink!80] { node[align=center]{\zh{经验}\\ expérience\\ \zh{工作 公司}\\ \zh{负责}} }
  child[concept color=popgold!90] { node[align=center]{\zh{技能}\\ compétences\\ \zh{翻译 会}\\ \zh{爱好}} };
\end{tikzpicture}
\legende{Carte mentale lexicale : les cinq rubriques du CV chinois et leur vocabulaire.}
\end{popfigure}

\courssub{Mini-texte de réinvestissement}

\begin{textebox}[Une candidate camerounaise]
\zh{玛丽是喀麦隆人，她出生于布埃亚。从2019年到2022年，她在雅温得第一大学学习，专业是汉语。毕业以后，她在一家中国公司工作，负责翻译。她会说法语、英语和汉语，也会用电脑。她的爱好是读书和踢足球。现在她想申请华为公司的工作，所以她写了一份中文简历。简历上有她的姓名、性别、国籍、电话、邮箱和地址。}\\
\emph{Mǎlì shì Kāmàilóng rén, tā chūshēng yú Bù'āiyà. Cóng 2019 nián dào 2022 nián, tā zài Yǎwēndé dì-yī dàxué xuéxí, zhuānyè shì Hànyǔ. Bìyè yǐhòu, tā zài yì jiā Zhōngguó gōngsī gōngzuò, fùzé fānyì. Tā huì shuō Fǎyǔ, Yīngyǔ hé Hànyǔ, yě huì yòng diànnǎo. Tā de àihào shì dúshū hé tī zúqiú. Xiànzài tā xiǎng shēnqǐng Huáwéi gōngsī de gōngzuò, suǒyǐ tā xiěle yí fèn Zhōngwén jiǎnlì. Jiǎnlì shang yǒu tā de xìngmíng, xìngbié, guójí, diànhuà, yóuxiāng hé dìzhǐ.}\\
« Marie est Camerounaise, elle est née à Buea. De 2019 à 2022, elle a étudié à l'Université de Yaoundé I, avec une spécialité en chinois. Après son diplôme, elle a travaillé dans une entreprise chinoise, où elle était responsable de la traduction. Elle sait parler français, anglais et chinois, et sait aussi utiliser l'ordinateur. Ses loisirs sont la lecture et le football. Maintenant, elle voudrait postuler à un emploi chez Huawei, c'est pourquoi elle a rédigé un CV en chinois. Sur le CV figurent son nom complet, son sexe, sa nationalité, son téléphone, son e-mail et son adresse. »
\end{textebox}

Questions rapides : 1) \zh{玛丽出生在哪里？} 2) \zh{她的专业是什么？} 3) \zh{她会说什么语言？} 4) \zh{她想申请什么公司的工作？}

Réponses attendues : 1) \zh{她出生于布埃亚。} 2) \zh{她的专业是汉语。} 3) \zh{她会说法语、英语和汉语。} 4) \zh{她想申请华为公司的工作。}

\courssub{Exercice résolu et entraînement}

\begin{exoresolu}[Compléter avec le mot du glossaire qui convient]
Énoncé — Complétez chaque phrase avec le mot approprié : 简历 / 国籍 / 专业 / 负责 / 爱好 / 申请.\\
1. 我的\_\_\_\_是喀麦隆。 2. 他在大学学习，\_\_\_\_是经济。 3. 我想\_\_\_\_这个工作。 4. 她在公司\_\_\_\_翻译。 5. 我的\_\_\_\_是打篮球。 6. 找工作的时候，你要写一份\_\_\_\_。
\tcblower
Solution détaillée — 1. \zh{国籍} (guójí) : \zh{国籍} se construit avec le nom du pays (\zh{喀麦隆}), jamais avec le gentilé ; on aurait pu dire aussi \zh{我是喀麦隆人。} 2. \zh{专业} (zhuānyè) : la filière étudiée à l'université. 3. \zh{申请} (shēnqǐng) : « postuler à cet emploi », verbe suivi de \zh{工作}. 4. \zh{负责} (fùzé) : « être responsable de la traduction ». 5. \zh{爱好} (àihào) : le basketball est un loisir. 6. \zh{简历} (jiǎnlì) : on rédige un CV pour chercher un emploi ; notez le classificateur \zh{份} utilisé pour les documents.
\end{exoresolu}

\begin{exercice}[À vous de jouer]
1. Traduisez en chinois : « Je suis né(e) à Douala. Je suis diplômé(e) de l'Université de Douala. Je sais parler français et chinois. Je voudrais postuler à cet emploi. »\\
2. Classez les mots suivants dans les cinq rubriques de la carte mentale : 邮箱、毕业、公司、爱好、性别、翻译、地址、专业、经验、电话。
\end{exercice}

\begin{corrige}[Exercice « À vous de jouer »]
1. \zh{我出生于杜阿拉。我毕业于杜阿拉大学。我会说法语和汉语。我想申请这个工作。} \emph{Wǒ chūshēng yú Dù'ālā. Wǒ bìyè yú Dù'ālā dàxué. Wǒ huì shuō Fǎyǔ hé Hànyǔ. Wǒ xiǎng shēnqǐng zhège gōngzuò.}\\
2. Identité : 性别. Coordonnées : 邮箱、地址、电话. Formation : 毕业、专业. Expérience : 公司、经验. Compétences et loisirs : 翻译、爱好.
\end{corrige}

\begin{savaistu}[Le CV en chinois : un atout au Cameroun]
Au Cameroun, de nombreuses entreprises chinoises — comme Huawei, ZTE ou les grandes sociétés de BTP qui construisent routes et stades — emploient du personnel camerounais. Pour ces employeurs, un CV rédigé en chinois est un \textbf{vrai atout à l'embauche} : il prouve immédiatement votre niveau de langue et votre motivation. En Chine, le CV (简历) est traditionnellement très bref : une seule page, avec parfois une photo d'identité, et les informations personnelles (âge, sexe) y figurent encore couramment, contrairement à l'usage français actuel.
\end{savaistu}

\begin{aretenir}[L'essentiel du glossaire]
\begin{itemize}
\item Le CV se dit \zh{简历} (jiǎnlì), classificateur \zh{份} ; ses rubriques : identité (\zh{姓名、性别、出生、国籍}), coordonnées (\zh{电话、邮箱、地址}), formation (\zh{教育、毕业、大学、专业}), expérience (\zh{经验、工作、公司、负责}), compétences et loisirs (\zh{技能、翻译、爱好}).
\item \zh{毕业于} + établissement ; \zh{负责} + domaine ; \zh{申请} + poste ; \zh{国籍} + nom de pays.
\item Périodes : \zh{从……到……} ; compétences : \zh{会} + verbe.
\item Jamais \zh{我工作老师} : dites \zh{我是老师} ou \zh{我当老师}.
\end{itemize}
\end{aretenir}

\coursec{Clés (radicaux) et ordre des traits}

Dans la section précédente, nous avons découvert le glossaire du CV : \zh{简历} (\emph{jiǎnlì}, « CV »), \zh{毕业} (\emph{bìyè}, « être diplômé »), \zh{经验} (\emph{jīngyàn}, « expérience »), \zh{职业} (\emph{zhíyè}, « profession »). Ces mots contiennent des caractères nouveaux qu'il faut maintenant apprendre à \cle{écrire} correctement. En chinois, on n'écrit pas un caractère « au feeling » : chaque caractère possède une \cle{clé} (radical) qui donne une indication de sens, et un \cle{ordre des traits} précis. C'est ce que nous allons observer sur cinq caractères du thème.

\courssub{Observation : à quoi sert une clé ?}

\begin{exemplebox}[Observer avant d'écrire]
Regardez ces trois caractères :\\
\zh{简历} \emph{jiǎnlì} « CV » — \zh{笔} \emph{bǐ} « stylo » — \zh{第一} \emph{dì yī} « premier »\\
Que remarquez-vous ? Les caractères \zh{简}, \zh{笔} et \zh{第} portent tous, en haut, la même clé \zh{⺮}, la clé du \cle{bambou}. Autrefois, en Chine, on écrivait sur des lamelles de bambou : les mots liés à l'écriture et aux documents portent donc souvent cette clé. Un CV (\zh{简历}), un stylo (\zh{笔}), un numéro d'ordre (\zh{第}) : tout se tient !
\end{exemplebox}

\begin{propriete}[La clé ⺮ (bambou)]
La clé \zh{⺮} (forme en tête de caractère de \zh{竹}, \emph{zhú}, « bambou ») apparaît dans de nombreux mots liés aux \cle{documents} et aux \cle{objets anciens en bambou} : \zh{简} (document, simple), \zh{笔} (stylo, pinceau), \zh{第} (préfixe des numéros d'ordre), \zh{等} (\emph{děng}, « attendre ; catégorie »), \zh{答} (\emph{dá}, « répondre »). Retenir la clé, c'est deviner le \cle{domaine de sens} du caractère avant même de le connaître.
\end{propriete}

\courssub{La règle générale de l'ordre des traits}

\begin{propriete}[Les trois grands principes]
Pour écrire n'importe quel caractère chinois, on applique trois principes fondamentaux :
\begin{enumerate}
\item \cle{De haut en bas} : on écrit d'abord la partie supérieure (ex. la clé \zh{⺮} de \zh{简}).
\item \cle{De gauche à droite} : on écrit d'abord la partie gauche (ex. \zh{马} dans \zh{验}).
\item \cle{L'extérieur avant l'intérieur} : pour les caractères à enveloppe, on trace d'abord le contour, puis l'intérieur, et on « ferme la porte » à la fin.
\end{enumerate}
À ces principes s'ajoutent des règles secondaires : l'horizontal (\zh{横}) avant le vertical (\zh{竖}) quand ils se croisent, le trait central avant les côtés dans certains caractères symétriques.
\end{propriete}

Contrairement au français, où l'on écrit des lettres de gauche à droite sur une ligne, le chinois compose chaque caractère dans un \cle{carré imaginaire} : il faut donc gérer l'espace en deux dimensions. C'est pourquoi l'ordre des traits n'est pas un détail : il garantit un caractère équilibré, lisible et rapide à tracer.

\courssub{Étude des cinq caractères du thème}

\textbf{1. 历 (lì) — « parcours, traverser », dans \zh{经历} et \zh{简历}}\par
Décomposition : clé \zh{厂} (\emph{hǎn}, « falaise », l'enveloppe extérieure) + \zh{力} (\emph{lì}, « force »). Nombre de traits : \cle{4}.\\
Ordre des traits : 1) \zh{横} (horizontal en haut) ; 2) \zh{撇} (le grand tombant à gauche, la « falaise ») ; 3) \zh{横折钩} (horizontal-crochet de \zh{力}) ; 4) \zh{撇} (tombant final de \zh{力}). On applique ici la règle « l'extérieur avant l'intérieur » : la falaise \zh{厂} d'abord, la force \zh{力} ensuite.

\textbf{2. 简 (jiǎn) — « simple ; document », dans \zh{简历}}\par
Décomposition : clé \zh{⺮} (bambou, en haut) + \zh{间} (\emph{jiān}, « entre », en bas). Nombre de traits : \cle{13}.\\
On écrit \cle{d'abord la clé du bambou en haut} (6 traits : trois séries « tombant-horizontal » suivies d'un point), puis \zh{间} en bas : la porte \zh{门} d'abord (extérieur : 3 traits), puis le soleil \zh{日} à l'intérieur (4 traits). C'est l'application parfaite des deux règles : haut avant bas, extérieur avant intérieur.

\textbf{3. 毕 (bì) — « terminer », dans \zh{毕业}}\par
Décomposition : \zh{比} (\emph{bǐ}, « comparer ») en haut + \zh{十} (« dix ») en bas. Nombre de traits : \cle{6}.\\
Ordre : on écrit \cle{les deux 匕} de la partie supérieure d'abord (le \zh{匕} de gauche : horizontal-crochet puis vertical ; puis le \zh{匕} de droite), puis le \zh{十} en bas : l'horizontal \zh{横} avant le vertical \zh{竖}. Erreur classique : commencer par le \zh{十} — non, on respecte « de haut en bas ».

\textbf{4. 验 (yàn) — « vérifier, éprouver », dans \zh{经验}}\par
Décomposition : clé \zh{马} (\emph{mǎ}, « cheval ») à gauche + \zh{佥} (\emph{qiān}) à droite. Nombre de traits : \cle{10}.\\
On écrit \cle{de gauche à droite} : d'abord le cheval \zh{马} (3 traits : horizontal-crochet double, horizontal, vertical), puis la partie droite \zh{佥} (7 traits, de haut en bas : le \zh{人} en haut (2 traits), l'horizontal \zh{横}, puis la partie inférieure \zh{又} (3 traits)). La clé du cheval, ici, n'indique plus le sens mais rappelle l'ancien usage (vérifier un cheval avant de l'acheter !) ; la partie droite donne le son.

\textbf{5. 职 (zhí) — « fonction, métier », dans \zh{职业}}\par
Décomposition : clé \zh{耳} (\emph{ěr}, « oreille ») à gauche + \zh{只} (\emph{zhǐ}) à droite. Nombre de traits : \cle{11}.\\
Gauche d'abord : l'oreille \zh{耳} (6 traits : horizontal, vertical, vertical, horizontal, horizontal, horizontal du bas). Droite ensuite : \zh{只} (5 traits : la bouche \zh{口} en haut — 3 traits : vertical, horizontal-crochet, horizontal —, puis les deux traits \zh{八} en bas : tombant gauche, tombant droit).

\begin{attention}[Homophones et caractères proches]
\begin{itemize}
\item Ne confondez pas \zh{历} (\emph{lì}, « parcours », dans \zh{简历}, \zh{经历}) et \zh{厉} (\emph{lì}, « sévère », dans \zh{厉害}, « redoutable »). Même prononciation, même enveloppe \zh{厂}, mais l'intérieur diffère : \zh{力} (force) contre \zh{万} (dix mille). \zh{简历} ne s'écrit jamais avec \zh{厉} !
\item Ne confondez pas \zh{毕} (\emph{bì}, « terminer ») et \zh{华} (\emph{huá}, « splendide ; Chine », dans \zh{华人}, « Chinois de souche »). Les deux contiennent \zh{十} en bas, mais le haut diffère : \zh{比} pour \zh{毕}, \zh{化} pour \zh{华}. \zh{毕业} (être diplômé) n'a rien à voir avec \zh{华}.
\end{itemize}
\end{attention}

\begin{center}
\begin{tabularx}{\linewidth}{|X|X|X|X|}
\hline
\rowcolor{popblueL} Caractère & Décomposition (clé + partie) & Traits & Ordre essentiel \\
\hline
\zh{历} lì & \zh{厂} (falaise) + \zh{力} & 4 & 横 → 撇 → 横折钩 → 撇 \\
\hline
\zh{简} jiǎn & \zh{⺮} (bambou) + \zh{间} & 13 & clé ⺮ (6 traits) → 门 → 日 \\
\hline
\zh{毕} bì & \zh{比} + \zh{十} & 6 & 比 (2 匕) → 十 (横 puis 竖) \\
\hline
\zh{验} yàn & \zh{马} (cheval) + \zh{佥} & 10 & gauche 马 → droite 佥 (人 → 横 → 又) \\
\hline
\zh{职} zhí & \zh{耳} (oreille) + \zh{只} & 11 & 耳 (6 traits) → 只 (口 → 八) \\
\hline
\end{tabularx}
\end{center}

\begin{popfigure}
\begin{tikzpicture}[font=\small]
% Décomposition de 简
\node[draw=popblue, fill=popblueL, rounded corners, minimum width=1.6cm, minimum height=1.6cm, font=\Huge] (jian) at (0,0) {简};
\node[draw=popgreen, fill=popgreenL, rounded corners, minimum width=1.1cm, minimum height=1.1cm, font=\huge] (zhu) at (2.8,0.9) {⺮};
\node[draw=poporange, fill=poporangeL, rounded corners, minimum width=1.1cm, minimum height=1.1cm, font=\huge] (jian2) at (2.8,-0.9) {间};
\draw[-{Stealth}, thick, popgreen] (jian.east) -- (zhu.west) node[midway, above, font=\footnotesize, text=popdark, align=center]{clé : bambou\\(6 traits)};
\draw[-{Stealth}, thick, poporange] (jian.east) -- (jian2.west) node[midway, below, font=\footnotesize, text=popdark, align=center]{partie phonétique\\间 (门 + 日)};
% Traits de 历
\node[draw=poppurple, fill=poppurpleL, rounded corners, minimum width=1.6cm, minimum height=1.6cm, font=\Huge] (li) at (6.2,0) {历};
\node[font=\footnotesize, text=popdark, align=left] at (9.6,0) {1. 横\quad 2. 撇\\3. 横折钩\quad 4. 撇};
\draw[-{Stealth}, thick, poppurple] (li.east) -- (8.0,0.35);
\draw[-{Stealth}, thick, poppurple] (li.east) -- (8.0,-0.35);
\node[font=\footnotesize, text=poppurple, align=center] at (6.2,-1.3) {extérieur 厂\\ avant intérieur 力};
\end{tikzpicture}
\legende{Décomposition de \zh{简} (clé + partie phonétique) et ordre des 4 traits de \zh{历}.}
\end{popfigure}

\begin{methode}[Mémoriser un caractère en trois étapes]
\begin{enumerate}
\item \cle{Identifier le radical} : il donne le domaine de sens (\zh{⺮} → bambou, documents ; \zh{耳} → oreille, écoute ; \zh{马} → ici, simple indice historique).
\item \cle{Identifier la partie phonétique} : elle suggère la prononciation (\zh{间} \emph{jiān} dans \zh{简} \emph{jiǎn} ; \zh{力} \emph{lì} dans \zh{历} \emph{lì} ; \zh{佥} \emph{qiān} dans \zh{验} \emph{yàn} ; \zh{只} \emph{zhǐ} dans \zh{职} \emph{zhí}).
\item \cle{Écrire 5 fois en comptant les traits} à voix haute : « un, deux, trois… » (\zh{一、二、三…}), en respectant l'ordre. La mémoire de la main aide la mémoire de la tête !
\end{enumerate}
\end{methode}

\begin{exoresolu}[Compter et ordonner les traits]
Énoncé : pour le caractère \zh{毕} (\emph{bì}), indiquez le nombre de traits et l'ordre général d'écriture, puis recopiez-le cinq fois sur votre ardoise.
\tcblower
Solution détaillée : \zh{毕} compte \cle{6 traits}. On écrit de haut en bas : d'abord la partie supérieure \zh{比}, c'est-à-dire le \zh{匕} de gauche (horizontal-crochet, puis vertical), puis le \zh{匕} de droite (horizontal-crochet, puis vertical) ; ensuite la partie inférieure \zh{十}, horizontal \zh{横} avant vertical \zh{竖}. Vérification : 2 + 2 + 2 = 6 traits. Piège : ne pas commencer par le \zh{十} et ne pas confondre avec \zh{华} (\emph{huá}), dont le haut est \zh{化}.
\end{exoresolu}

\begin{experience}[Écriture guidée : du tableau à l'ardoise]
Le professeur écrit au tableau, trait par trait, les cinq caractères \zh{历、简、毕、验、职} en annonçant chaque trait en chinois (\zh{横、竖、撇、横折钩…}). Les élèves lèvent le doigt et « écrivent dans l'air » en même temps, puis reproduisent chaque caractère cinq fois sur l'ardoise ou le cahier, en comptant les traits à voix haute. En binôme, chacun dicte ensuite un caractère à son voisin, qui doit l'écrire et nommer sa clé. Objectif : écrire sans faute le mot complet \zh{简历} en fin de séance.
\end{experience}

\begin{aretenir}[L'essentiel de la section]
Un caractère chinois se construit selon trois principes : de haut en bas, de gauche à droite, l'extérieur avant l'intérieur. La \cle{clé} (radical) renseigne sur le sens — \zh{⺮} (bambou) pour les documents comme \zh{简}, \zh{笔}, \zh{第} — et la partie restante donne souvent le son. Pour le thème du CV, retenez : \zh{历} (4 traits), \zh{简} (13 traits), \zh{毕} (6 traits), \zh{验} (10 traits), \zh{职} (11 traits), et méfiez-vous des sosies \zh{厉} et \zh{华}.
\end{aretenir}

\coursec{Compréhension du texte et collocations}

Dans la section précédente, nous avons appris à écrire correctement les caractères du thème en respectant les clés et l'ordre des traits. Savoir écrire ne suffit pas : il faut maintenant \cle{comprendre} le texte du CV de Li Ming dans le détail et maîtriser les \cle{collocations}, c'est-à-dire les associations de mots qui reviennent toujours ensemble dans le monde du travail (\zh{找工作}, \zh{毕业于}, \zh{工作经验}…). C'est cette double compétence — lire et réutiliser — qui sera évaluée.

\courssub{Relecture du texte support : le CV de Li Ming}

Avant de répondre aux questions, relisons attentivement le texte support de la leçon. Chaque phrase contient une information qui sera testée.

\begin{textebox}[简历 — Le CV de Li Ming]
个人简历\\
Gèrén jiǎnlì\\
« Curriculum vitæ »\\[0.4em]
我叫李明，是中国人。我1998年出生在北京。\\
Wǒ jiào Lǐ Míng, shì Zhōngguó rén. Wǒ 1998 nián chūshēng zài Běijīng.\\
« Je m'appelle Li Ming, je suis chinois. Je suis né en 1998 à Pékin. »\\[0.4em]
2023年，我毕业于北京大学，我的专业是汉语国际教育。\\
2023 nián, wǒ bìyè yú Běijīng Dàxué, wǒ de zhuānyè shì Hànyǔ guójì jiàoyù.\\
« En 2023, j'ai été diplômé de l'Université de Pékin ; ma spécialité est l'enseignement international du chinois. »\\[0.4em]
从2023年到2024年，我在一家公司工作，负责翻译工作。\\
Cóng 2023 nián dào 2024 nián, wǒ zài yì jiā gōngsī gōngzuò, fùzé fānyì gōngzuò.\\
« De 2023 à 2024, j'ai travaillé dans une entreprise, où j'étais responsable de la traduction. »\\[0.4em]
我会说三种语言：汉语、英语和法语。\\
Wǒ huì shuō sān zhǒng yǔyán : Hànyǔ, Yīngyǔ hé Fǎyǔ.\\
« Je sais parler trois langues : le chinois, l'anglais et le français. »\\[0.4em]
现在我想在喀麦隆找工作，申请一份教汉语的工作。\\
Xiànzài wǒ xiǎng zài Kāmàilóng zhǎo gōngzuò, shēnqǐng yí fèn jiāo Hànyǔ de gōngzuò.\\
« Maintenant, je voudrais chercher un emploi au Cameroun et postuler pour un poste de professeur de chinois. »
\end{textebox}

\begin{definition}[Glossaire de rappel]
\begin{itemize}
\item \zh{简历} (jiǎnlì) : CV, résumé
\item \zh{出生} (chūshēng) : naître
\item \zh{毕业} (bìyè) : être diplômé
\item \zh{专业} (zhuānyè) : spécialité, filière
\item \zh{负责} (fùzé) : être responsable de
\item \zh{申请} (shēnqǐng) : postuler, faire une demande
\item \zh{份} (fèn) : classificateur des documents, postes, journaux
\end{itemize}
\end{definition}

\courssub{Méthode : répondre aux questions de compréhension}

\begin{methode}[Trouver la réponse dans un texte chinois en quatre étapes]
\begin{enumerate}
\item \textbf{Repérer le mot interrogatif} : \zh{什么} (quoi, quel), \zh{哪} (quel, lequel), \zh{什么时候} (quand), \zh{几} (combien, petit nombre), \zh{哪里} (où). Chaque mot interrogatif annonce le \emph{type} d'information attendue.
\item \textbf{Repérer le mot-clé de la question dans le texte} : si la question contient \zh{毕业} (diplôme), cherchez ce même caractère dans le texte.
\item \textbf{Localiser la phrase entière} qui contient l'information, puis la relire avec le pinyin.
\item \textbf{Rédiger la réponse en phrase complète}, en recopiant la structure du texte et en remplaçant seulement ce qui doit l'être (je $\rightarrow$ il, etc.). Une réponse d'un seul mot est acceptée à l'oral, mais la phrase complète est exigée à l'écrit.
\end{enumerate}
\end{methode}

\courssub{Questions de compréhension}

Appliquons la méthode aux cinq questions de la fiche de progression.

\begin{exoresolu}[Questions de compréhension sur le CV de Li Ming]
Énoncé — Répondez en chinois, par des phrases complètes :\\
1. 李明是哪国人？(Lǐ Míng shì nǎ guó rén ? — De quel pays est Li Ming ?)\\
2. 他什么时候毕业的？(Tā shénme shíhou bìyè de ? — Quand a-t-il obtenu son diplôme ?)\\
3. 他的专业是什么？(Tā de zhuānyè shì shénme ? — Quelle est sa spécialité ?)\\
4. 他在哪里工作过？(Tā zài nǎlǐ gōngzuò guo ? — Où a-t-il travaillé ?)\\
5. 他会说几种语言？(Tā huì shuō jǐ zhǒng yǔyán ? — Combien de langues sait-il parler ?)
\tcblower
Solution détaillée —\\
1. Le mot interrogatif est \zh{哪} (quel pays) ; le texte dit \zh{是中国人}. Réponse : 李明是中国人。(Lǐ Míng shì Zhōngguó rén.) — Li Ming est chinois.\\
2. Le mot interrogatif est \zh{什么时候} (quand) ; le texte dit \zh{2023年，我毕业于北京大学}. Réponse : 他2023年毕业的。(Tā 2023 nián bìyè de.) — Il a été diplômé en 2023.\\
3. Le mot interrogatif est \zh{什么} ; le texte dit \zh{我的专业是汉语国际教育}. Réponse : 他的专业是汉语国际教育。(Tā de zhuānyè shì Hànyǔ guójì jiàoyù.) — Sa spécialité est l'enseignement international du chinois.\\
4. Le mot interrogatif est \zh{哪里} (où) ; le texte dit \zh{我在一家公司工作}. Réponse : 他在一家公司工作过。(Tā zài yì jiā gōngsī gōngzuò guo.) — Il a travaillé dans une entreprise.\\
5. Le mot interrogatif est \zh{几} (combien) ; le texte dit \zh{我会说三种语言}. Réponse : 他会说三种语言。(Tā huì shuō sān zhǒng yǔyán.) — Il sait parler trois langues.\\
Remarque : dans la question 4 et la réponse 4, le caractère \zh{过} (guo) marque l'\emph{expérience vécue} : « avoir déjà travaillé (au moins une fois) ».
\end{exoresolu}

\begin{exercice}[Vrai ou faux ?]
Lisez chaque phrase et écrivez 对 (duì, vrai) ou 错 (cuò, faux), puis corrigez la phrase fausse.\\
1. 李明2024年毕业于北京大学。\\
2. 他会说法语。\\
3. 李明负责翻译工作。\\
4. 他想在喀麦隆找工作。
\end{exercice}

\begin{corrige}[Exercice : Vrai ou faux]
1. 错 (cuò) — Faux. Le texte dit \zh{2023年，我毕业于北京大学} : il a été diplômé en \textbf{2023}, pas en 2024. Phrase corrigée : 李明2023年毕业于北京大学。\\
2. 对 (duì) — Vrai. Le texte dit \zh{我会说三种语言：汉语、英语和法语} : le français fait bien partie de ses langues.\\
3. 对 (duì) — Vrai : \zh{负责翻译工作}.\\
4. 对 (duì) — Vrai : \zh{我想在喀麦隆找工作}.
\end{corrige}

\courssub{Les collocations clés du monde du travail}

Une \cle{collocation} est une association stable de mots : en français on dit « chercher un emploi » et non « trouver un emploi » quand on est en recherche ; le chinois fonctionne de la même façon. Voici les six collocations du thème, observées d'abord dans le texte, puis réutilisées.

\begin{center}
\begin{tabularx}{\linewidth}{|X|X|X|X|}
\hline
\rowcolor{popblueL} Caractères & Pinyin & Nature & Traduction \\
\hline
找工作 & zhǎo gōngzuò & verbe + nom & chercher un emploi \\
\hline
工作经验 & gōngzuò jīngyàn & locution nominale & expérience professionnelle \\
\hline
毕业于 & bìyè yú & verbe + préposition & être diplômé de \\
\hline
申请工作 & shēnqǐng gōngzuò & verbe + nom & postuler (à un emploi) \\
\hline
会说……语 & huì shuō … yǔ & verbe + verbe + nom & savoir parler telle langue \\
\hline
负责…… & fùzé … & verbe & être responsable de \\
\hline
\end{tabularx}
\end{center}

\begin{exemplebox}[Chaque collocation en contexte]
\zh{他正在找工作。} — \emph{Tā zhèngzài zhǎo gōngzuò.} — « Il est en train de chercher un emploi. »\\[0.3em]
\zh{她有两年的工作经验。} — \emph{Tā yǒu liǎng nián de gōngzuò jīngyàn.} — « Elle a deux ans d'expérience professionnelle. »\\[0.3em]
\zh{我毕业于雅温得第一大学。} — \emph{Wǒ bìyè yú Yǎwēndé Dì-yī Dàxué.} — « Je suis diplômé de l'Université de Yaoundé I. »\\[0.3em]
\zh{他想申请这份工作。} — \emph{Tā xiǎng shēnqǐng zhè fèn gōngzuò.} — « Il voudrait postuler à ce poste. » (notez le classificateur \zh{份})\\[0.3em]
\zh{她会说英语和法语。} — \emph{Tā huì shuō Yīngyǔ hé Fǎyǔ.} — « Elle sait parler anglais et français. »\\[0.3em]
\zh{他负责销售工作。} — \emph{Tā fùzé xiāoshòu gōngzuò.} — « Il est responsable des ventes. »
\end{exemplebox}

\begin{exemplebox}[La durée avec 从……到……]
\zh{他从2023年到2024年在一家公司工作。} — \emph{Tā cóng 2023 nián dào 2024 nián zài yì jiā gōngsī gōngzuò.} — « Il a travaillé dans une entreprise de 2023 à 2024. »\\[0.3em]
\zh{我们从星期一到星期五上课。} — \emph{Wǒmen cóng xīngqīyī dào xīngqīwǔ shàngkè.} — « Nous avons cours du lundi au vendredi. »\\[0.3em]
\zh{从杜阿拉到雅温得有多远？} — \emph{Cóng Dù'ālā dào Yǎwēndé yǒu duō yuǎn ?} — « Quelle distance y a-t-il de Douala à Yaoundé ? »
\end{exemplebox}

\begin{attention}[N'oubliez pas le 到 final !]
La structure \zh{从……到……} (cóng … dào …, « de … à … ») \textbf{encadre} la durée ou le trajet : le mot \zh{从} ouvre le cadre et \zh{到} le ferme. Les francophones, influencés par « de 2023 à 2024 » où « à » est un petit mot discret, oublient très souvent le \zh{到} final et écrivent \zh{他从2023年2024年工作}, phrase incorrecte. Retenez : pas de \zh{从} sans \zh{到}, comme pas de parenthèse ouvrante sans parenthèse fermante. La structure se place \emph{avant} le verbe (et avant le complément de lieu \zh{在…} s'il y en a un).
\end{attention}

\begin{propriete}[会 + verbe : la compétence acquise]
\textbf{Structure} : Sujet + 会 (huì) + verbe.\\
\textbf{Emploi} : \zh{会} exprime une \cle{compétence apprise}, un savoir-faire : \zh{会说汉语} (savoir parler chinois), \zh{会开车} (savoir conduire), \zh{会游泳} (savoir nager). La négation est \zh{不会} : \zh{我不会说德语。} (Wǒ bú huì shuō Déyǔ. — Je ne sais pas parler allemand.)\\
\textbf{Différence avec 能} (néng) : \zh{能} exprime une \emph{capacité ponctuelle} ou une \emph{possibilité due aux circonstances}. Comparez :\\
\zh{他会说汉语。} — Il sait parler chinois (compétence acquise).\\
\zh{他今天能来。} — Il peut venir aujourd'hui (les circonstances le permettent).\\
\textbf{Exception} : pour les langues, on utilise toujours \zh{会}, jamais \zh{能} : on dit \zh{我会说法语}, et non \zh{我能说法语} pour parler de la compétence.
\end{propriete}

\courssub{Les étapes de la recherche d'emploi}

Le CV de Li Ming suit en réalité un parcours logique, le même au Cameroun et en Chine : on cherche un emploi, on rédige son CV, on postule, puis on passe un entretien. La figure suivante résume ce parcours avec le vocabulaire chinois correspondant.

\begin{popfigure}
\begin{center}
\begin{tikzpicture}[node distance=0.9cm]
\node[draw=popblue, fill=popblueL, rounded corners, minimum width=2.6cm, minimum height=1cm, align=center] (a) {找工作\\ \emph{zhǎo gōngzuò}\\ chercher un emploi};
\node[draw=poporange, fill=poporangeL, rounded corners, minimum width=2.6cm, minimum height=1cm, align=center, right=of a] (b) {写简历\\ \emph{xiě jiǎnlì}\\ rédiger un CV};
\node[draw=popgreen, fill=popgreenL, rounded corners, minimum width=2.6cm, minimum height=1cm, align=center, right=of b] (c) {申请工作\\ \emph{shēnqǐng gōngzuò}\\ postuler};
\node[draw=poppurple, fill=poppurpleL, rounded corners, minimum width=2.6cm, minimum height=1cm, align=center, right=of c] (d) {面试\\ \emph{miànshì}\\ passer un entretien};
\draw[-{Stealth}, very thick, popdark] (a) -- (b);
\draw[-{Stealth}, very thick, popdark] (b) -- (c);
\draw[-{Stealth}, very thick, popdark] (c) -- (d);
\end{tikzpicture}
\end{center}
\legende{Les quatre étapes de la recherche d'emploi : du besoin initial à l'entretien.}
\end{popfigure}

\begin{aretenir}[Le parcours en une phrase]
\zh{我先找工作，然后写简历，申请工作，最后去面试。} — \emph{Wǒ xiān zhǎo gōngzuò, ránhòu xiě jiǎnlì, shēnqǐng gōngzuò, zuìhòu qù miànshì.} — « Je cherche d'abord un emploi, ensuite je rédige mon CV, je postule, et enfin je vais à l'entretien. » Notez les marqueurs chronologiques \zh{先} (d'abord), \zh{然后} (ensuite), \zh{最后} (enfin), très utiles pour raconter un parcours.
\end{aretenir}

\courssub{Mini-dialogue oral : parler de ses langues}

\begin{experience}[Jeu de rôle : l'entretien d'embauche]
Par groupes de deux, l'un joue le recruteur, l'autre le candidat. Utilisez le modèle suivant, puis remplacez les langues par les vôtres.\\[0.4em]
A : 你会说什么语言？— \emph{Nǐ huì shuō shénme yǔyán ?} — « Quelles langues sais-tu parler ? »\\
B : 我会说法语、英语和汉语。— \emph{Wǒ huì shuō Fǎyǔ, Yīngyǔ hé Hànyǔ.} — « Je sais parler français, anglais et chinois. »\\
A : 你有工作经验吗？— \emph{Nǐ yǒu gōngzuò jīngyàn ma ?} — « As-tu de l'expérience professionnelle ? »\\
B : 有，我在一家公司工作过。— \emph{Yǒu, wǒ zài yì jiā gōngsī gōngzuò guo.} — « Oui, j'ai travaillé dans une entreprise. »\\[0.4em]
Attention à la prononciation : \zh{语} (yǔ, 3\textsuperscript{e} ton) dans les noms de langues, et à l'énumération avec la virgule chinoise \zh{、} (dunhao), réservée aux listes de mots parallèles.
\end{experience}

\begin{exoresolu}[Traduction guidée : de Yaoundé à Pékin]
Énoncé — Traduisez en chinois : « Marie est diplômée de l'Université de Yaoundé I. Elle sait parler français et chinois. Maintenant elle cherche un emploi à Douala. »
\tcblower
Solution détaillée — On découpe phrase par phrase et on choisit les collocations du jour :\\
1. « être diplômé de » $\rightarrow$ \zh{毕业于} : 玛丽毕业于雅温得第一大学。(Mǎlì bìyè yú Yǎwēndé Dì-yī Dàxué.)\\
2. « savoir parler » $\rightarrow$ \zh{会说} (compétence, pas \zh{能}) : 她会说法语和汉语。(Tā huì shuō Fǎyǔ hé Hànyǔ.)\\
3. « chercher un emploi » $\rightarrow$ \zh{找工作} ; le lieu se place avant le verbe avec \zh{在} : 现在她在杜阿拉找工作。(Xiànzài tā zài Dù'ālā zhǎo gōngzuò.)\\
Piège évité : ne dites pas \zh{她找工作在杜阿拉} — en chinois, le complément de lieu avec \zh{在} se place \emph{avant} le verbe, contrairement au français où « à Douala » vient après.
\end{exoresolu}

\begin{attention}[Les trois pièges des francophones dans cette leçon]
\begin{enumerate}
\item \textbf{Ordre des mots} : le lieu (\zh{在杜阿拉}) et la durée (\zh{从2023年到2024年}) se placent \emph{avant} le verbe, jamais après.
\item \textbf{毕业于} : la préposition \zh{于} (yú) est inséparable de \zh{毕业} quand on nomme l'établissement ; on ne peut pas dire \zh{我毕业北京大学}.
\item \textbf{会 / 能} : pour les langues et les savoir-faire, toujours \zh{会} ; réservez \zh{能} aux possibilités du moment.
\end{enumerate}
\end{attention}

Vous savez désormais lire un CV chinois, répondre à des questions de compréhension en phrases complètes et réemployer les six collocations clés du monde du travail. Dans la prochaine section, vous passerez à la production guidée : la rédaction de votre propre mini-CV en chinois.

\coursec{Production guidée : mon mini-CV}

Dans la section précédente, nous avons vérifié notre compréhension du CV de Li Ming et appris les collocations essentielles (\zh{申请工作} « postuler à un emploi », \zh{会说汉语} « savoir parler chinois », \zh{教育经历} « formation »). Il est maintenant temps de passer de la lecture à l'écriture : vous allez rédiger votre propre mini-CV en chinois. C'est la compétence finale de cette leçon : \cle{réutiliser le vocabulaire et les structures observés pour produire un texte personnel, correct et présenté selon les normes chinoises}.

\courssub{Observer avant d'écrire : le modèle de présentation personnelle}

Avant de remplir le gabarit, observons comment un élève camerounais peut se présenter en quelques phrases simples. Lisez d'abord le texte, puis repérez les structures que vous connaissez déjà.

\begin{textebox}[Modèle de présentation personnelle]
我叫玛丽，我是喀麦隆人。我住在雅温得。\\
Wǒ jiào Mǎlì, wǒ shì Kāmàilóng rén. Wǒ zhù zài Yǎwēndé.\\
Je m'appelle Marie, je suis Camerounaise. J'habite à Yaoundé.\\[4pt]
我在雅温得第二中学学习，今年我上毕业班。\\
Wǒ zài Yǎwēndé Dì-èr Zhōngxué xuéxí, jīnnián wǒ shàng bìyè bān.\\
J'étudie au lycée de Yaoundé II, cette année je suis en classe de Terminale.\\[4pt]
我会说法语、英语和汉语。我的爱好是踢足球和听音乐。\\
Wǒ huì shuō Fǎyǔ, Yīngyǔ hé Hànyǔ. Wǒ de àihào shì tī zúqiú hé tīng yīnyuè.\\
Je parle français, anglais et chinois. Mes loisirs sont le football et la musique.\\[4pt]
我想申请一份翻译的工作。\\
Wǒ xiǎng shēnqǐng yī fèn fānyì de gōngzuò.\\
Je voudrais postuler à un emploi de traducteur.
\end{textebox}

\begin{definition}[Les quatre phrases modèles à réutiliser]
\begin{itemize}
\item \zh{我是喀麦隆人。} (Wǒ shì Kāmàilóng rén.) — Je suis Camerounais(e). Structure : Sujet + 是 + nationalité + 人.
\item \zh{我在……中学学习。} (Wǒ zài ... zhōngxué xuéxí.) — J'étudie au lycée de... Structure : Sujet + 在 + lieu + verbe.
\item \zh{我会说法语、英语和汉语。} (Wǒ huì shuō Fǎyǔ, Yīngyǔ hé Hànyǔ.) — Je sais parler français, anglais et chinois. Structure : Sujet + 会 + 说 + langues.
\item \zh{我的爱好是……} (Wǒ de àihào shì ...) — Mon loisir est... Structure : Possesseur + 的 + 爱好 + 是 + activité.
\end{itemize}
\end{definition}

\begin{attention}[喀麦隆 ou 喀麦隆人 ?]
\zh{喀麦隆} (Kāmàilóng) désigne le \textbf{pays}, le Cameroun. Pour dire « je suis Camerounais(e) », il faut ajouter \zh{人} (rén, « personne ») : \zh{我是喀麦隆人。} Dire \zh{我是喀麦隆} reviendrait à dire « je suis le Cameroun », ce qui est une erreur fréquente des francophones, car en français « je suis camerounais » n'a pas de mot supplémentaire. Même logique : \zh{中国人} (Chinois), \zh{法国人} (Français), \zh{美国人} (Américain).
\end{attention}

\courssub{Le gabarit du CV chinois : les neuf rubriques}

Un CV chinois (简历, jiǎnlì) se présente comme un tableau à deux colonnes : à gauche la \cle{rubrique}, à droite l'\cle{information}. Voici les neuf rubriques que vous devez connaître et savoir remplir.

\begin{center}
\begin{tabularx}{\linewidth}{|X|X|X|X|}
\hline
\rowcolor{popblueL} Caractères & Pinyin & Nature & Traduction \\
\hline
姓名 & xìngmíng & nom & nom et prénom \\
\hline
性别 & xìngbié & nom & sexe (男 homme / 女 femme) \\
\hline
出生日期 & chūshēng rìqī & locution nominale & date de naissance \\
\hline
国籍 & guójí & nom & nationalité \\
\hline
电话 & diànhuà & nom & téléphone \\
\hline
邮箱 & yóuxiāng & nom & adresse électronique \\
\hline
教育经历 & jiàoyù jīnglì & locution nominale & formation, parcours scolaire \\
\hline
技能 & jìnéng & nom & compétences \\
\hline
爱好 & àihào & nom & loisirs, passe-temps \\
\hline
\end{tabularx}
\end{center}

\begin{popfigure}
\begin{tikzpicture}[
  rub/.style={fill=popblueL, draw=popline, minimum width=3.4cm, minimum height=0.85cm, anchor=west, font=\small},
  cel/.style={fill=white, draw=popline, minimum width=6.6cm, minimum height=0.85cm, anchor=west, font=\small},
  tit/.style={font=\bfseries}
]
\node[tit] at (5,0.7) {个人简历 Gèrén jiǎnlì — CV};
\foreach \i/\txt in {0/姓名 xìngmíng, 1/性别 xìngbié, 2/出生日期 chūshēng rìqī, 3/国籍 guójí, 4/电话 diànhuà, 5/邮箱 yóuxiāng, 6/教育经历 jiàoyù jīnglì, 7/技能 jìnéng, 8/爱好 àihào}{
  \node[rub] at (0,-0.9*\i) {\txt};
  \node[cel] at (3.4,-0.9*\i) {};
}
\node[anchor=west, text=popmuted, font=\itshape\small] at (0,-8.6) {Gabarit vierge : recopiez-le et complétez la colonne de droite avec vos informations.};
\end{tikzpicture}
\legende{Gabarit vierge d'un CV chinois à compléter (rubriques à gauche, informations personnelles à droite).}
\end{popfigure}

\begin{propriete}[Comment remplir chaque rubrique]
\begin{itemize}
\item \zh{姓名} : écrivez votre nom en caractères phonétiques ou en lettres latines, par exemple \zh{玛丽} (Mǎlì, Marie).
\item \zh{性别} : une seule réponse possible : \zh{男} (nán, masculin) ou \zh{女} (nǚ, féminin).
\item \zh{出生日期} : en chinois, la date se lit \textbf{du plus grand au plus petit} : année + 年 + mois + 月 + jour + 日. Exemple : \zh{2007年3月15日} (èr líng líng qī nián sān yuè shíwǔ rì) — le 15 mars 2007. C'est l'inverse de l'ordre français (jour/mois/année).
\item \zh{国籍} : \zh{喀麦隆} (Cameroun), sans 人 ici, car la rubrique demande le pays, pas la personne.
\item \zh{教育经历} : nom du lycée + classe, par exemple \zh{雅温得第二中学，毕业班} (lycée de Yaoundé II, Terminale).
\item \zh{技能} : les langues d'abord : \zh{法语、英语、汉语} ; on peut ajouter \zh{电脑} (diànnǎo, informatique).
\item \zh{爱好} : verbe + complément : \zh{踢足球} (jouer au football), \zh{听音乐} (écouter de la musique), \zh{看书} (lire), \zh{打篮球} (jouer au basket).
\end{itemize}
\end{propriete}

\begin{methode}[Rédiger son mini-CV en quatre étapes]
\begin{enumerate}
\item \textbf{Compléter d'abord les rubriques faciles} : nom (姓名), sexe (性别), nationalité (国籍), téléphone (电话). Ce sont des mots isolés, sans phrase à construire.
\item \textbf{Rédiger la date de naissance} en respectant l'ordre chinois : année → mois → jour, avec 年, 月, 日.
\item \textbf{Rédiger deux ou trois phrases complètes} pour les rubriques 教育经历, 技能 et 爱好, en recopiant les modèles : \zh{我在……中学学习。} / \zh{我会说法语、英语和汉语。} / \zh{我的爱好是……}
\item \textbf{Relire} : vérifier les caractères (traits manquants ?), le pinyin (tons placés sur la bonne voyelle ?) et les rubriques (rien d'oublié ?).
\end{enumerate}
\end{methode}

\courssub{Exercice résolu : le CV d'André, élève de Terminale à Douala}

\begin{exoresolu}[Compléter un gabarit]
Énoncé. André Mbarga est un élève de Terminale A au lycée de Douala. Il est né le 2 juin 2007. Il parle français, anglais et chinois ; il aime jouer au football et lire. Complétez les rubriques 国籍, 出生日期, 技能 et 爱好 de son CV, puis rédigez une phrase de présentation complète.
\tcblower
Solution détaillée.
\begin{itemize}
\item 国籍 : \zh{喀麦隆} (Kāmàilóng) — le pays, sans 人, car la rubrique demande la nationalité.
\item 出生日期 : \zh{2007年6月2日} (èr líng líng qī nián liù yuè èr rì) — ordre année → mois → jour.
\item 技能 : \zh{法语、英语、汉语} (Fǎyǔ, Yīngyǔ, Hànyǔ).
\item 爱好 : \zh{踢足球、看书} (tī zúqiú, kàn shū).
\end{itemize}
Phrase de présentation : \zh{我叫安德烈，我是喀麦隆人。我在杜阿拉中学学习。我会说法语、英语和汉语。我的爱好是踢足球和看书。} (Wǒ jiào Āndéliè, wǒ shì Kāmàilóng rén. Wǒ zài Dù'ālā zhōngxué xuéxí. Wǒ huì shuō Fǎyǔ, Yīngyǔ hé Hànyǔ. Wǒ de àihào shì tī zúqiú hé kàn shū.) — « Je m'appelle André, je suis Camerounais. J'étudie au lycée de Douala. Je parle français, anglais et chinois. Mes loisirs sont le football et la lecture. »
\end{exoresolu}

\begin{attention}[Les trois pièges du CV des francophones]
\begin{enumerate}
\item \textbf{L'ordre de la date} : ne jamais écrire \zh{2日6月2007年} (calque du français). En chinois : \zh{2007年6月2日}, du plus grand au plus petit.
\item \textbf{会 + verbe} : pour « je parle chinois », on dit \zh{我会说汉语} ; 会 (huì) exprime la compétence acquise par l'apprentissage. Sans 会, la phrase perd l'idée de « savoir-faire ».
\item \textbf{La virgule de liste} : entre des mots d'une énumération, le chinois utilise \zh{、} (dunhao) et non la virgule \zh{，} : \zh{法语、英语和汉语}. Le mot \zh{和} (hé, « et ») ne se place qu'entre les deux derniers éléments.
\end{enumerate}
\end{attention}

\courssub{Complément utile à l'examen : la lettre de motivation}

\begin{definition}[La formule d'ouverture 我想申请……]
Pour ouvrir une lettre de motivation (求职信, qiúzhíxìn), la formule standard est :\\
\zh{我想申请……} (Wǒ xiǎng shēnqǐng ...) — « Je voudrais postuler à... »\\
Structure : Sujet + 想 (vouloir) + 申请 (postuler) + complément.\\
Exemples : \zh{我想申请这份工作。} (Wǒ xiǎng shēnqǐng zhè fèn gōngzuò.) — Je voudrais postuler à cet emploi. \zh{我想申请翻译的工作。} (Wǒ xiǎng shēnqǐng fānyì de gōngzuò.) — Je voudrais postuler à un emploi de traducteur. Notez le classificateur \zh{份} (fèn) devant 工作.
\end{definition}

\begin{savaistu}[Trois langues, un atout majeur]
Maîtriser le français, l'anglais et le chinois est un atout considérable sur le marché du travail camerounais et international. Le Cameroun est officiellement bilingue (français et anglais), et la Chine est devenue l'un de ses principaux partenaires économiques : entreprises de construction, télécommunications, commerce. Les traducteurs, interprètes et employés bilingues ou trilingues sont très recherchés. Votre CV trilingue, commencé aujourd'hui en classe, est déjà une vraie compétence professionnelle.
\end{savaistu}

\courssub{À vous : remplissage guidé, relecture en binôme et mise en commun}

\begin{experience}[Mon mini-CV en trois temps]
\begin{enumerate}
\item \textbf{Production individuelle (10 min)} : recopiez le gabarit vierge et complétez les neuf rubriques avec vos informations personnelles (votre ville : Yaoundé, Douala, Buea... ; votre lycée ; vos vraies langues et loisirs). Ajoutez au bas du CV une phrase avec \zh{我想申请……}.
\item \textbf{Relecture en binôme (5 min)} : échangez votre feuille avec votre voisin. Vérifiez trois points : les \textbf{caractères} (aucun trait manquant, 喀麦隆 bien écrit), les \textbf{tons} du pinyin (Kāmàilóng : 1-4-2), les \textbf{rubriques} (les neuf sont-elles remplies ? la date est-elle dans l'ordre chinois ?). Signalez les erreurs avec bienveillance et proposez une correction.
\item \textbf{Mise en commun (5 min)} : quelques volontaires lisent leur CV à voix haute. La classe écoute, puis la correction collective se fait au tableau : on relève une ou deux erreurs par CV, jamais plus, et on félicite les réussites (belle écriture, tons corrects, phrase 我想申请 bien construite).
\end{enumerate}
\end{experience}

\begin{exercice}[Pour s'entraîner à la maison]
Rédigez le mini-CV complet d'un personnage imaginaire : une élève de Terminale à Buea, née le 20 octobre 2007, qui parle anglais, français et chinois, aime jouer au basket et écouter de la musique, et veut postuler à un emploi de guide touristique (导游, dǎoyóu). Utilisez les neuf rubriques et au moins trois phrases modèles.
\end{exercice}

\begin{corrige}[Pour s'entraîner à la maison]
Éléments de correction : 国籍 : \zh{喀麦隆} ; 出生日期 : \zh{2007年10月20日} (èr líng líng qī nián shí yuè èrshí rì) ; 技能 : \zh{英语、法语、汉语} ; 爱好 : \zh{打篮球、听音乐}. Phrase modèle : \zh{我在布埃亚中学学习。我会说英语、法语和汉语。我想申请导游的工作。} (Wǒ zài Bù'āiyà zhōngxué xuéxí. Wǒ huì shuō Yīngyǔ, Fǎyǔ hé Hànyǔ. Wǒ xiǎng shēnqǐng dǎoyóu de gōngzuò.) — « J'étudie au lycée de Buea. Je parle anglais, français et chinois. Je voudrais postuler à un emploi de guide touristique. »
\end{corrige}

\begin{aretenir}[L'essentiel de la production]
Un CV chinois se présente en tableau à deux colonnes avec neuf rubriques clés : 姓名, 性别, 出生日期, 国籍, 电话, 邮箱, 教育经历, 技能, 爱好. Trois règles d'or : la date se lit année → mois → jour ; « je suis Camerounais » se dit \zh{我是喀麦隆人} avec 人 ; les énumérations utilisent \zh{、} et \zh{和}. Avec les quatre phrases modèles et la formule \zh{我想申请……}, vous savez désormais vous présenter et postuler en chinois : une compétence concrète pour l'examen et pour la vie professionnelle.
\end{aretenir}

\newpage

\coursec{Activité d'intégration}

\coursec{Leçon 38 — Vocabulaire et compréhension de texte : rédaction d'un CV}

\courssub{Mise en route et découverte du thème}

\begin{experience}[Activation des connaissances : « Qu'est-ce qu'un bon CV ? »]
En classe entière, le professeur projette ou affiche deux CV simplifiés en chinois (un clair et structuré, un brouillon désordonné). Les élèves identifient en binôme les rubriques indispensables (\zh{姓名}, \zh{学历}, \zh{技能}, \zh{爱好}…) et les erreurs à éviter (ordre des mots, classificateurs, tons). Mise en commun au tableau : construction collective du \cle{gabarit type} d'un CV chinois pour lycéen.
\end{experience}

\begin{aretenir}[Objectifs de la leçon]
À l'issue de cette leçon, tu seras capable de :
\begin{itemize}
\item \cle{lire et comprendre} un CV standard en chinois (niveau HSK 3) et une annonce de recrutement ;
\item \cle{maîtriser le vocabulaire} du recrutement et de la présentation professionnelle (identité, formation, langues, compétences, disponibilité) ;
\item \cle{reconnaître et écrire} 8 à 10 caractères clés du thème (radicaux, ordre des traits) ;
\item \cle{répondre à des questions de compréhension} explicites et implicites sur un texte de candidature ;
\item \cle{rédiger ton propre mini-CV} en chinois (15--20 lignes) en respectant la structure, l'ordre des mots (lieu avant verbe) et les mots-outils (\zh{会}, \zh{喜欢}, \zh{的}, \zh{得}) ;
\item \cle{te présenter oralement} en situation d'entretien simulé (job dating) avec une prononciation tonale correcte.
\end{itemize}
\end{aretenir}

\courssub{Texte support : le CV de Li Ming et l'annonce de Zhongya Telecom}

\begin{textebox}[简历 — CV de Li Ming (modèle de référence)]
简历\\[2mm]
Jiǎnlì\\[1mm]
« CV »\\[3mm]
姓名：李明 \quad 年龄：二十岁 \quad 性别：男\\
Xìngmíng : Lǐ Míng \quad Niánlíng : èrshí suì \quad Xìngbié : nán\\
« Nom : Li Ming \quad Âge : 20 ans \quad Sexe : Masculin »\\[2mm]
住址：北京市海淀区 \quad 电话：138-1234-5678 \quad 邮箱：liming@email.cn\\
Zhùzhǐ : Běijīng shì Hǎidiàn qū \quad Diànhuà : 138-1234-5678 \quad Yóuxiāng : liming@email.cn\\
« Adresse : District de Haidian, Pékin \quad Tél : ... \quad Courriel : ... »\\[3mm]
\textbf{教育背景} (Jiàoyù bèijǐng — « Formation ») :\\
2018--2022 \quad 北京大学 \quad 中文系 \quad 本科\\
Èr líng yī bā -- èr líng èr èr \quad Běijīng Dàxué \quad Zhōngwén xì \quad běnkē\\
« 2018--2022 \quad Université de Pékin \quad Département de chinois \quad Licence »\\[2mm]
\textbf{工作经验} (Gōngzuò jīngyàn — « Expérience professionnelle ») :\\
2021 \quad 中非贸易公司 \quad 实习生\\
Èr líng èr yī \quad Zhōng-Fēi Màoyì Gōngsī \quad shíxíshēng\\
« 2021 \quad Société de commerce Sino-Africaine \quad Stagiaire »\\[2mm]
\textbf{语言技能} (Yǔyán jìnéng — « Compétences linguistiques ») :\\
汉语 (母语) 、 英语 (流利) 、 法语 (良好)\\
Hànyǔ (mǔyǔ) 、 Yīngyǔ (liúlì) 、 Fǎyǔ (liánghǎo)\\
« Chinois (langue maternelle), Anglais (courant), Français (bon niveau) »\\[2mm]
\textbf{专业技能} (Zhuānyè jìnéng — « Compétences techniques ») :\\
熟练使用办公软件 (Word, Excel, PPT)，会开车。\\
Shúliàn shǐyòng bàngōng ruǎnjiàn (Word, Excel, PPT), huì kāichē.\\
« Maîtrise des logiciels bureautiques (Word, Excel, PPT), sait conduire. »\\[2mm]
\textbf{爱好与特长} (Àihào yǔ tècháng — « Loisirs et talents ») :\\
喜欢阅读、踢足球、帮助别人。周末有时间。\\
Xǐhuan yuèdú, tī zúqiú, bāngzhù biérén. Zhōumò yǒu shíjiān.\\
« Aime la lecture, le football, aider les autres. Disponible le week-end. »\\[2mm]
\textbf{自我评价} (Zìwǒ píngjià — « Auto-évaluation ») :\\
勤奋、诚实、有团队精神。\\
Qínfèn, chéngshí, yǒu tuánduì jīngshén.\\
« Travailleur, honnête, esprit d'équipe. »
\end{textebox}

\begin{textebox}[招聘启事 — Annonce de recrutement Zhongya Telecom]
招聘启事\\[2mm]
Zhāopìn qǐshì\\[1mm]
« Annonce de recrutement »\\[3mm]
中亚电信公司招聘兼职员工。\\
Zhōngyà Diànxìn Gōngsī zhāopìn jiānzhí yuángōng.\\
« La société Zhongya Telecom recrute des employés à temps partiel. »\\[3mm]
要求：会说汉语和法语，喜欢帮助别人，周末有时间。\\
Yāoqiú : huì shuō Hànyǔ hé Fǎyǔ, xǐhuan bāngzhù biérén, zhōumò yǒu shíjiān.\\
« Profil recherché : savoir parler chinois et français, aimer aider les autres, être disponible le week-end. »\\[3mm]
请把您的简历发到：zhongya@douala.com\\
Qǐng bǎ nín de jiǎnlì fā dào : zhongya@douala.com\\
« Veuillez envoyer votre CV à : zhongya@douala.com »
\end{textebox}

\courssub{Glossaire du thème : vocabulaire du CV et du recrutement}

\begin{center}
\begin{tabularx}{\linewidth}{|X|X|X|X|}
\hline
\rowcolor{popblueL} \textbf{Caractères} & \textbf{Pinyin} & \textbf{Nature} & \textbf{Traduction} \\ \hline
简历 & jiǎnlì & nom & CV, curriculum vitæ \\ \hline
招聘 & zhāopìn & verbe & recruter, embaucher \\ \hline
启事 & qǐshì & nom & avis, annonce, communiqué \\ \hline
公司 & gōngsī & nom & société, entreprise \\ \hline
兼职 & jiānzhí & adj./adv. & à temps partiel \\ \hline
全职 & quánzhí & adj./adv. & à temps plein \\ \hline
员工 & yuángōng & nom & employé, personnel \\ \hline
要求 & yāoqiú & verbe/nom & exiger / exigence, critère \\ \hline
专业 & zhuānyè & nom/adj. & spécialité, majeur / professionnel \\ \hline
学历 & xuélì & nom & niveau d'études, diplôme \\ \hline
经验 & jīngyàn & nom & expérience \\ \hline
能力 & nénglì & nom & capacité, compétence \\ \hline
技能 & jìnéng & nom & compétence technique, savoir-faire \\ \hline
母语 & mǔyǔ & nom & langue maternelle \\ \hline
流利 & liúlì & adj. & courant (langue) \\ \hline
良好 & liánghǎo & adj. & bon, satisfaisant \\ \hline
熟练 & shúliàn & adj. & maîtriser, être habile \\ \hline
办公软件 & bàngōng ruǎnjiàn & nom & logiciels bureautiques \\ \hline
实习生 & shíxíshēng & nom & stagiaire \\ \hline
自荐信 & zìjiànxìn & nom & lettre de motivation \\ \hline
面试 & miànshì & verbe/nom & passer un entretien / entretien \\ \hline
\end{tabularx}
\end{center}

\begin{propriete}[Collocations fréquentes (搭配) — À apprendre par cœur]
\begin{center}
\begin{tabularx}{\linewidth}{|X|X|X|}
\hline
\rowcolor{poporangeL} \textbf{Chinois (Pinyin)} & \textbf{Français} & \textbf{Exemple contextuel} \\ \hline
\zh{发简历} (fā jiǎnlì) & envoyer un CV & \zh{请发简历到邮箱。} (Qǐng fā jiǎnlì dào yóuxiāng.) \\ \hline
\zh{找工作} (zhǎo gōngzuò) & chercher du travail & \zh{我在杜阿拉找工作。} (Wǒ zài Dù'ālā zhǎo gōngzuò.) \\ \hline
\zh{参加面试} (cānjiā miànshì) & passer un entretien & \zh{下周我参加面试。} (Xià zhōu wǒ cānjiā miànshì.) \\ \hline
\zh{专业对口} (zhuānyè duìkǒu) & correspondre à la spécialité & \zh{这个工作专业对口。} (Zhège gōngzuò zhuānyè duìkǒu.) \\ \hline
\zh{薪资待遇} (xīnzī dàiyù) & salaire et avantages & \zh{薪资待遇面议。} (Xīnzī dàiyù miànyì.) \\ \hline
\end{tabularx}
\end{center}
\end{propriete}

\courssub{Clés (radicaux) et ordre des traits : 8 caractères clés}

\begin{definition}[Rappel méthodologique : ordre des traits de base]
Règles fondamentales : \textbf{1. Horizontal avant vertical} (十) \textbf{2. Gauche avant droite} (川) \textbf{3. Haut avant bas} (三) \textbf{4. Extérieur avant intérieur} (同) \textbf{5. Fermer en dernier} (回) \textbf{6. Trait central avant les ailes} (小) \textbf{7. Point en haut à gauche en premier, en haut à droite en dernier} (门).
\end{definition}

\begin{center}
\begin{tabularx}{\linewidth}{|c|X|X|X|}
\hline
\rowcolor{popgreenL} \textbf{Caractère} & \textbf{Radical (Clé)} & \textbf{Nombre de traits} & \textbf{Ordre des traits (description) \& Décomposition} \\ \hline
\zh{简} & \zh{竹} (zhú, bambou) & 13 & Haut (bambou) → Bas (简 sans bambou : 间). Le bambou s'écrit : 6 traits (haut→bas, gauche→droite). \\ \hline
\zh{历} & \zh{止} (zhǐ, s'arrêter) / \zh{厂} (hǎn, falaise) & 8 & \zh{厂} (large haut→bas gauche→bas droit) → \zh{止} (horizontal, vertical, crochet, horizontal, horizontal). *Note : clé moderne souvent \zh{止}*. \\ \hline
\zh{招} & \zh{扌} (shǒu, main) & 8 & Main (3 traits : horizontal, crochet vertical, point) → \zh{召} (haut→bas : 刀 + 口 + 刂). \\ \hline
\zh{聘} & \zh{言} (yán, parole) & 14 & \zh{言} (7 traits : point, horizontaux, vertical, horizontaux) → \zh{并} (haut→bas : 一, 八, 一, 一). \\ \hline
\zh{专} & \zh{寸} (cùn, pouce/unité) & 4 & \zh{叀} (haut→bas : 一, 丨, 丨, 一) → \zh{寸} (horizontal, crochet vertical, point, horizontal). \\ \hline
\zh{业} & \zh{一} (yī, un) / \zh{木} (mù, bois) & 5 & Version simplifiée : \zh{一} → \zh{丨} → \zh{丿} → \zh{丶} → \zh{一} (base). *Clé traditionnelle \zh{木}*. \\ \hline
\zh{语} & \zh{言} (yán, parole) & 14 & \zh{言} (7 traits) → \zh{吾} (7 traits : horizontal, vertical, horizontal, horizontal, vertical, horizontal, horizontal). \\ \hline
\zh{简} & \zh{竹} (zhú) & 13 & (Rappel) Le radical \zh{竹} évoque les tablettes de bambou pour écrire. \\ \hline
\end{tabularx}
\end{center}

\begin{exercice}[Entraînement à l'écriture]
Sur ton cahier d'écriture (纸格子), écris chaque caractère ci-dessus \textbf{5 fois} en respectant l'ordre des traits. Surline le radical en rouge. Dicte à ton voisin : « \zh{请写“招聘”} » / « \zh{请写“专业”} ».
\end{exercice}

\courssub{Compréhension du texte et collocations}

\begin{exoresolu}[Questions de compréhension sur le CV de Li Ming]
Énoncé : Lis le CV de Li Ming (texte support 1) et réponds en phrases complètes en chinois (caractères + pinyin).
\tcblower
\textbf{1. 李明多大了？} — \zh{他二十岁。} (Tā èrshí suì.)\\
\textbf{2. 他在哪里上大学？} — \zh{他在北京大学上大学。} (Tā zài Běijīng Dàxué shàng dàxué.)\\
\textbf{3. 他的专业是什么？} — \zh{他的专业是中文。} (Tā de zhuānyè shì Zhōngwén.)\\
\textbf{4. 他会说哪些语言？} — \zh{他会说法语、英语和汉语。} (Tā huì shuō Fǎyǔ, Yīngyǔ hé Hànyǔ.)\\
\textbf{5. 他有什么工作经验？} — \zh{他在中非贸易公司做过实习生。} (Tā zài Zhōng-Fēi Màoyì Gōngsī zuò guò shíxíshēng.)\\
\textbf{6. Quels sont ses loisirs ?} — \zh{他喜欢阅读、踢足球、帮助别人。} (Tā xǐhuan yuèdú, tī zúqiú, bāngzhù biérén.)\\
\textbf{7. Est-il disponible le week-end pour le poste de Zhongya Telecom ?} — \zh{是的，他周末有时间。} (Shì de, tā zhōumò yǒu shíjiān.)\\
\textbf{8. Correspond-il au profil ?} — \zh{是的，完全符合要求。} (Shì de, wánquán fúhé yāoqiú.)
\end{exoresolu}

\begin{exercice}[Compréhension de l'annonce — Vrai ou Faux ? Justifie en chinois.]
\begin{enumerate}
\item L'entreprise s'appelle « Zhongya Telecom ». (\zh{中亚电信公司})
\item Le poste est à temps plein (\zh{全职}).
\item Il faut savoir parler anglais.
\item Il faut aimer aider les autres (\zh{喜欢帮助别人}).
\item Il faut envoyer le CV par la poste.
\end{enumerate}
\end{exercice}

\begin{corrige}[Exercice Compréhension annonce]
\begin{enumerate}
\item \textbf{Vrai.} \zh{公司叫中亚电信公司。}
\item \textbf{Faux.} \zh{是兼职，不是全职。} (Shì jiānzhí, bù shì quánzhí.)
\item \textbf{Faux.} \zh{要求会说法语和汉语，没提英语。} (Yāoqiú huì shuō Fǎyǔ hé Hànyǔ, méi tí Yīngyǔ.)
\item \textbf{Vrai.} \zh{要求喜欢帮助别人。}
\item \textbf{Faux.} \zh{发到邮箱。} (Fā dào yóuxiāng.)
\end{enumerate}
\end{corrige}

\courssub{Production guidée : mon mini-CV (Job Dating Zhongya Telecom)}

\begin{methode}[Comment rédiger ton mini-CV en chinois — 5 étapes]
\begin{enumerate}
\item \textbf{Identité} : \zh{姓名} (Nom), \zh{年龄} (Âge : \zh{我今年XX岁}), \zh{住址} (Ville), \zh{学校} (Lycée). \textit{Astuce : pas de \zh{是} avant l'âge.}
\item \textbf{Formation} : \zh{我在[Ville]的[École]学习。} (Lieu avant verbe !) \zh{我的专业是[Série A/Littérature/Sciences].}
\item \textbf{Langues} : \zh{我会说[Langue1]、[Langue2]和[Langue3].} (Virgule chinoise \zh{、}). Niveau : \zh{我的[Langue]说得不错/流利/一般。} (Complément de degré \zh{得}).
\item \textbf{Compétences \& Qualités} : \zh{我会用电脑。} \zh{我喜欢帮助别人。} \zh{我周末有时间。} (Reprendre les mots de l'annonce).
\item \textbf{Conclusion} : \zh{谢谢您！} (Formule de politesse formelle \zh{您}).
\end{enumerate}
\end{methode}

\courssub{Situation : Job Dating au Lycée}

La société chinoise \emph{Zhongya Telecom} (中亚电信), installée à Bonanjo, Douala, organise un \emph{job dating} dans ton établissement. Tu es élève de Terminale A, tu apprends le chinois depuis trois ans. C'est l'occasion de valoriser tes compétences.

\begin{experience}[Job Dating en binôme — Mise en situation professionnelle]
\textbf{Étape 1 : Rédaction individuelle (15 min).} Sur feuille propre, rédige ton CV (15--20 lignes) selon le gabarit ci-dessus. Vérifie : tons, ordre des traits, classificateurs (\zh{一个公司}, \zh{三种语言}), structure \zh{在+Lieu+Verbe}.

\textbf{Étape 2 : Échange et entretien (10 min).} Échange ton CV avec ton voisin.
\begin{itemize}
\item \textbf{Rôle A (Recruteur Zhongya)} : Lis le CV en silence. Pose les deux questions obligatoires :
      \zh{你的专业是什么？} \quad \zh{你会说什么语言？}
\item \textbf{Rôle B (Candidat)} : Réponds \textbf{à voix haute, en phrases complètes, sans lire}.
\end{itemize}
Inversez les rôles. Le recruteur remplit la grille d'évaluation :
\begin{center}
\begin{tabular}{|l|c|c|}
\hline
\rowcolor{poppurpleL} \textbf{Critère} & \textbf{Oui} & \textbf{Non} \\ \hline
CV complet (4 rubriques + conclusion) & & \\ \hline
Tons corrects (auto-éval. + pair) & & \\ \hline
Réponses en phrases complètes & & \\ \hline
Ordre des mots (lieu avant verbe) & & \\ \hline
Emploi correct de \zh{会} / \zh{喜欢} & & \\ \hline
\end{tabular}
\end{center}

\textbf{Étape 3 : Sélection et affichage (5 min).} Chaque binôme désigne le « meilleur CV » (critères : exactitude, lisibilité, adéquation à l'annonce). Affichage au mur de la classe. Les lauréats présentent leur parcours devant la classe (1 min).
\end{experience}

\begin{attention}[Pièges à éviter dans ton CV — Erreurs fréquentes des francophones]
\begin{itemize}
\item \textbf{Âge :} \zh{我今年十八岁。} (Wǒ jīnnián shíbā suì.) \textbf{Pas} \zh{我是十八岁} (\zh{是} interdit avec nombre + \zh{岁}).
\item \textbf{Lieu + Verbe :} \zh{我在杜阿拉学习。} (Wǒ \textbf{zài Dù'ālā} xuéxí.) \textbf{Pas} \zh{我学习在杜阿拉。} (Ordre SVO strict : Complément de lieu avant verbe).
\item \textbf{\zh{会} vs \zh{喜欢} :} \zh{会} = compétence acquise (\zh{会说汉语}, \zh{会开车}, \zh{会用电脑}). \zh{喜欢} = goût, préférence (\zh{喜欢足球}, \zh{喜欢帮助别人}). Ne pas confondre.
\item \textbf{Classificateurs obligatoires :} \zh{一\textbf{个}公司} (yí gè gōngsī), \zh{三\textbf{种}语言} (sān zhǒng yǔyán), \zh{一\textbf{份}简历} (yí fèn jiǎnlì).
\item \textbf{Complément de degré :} \zh{说得好/不错/流利} (shuō \textbf{de} hǎo/búcuò/liúlì). Le \zh{得} est obligatoire pour évaluer la manière.
\item \textbf{Ponctuation :} Utilise la virgule de liste \zh{、} pour énumérer des noms (\zh{法语、英语和汉语}), le point final \zh{。}.
\end{itemize}
\end{attention}

\begin{aretenir}[Structures et lexique utiles — Fiche mémo pour l'entretien]
\begin{itemize}
\item \textbf{Identité :} \zh{我叫……} (Wǒ jiào…) ; \zh{我是喀麦隆人。} (Wǒ shì Kāmàilóng rén.) ; \zh{我今年十八岁。}
\item \textbf{Formation :} \zh{我在……中学学习。} (Wǒ zài … zhōngxué xuéxí.) ; \zh{我的专业是文学/科学。} (Wǒ de zhuānyè shì wénxué/kēxué.)
\item \textbf{Langues :} \zh{我会说法语、英语和汉语。} (Wǒ huì shuō Fǎyǔ, Yīngyǔ hé Hànyǔ.) ; \zh{我的汉语说得不错。} (Wǒ de Hànyǔ shuō de búcuò.)
\item \textbf{Qualités / Disponibilité :} \zh{我喜欢帮助别人。} (Wǒ xǐhuan bāngzhù biérén.) ; \zh{我会用电脑。} (Wǒ huì yòng diànnǎo.) ; \zh{我周末有时间。} (Wǒ zhōumò yǒu shíjiān.)
\item \textbf{Mots-clés recrutement :} \zh{简历} (jiǎnlì), \zh{招聘} (zhāopìn), \zh{公司} (gōngsī), \zh{兼职} (jiānzhí), \zh{面试} (miànshì), \zh{录取} (lùqǔ, être admis/embauché).
\end{itemize}
\end{aretenir}

\begin{savaistu}[CV en Chine vs Cameroun : codes culturels]
\begin{itemize}
\item \textbf{Photo :} En Chine, la photo d'identité (fond bleu/rouge, tenue formelle) est \textbf{quasi obligatoire} sur le CV. Au Cameroun, elle est recommandée mais pas toujours exigée pour la fonction publique.
\item \textbf{Informations personnelles :} Les CV chinois incluent souvent \zh{出生日期} (date de naissance), \zh{籍贯} (lieu d'origine/famille), \zh{婚姻状况} (situation matrimoniale), \zh{政治面貌} (appartenance au Parti/Jeunesses communistes). En France/Cameroun (RGPD), ces infos sont \textbf{interdites} ou déconseillées (discrimination).
\item \textbf{Longueur :} Un jeune diplômé chinois fait souvent 1 page dense. En Europe, 1 page max est la norme stricte.
\item \textbf{Réseaux :} \zh{领英} (LinkedIn) est peu utilisé en Chine continentale. On utilise \zh{前程无忧} (51job), \zh{智联招聘} (Zhaopin), \zh{BOSS直聘} (Boss Zhipin) ou \zh{微信} (WeChat) pour le networking.
\item \textbf{Conseil :} Pour postuler chez \emph{Zhongya Telecom} à Douala, adapte-toi : mets une photo pro, indique ton âge et ta disponibilité (critères locaux), mais garde la structure claire « à la chinoise » apprise en classe.
\end{itemize}
\end{savaistu}

\begin{exoresolu}[Modèle de mini-CV et d'entretien — Aïcha, Terminale A, Douala]
Énoncé : Rédige le mini-CV d'Aïcha, candidate au poste de Zhongya Telecom, puis écris ses réponses aux deux questions du recruteur.
\tcblower
\textbf{CV modèle :}\\[2mm]
简历 — Jiǎnlì — « CV »\\
姓名：艾莎 — Xìngmíng : Àishā — « Nom : Aïcha »\\
年龄：十八岁 — Niánlíng : shíbā suì — « Âge : 18 ans »\\
住址：杜阿拉 — Zhùzhǐ : Dù'ālā — « Adresse : Douala »\\[2mm]
我是喀麦隆人，在杜阿拉的一所中学学习，现在是毕业班的学生。\\
Wǒ shì Kāmàilóng rén, zài Dù'ālā de yì suǒ zhōngxué xuéxí, xiànzài shì bìyèbān de xuésheng.\\
« Je suis Camerounaise, j'étudie dans un lycée de Douala et je suis maintenant en classe de Terminale. »\\
我的专业是文学。\\
Wǒ de zhuānyè shì wénxué.\\
« Ma spécialité est la littérature (série A). »\\
我会说法语、英语和汉语。我的汉语说得不错。\\
Wǒ huì shuō Fǎyǔ, Yīngyǔ hé Hànyǔ. Wǒ de Hànyǔ shuō de búcuò.\\
« Je sais parler français, anglais et chinois. Je parle plutôt bien chinois. »\\
我喜欢帮助别人，也会用电脑。我周末有时间。\\
Wǒ xǐhuan bāngzhù biérén, yě huì yòng diànnǎo. Wǒ zhōumò yǒu shíjiān.\\
« J'aime aider les autres et je sais aussi utiliser l'ordinateur. Je suis disponible le week-end. »\\
谢谢您！— Xièxie nín ! — « Merci ! »\\[3mm]
\textbf{Réponses à l'entretien :}\\
— \zh{你的专业是什么？} — \zh{我的专业是文学。}\\
— \zh{你会说什么语言？} — \zh{我会说法语、英语和汉语。}\\[2mm]
\textbf{Commentaire des choix de langue :} L'âge est donné sans \zh{是} (\zh{十八岁}) ; le complément de lieu \zh{在杜阿拉} précède le verbe \zh{学习} ; l'énumération des langues utilise la virgule chinoise de liste \zh{、}) ; le complément de degré \zh{说得不错} évalue la manière de parler ; le CV reprend exactement les trois critères de l'annonce (langues, aider les autres, disponibilité le week-end), ce qui montre une vraie compréhension du document de départ.
\end{exoresolu}

\courssub{Bilan de la leçon}

\begin{aretenir}[Auto-évaluation — Coche ce que tu sais faire]
\begin{itemize}
\item[ ] Je sais \cle{rédiger un CV en chinois} avec les 4 rubriques du gabarit (identité, formation, langues, loisirs/qualités) et la formule de politesse finale.
\item[ ] Je sais \cle{lire le CV d'un camarade} et en dégager les informations principales (nom, âge, spécialité, langues, dispo).
\item[ ] Je sais \cle{répondre oralement} aux questions \zh{你的专业是什么？} et \zh{你会说什么语言？} en phrases complètes, sans lire.
\item[ ] Je fais attention aux \cle{tons} (ex: \zh{zhāopìn} 1-4, \zh{jiānzhí} 1-2, \zh{zhuānyè} 1-4), à l'\cle{ordre des mots} (lieu \zh{在...} avant verbe) et à l'emploi de \zh{会} (compétence) vs \zh{喜欢} (goût).
\item[ ] J'ai écrit correctement les 8 caractères clés (\zh{简历招聘专业语言}…) en respectant l'ordre des traits et identifié leurs radicaux.
\item[ ] J'ai découvert comment le chinois devient un \cle{atout professionnel concret} au Cameroun (entreprises chinoises, télécoms, BTP, commerce), et je connais les différences culturelles CV Chine / Cameroun.
\end{itemize}
\end{aretenir}

Garde ton CV affiché ou rangé dans ton portfolio : il te servira de base pour la leçon d'expression écrite (Lettre de motivation) et pour de vraies candidatures futures.

\newpage

\coursec{Méthodes et exercices résolus}

\coursec{Leçon 38 — Vocabulaire et compréhension de texte : rédaction d’un CV}

\courssub{Mise en route et découverte du thème}

\begin{textebox}[Texte support : Le CV de Li Ming]
李明，男，2004年5月12日出生，住在雅温得。2020年至2023年，他在一所中学学习，成绩很好。他会说中文、法文和英文，还会用电脑。他喜欢足球和音乐。他想找一份翻译的工作。

\emph{Lǐ Míng, nán, 2004 nián 5 yuè 12 rì chūshēng, zhù zài Yǎwēndé. 2020 nián zhì 2023 nián, tā zài yì suǒ zhōngxué xuéxí, chéngjì hěn hǎo. Tā huì shuō Zhōngwén, Fǎwén hé Yīngwén, hái huì yòng diànnǎo. Tā xǐhuan zúqiú hé yīnyuè. Tā xiǎng zhǎo yí fèn fānyì de gōngzuò.}

« Li Ming, masculin, né le 12 mai 2004, habite à Yaoundé. De 2020 à 2023, il a étudié dans un lycée, ses résultats étaient très bons. Il sait parler chinois, français et anglais, et sait aussi utiliser l'ordinateur. Il aime le football et la musique. Il veut trouver un emploi de traducteur. »
\end{textebox}

\begin{popfigure}
\begin{tikzpicture}[mindmap, grow cyclic, every node/.style=concept, concept color=popblue, level 1/.append style={level distance=4cm, sibling angle=90}, level 2/.append style={level distance=3cm, sibling angle=45}]
\node[root concept] {简历 \emph{jiǎnlì} (CV)}
child [concept color=popgreen] {node {身份 \emph{shēnfèn}} child {node {姓名 \emph{xìngmíng}}} child {node {性别 \emph{xìngbié}}} child {node {出生 \emph{chūshēng}}}}
child [concept color=poporange] {node {教育 \emph{jiàoyù}} child {node {学校 \emph{xuéxiào}}} child {node {成绩 \emph{chéngjì}}} child {node {毕业 \emph{bìyè}}}}
child [concept color=poppurple] {node {技能 \emph{jìnéng}} child {node {语言 \emph{yǔyán}}} child {node {电脑 \emph{diànnǎo}}} child {node {经验 \emph{jīngyàn}}}}
child [concept color=poppink] {node {项目 \emph{xiàngmù}} child {node {工作 \emph{gōngzuò}}} child {node {爱好 \emph{àihào}}} child {node {求职 \emph{qiúzhí}}}}
;
\end{tikzpicture}
\legende{Structure mentale d'un CV chinois (简历) : les quatre rubriques principales et leur vocabulaire clé.}
\end{popfigure}

\courssub{Glossaire du thème : Le CV (简历)}

\begin{definition}[Vocabulaire thématique — Tableau récapitulatif]
\begin{center}
\begin{tabularx}{\linewidth}{|X|X|X|X|}
\hline
\rowcolor{popblueL} \textbf{Caractères} & \textbf{Pinyin} & \textbf{Nature} & \textbf{Traduction} \\ \hline
简历 & jiǎnlì & nom & CV, curriculum vitæ \\ \hline
姓名 & xìngmíng & nom & nom et prénom \\ \hline
性别 & xìngbié & nom & sexe, genre \\ \hline
男 & nán & adj/nom & masculin, homme \\ \hline
女 & nǚ & adj/nom & féminin, femme \\ \hline
出生 & chūshēng & verbe & naître \\ \hline
日期 & rìqī & nom & date \\ \hline
住在 & zhù zài & loc. verbale & habiter à \\ \hline
雅温得 & Yǎwēndé & n. propre & Yaoundé \\ \hline
杜阿拉 & Dù'ālā & n. propre & Douala \\ \hline
布埃亚 & Bù'āiyà & n. propre & Buea \\ \hline
中学 & zhōngxué & nom & lycée, collège \\ \hline
学习 & xuéxí & verbe & étudier, apprendre \\ \hline
成绩 & chéngjì & nom & résultats, notes \\ \hline
很好 & hěn hǎo & adj & très bon \\ \hline
毕业 & bìyè & verbe & obtenir son diplôme, finir ses études \\ \hline
会 & huì & verbe aux. & savoir (compétence acquise) \\ \hline
说 & shuō & verbe & parler \\ \hline
中文 & Zhōngwén & nom & chinois (langue) \\ \hline
法文 & Fǎwén & nom & français (langue) \\ \hline
英文 & Yīngwén & nom & anglais (langue) \\ \hline
语言 & yǔyán & nom & langue \\ \hline
电脑 & diànnǎo & nom & ordinateur \\ \hline
用 & yòng & verbe & utiliser \\ \hline
还 & hái & adv. & aussi, en plus \\ \hline
喜欢 & xǐhuan & verbe & aimer \\ \hline
足球 & zúqiú & nom & football \\ \hline
音乐 & yīnyuè & nom & musique \\ \hline
找 & zhǎo & verbe & chercher \\ \hline
工作 & gōngzuò & nom/verbe & travail, emploi / travailler \\ \hline
翻译 & fānyì & nom/verbe & traduction, traducteur / traduire \\ \hline
当 & dāng & verbe & être (métier), exercer comme \\ \hline
记者 & jìzhě & nom & journaliste \\ \hline
经验 & jīngyàn & nom & expérience \\ \hline
份 & fèn & classif. & classif. pour emplois, documents \\ \hline
所 & suǒ & classif. & classif. pour écoles, hôpitaux \\ \hline
种 & zhǒng & classif. & classif. pour types, sortes \\ \hline
一点儿 & yìdiǎnr & quantif. & un peu \\ \hline
现在 & xiànzài & adv. & maintenant \\ \hline
毕业班 & bìyè bān & nom & classe de terminale, classe d'examen \\ \hline
学生 & xuésheng & nom & élève, étudiant \\ \hline
\end{tabularx}
\end{center}
\end{definition}

\courssub{Clés (radicaux) et ordre des traits}

\begin{propriete}[Clés sémantiques et ordre d'écriture des caractères clés]
\begin{center}
\begin{tabularx}{\linewidth}{|X|X|X|X|}
\hline
\rowcolor{popblueL} \textbf{Caractère} & \textbf{Clé (radical)} & \textbf{Sens de la clé} & \textbf{Ordre des traits (règles principales)} \\ \hline
\zh{语} \emph{yǔ} & \zh{讠} (yán) & parole, langage & 1. \zh{讠} (point, trait horizontal pliant) — 2. \zh{吾} (horizontal, vertical, horizontaux) \\ \hline
\zh{译} \emph{yì} & \zh{讠} (yán) & parole, langage & 1. \zh{讠} — 2. \zh{睪} (haut → bas : \zh{目}, \zh{日}, \zh{日}) \\ \hline
\zh{说} \emph{shuō} & \zh{讠} (yán) & parole, langage & 1. \zh{讠} — 2. \zh{兑} (vertical, horizontaux, crochet) \\ \hline
\zh{简} \emph{jiǎn} & \zh{⺮} (zhú) & bambou (lamelles d'écriture) & 1. \zh{⺮} (haut, gauche → droite) — 2. \zh{间} : \zh{门} (extérieur : vertical gauche, haut+droite, bas) puis \zh{日} (intérieur) \\ \hline
\zh{历} \emph{lì} & \zh{止} (zhǐ) & pied, marche & 1. \zh{厂} (haut, gauche) — 2. \zh{止} (horizontal, vertical, crochet, horizontal) \\ \hline
\zh{住} \emph{zhù} & \zh{亻} (rén) & personne & 1. \zh{亻} (traits gauche/droite) — 2. \zh{主} (point, horizontaux, vertical, horizontal) \\ \hline
\zh{注} \emph{zhù} & \zh{氵} (shuǐ) & eau & 1. \zh{氵} (trois traits : point, point, trait montant) — 2. \zh{主} \\ \hline
\zh{作} \emph{zuò} & \zh{亻} (rén) & personne & 1. \zh{亻} — 2. \zh{乍} (haut → bas) \\ \hline
\zh{想} \emph{xiǎng} & \zh{心} (xīn) & cœur, esprit & 1. \zh{相} (haut \zh{木}\zh{目}, bas \zh{心}) — 2. \zh{心} écrit en dernier (bas) \\ \hline
\zh{翻} \emph{fān} & \zh{羽} (yǔ) & plume, aile & 1. \zh{羽} (haut complexe) — 2. \zh{凡} (bas) \\ \hline
\end{tabularx}
\end{center}
\emph{Rappel des règles d'or :} haut → bas ; gauche → droite ; horizontal avant vertical ; extérieur avant intérieur ; on ferme la boîte en dernier (ex: \zh{口}, \zh{门}, \zh{间}).
\end{propriete}

\courssub{Méthode 1 : Lire et comprendre un court texte chinois (un CV)}

\begin{methode}[Lire un CV chinois en 5 étapes]
\begin{enumerate}
\item \textbf{Repérer la structure du document.} Un CV chinois (\zh{简历} \emph{jiǎnlì}) se lit de haut en bas : \zh{姓名} (nom), \zh{性别} (sexe), \zh{出生日期} (date de naissance), \zh{住址} / \zh{地址} (adresse), \zh{教育背景} (formation), \zh{工作经验} (expérience), \zh{技能} (compétences), \zh{爱好} (loisirs), \zh{求职意向} (projet professionnel). Repérer ces rubriques avant de lire les phrases.
\item \textbf{Lire une première fois sans dictionnaire.} Entourer les mots connus (\cle{mots connus}), souligner les inconnus, deviner leur sens grâce au contexte et aux \cle{clés (radicaux)} : \zh{语言} contient la clé \zh{讠} (parole) → lien avec les langues ; \zh{翻译} contient \zh{羽} (aile/plume → tourner/retourner) + \zh{讠} → « retourner les mots » = traduire.
\item \textbf{Identifier les informations factuelles.} Qui ? (\zh{李明}) Où ? (\zh{雅温得}) Quand ? (\zh{2020年}) Quoi ? (\zh{学中文}). Ce sont les réponses aux questions de compréhension.
\item \textbf{Relire phrase par phrase} en découpant selon l'ordre chinois : \cle{Sujet + Temps + Lieu + Verbe + Complément}. Ne jamais traduire mot à mot.
\item \textbf{Vérifier les tons et les mots-outils} : \zh{的} après un nom (\zh{我的简历}), \zh{了} pour une action accomplie (\zh{我学了三年中文}), \zh{过} pour l'expérience passée (\zh{我去过中国}).
\end{enumerate}
\end{methode}

\courssub{Méthode 2 : Reconnaître, prononcer et employer le vocabulaire du CV}

\begin{methode}[Mémoriser et réemployer le lexique du CV]
\begin{enumerate}
\item \textbf{Classer le vocabulaire par rubrique} : identité (\zh{姓名, 性别, 出生}), formation (\zh{学习, 成绩, 毕业, 学校}), compétences (\zh{会, 语言, 电脑, 用}), expérience et projet (\zh{工作, 经验, 找, 当, 想}).
\item \textbf{Apprendre les mots avec leur collocation} : on n'apprend pas \zh{找} seul, mais \zh{找工作} (« chercher un emploi ») ; pas \zh{经验} seul, mais \zh{有经验} (« avoir de l'expérience ») ; \zh{当记者} (« être journaliste »).
\item \textbf{Apprendre chaque mot avec son classificateur} : \zh{一份工作} (un emploi), \zh{一所学校} (une école), \zh{三种语言} (trois langues), \zh{一封信} (une lettre), \zh{一个小时} (une heure).
\item \textbf{Vérifier la prononciation} en lisant le pinyin à voix haute : \zh{简历} \emph{jiǎnlì} (3\textsuperscript{e} + 4\textsuperscript{e} ton), \zh{性别} \emph{xìngbié} (4\textsuperscript{e} + 2\textsuperscript{e}), \zh{经验} \emph{jīngyàn} (1\textsuperscript{er} + 4\textsuperscript{e}), \zh{翻译} \emph{fānyì} (1\textsuperscript{er} + 4\textsuperscript{e}).
\item \textbf{Réemployer dans une phrase personnelle} : remplacer le nom, la date, la ville par ses propres informations.
\end{enumerate}
\end{methode}

\begin{exoresolu}[Vocabulaire : compléter un CV à trous]
\textbf{Énoncé.} Complète avec le mot convenable : \zh{简历 / 出生 / 经验 / 会 / 份}

这是我的\underline{\hspace{1.5cm}}。我2005年在杜阿拉\underline{\hspace{1.5cm}}。我\underline{\hspace{1.5cm}}说法文和中文。我有一年的工作\underline{\hspace{1.5cm}}。我想找一\underline{\hspace{1.5cm}}工作。

\emph{Zhè shì wǒ de \underline{\hspace{1.5cm}}. Wǒ 2005 nián zài Dù'ālā \underline{\hspace{1.5cm}}. Wǒ \underline{\hspace{1.5cm}} shuō Fǎwén hé Zhōngwén. Wǒ yǒu yì nián de gōngzuò \underline{\hspace{1.5cm}}. Wǒ xiǎng zhǎo yí \underline{\hspace{1.5cm}} gōngzuò.}
\tcblower
\textbf{Solution détaillée.}
\begin{enumerate}
\item \underline{\hspace{1.5cm}} = \zh{简历} (\emph{jiǎnlì}) : « Ceci est mon CV ». Le contexte (tout le texte décrit une identité professionnelle) impose ce mot.
\item \underline{\hspace{1.5cm}} = \zh{出生} (\emph{chūshēng}) : après une date et un lieu (\zh{在杜阿拉}), le verbe « naître » est attendu : « Je suis né en 2005 à Douala. »
\item \underline{\hspace{1.5cm}} = \zh{会} (\emph{huì}) : devant le verbe \zh{说}, il faut l'auxiliaire de capacité acquise « savoir » : \zh{会说中文} = « savoir parler chinois ». On ne peut pas utiliser \zh{能} ici, car \zh{能} exprime la possibilité physique ou la permission, pas la compétence apprise.
\item \underline{\hspace{1.5cm}} = \zh{经验} (\emph{jīngyàn}) : la collocation figée est \zh{工作经验} (« expérience professionnelle ») : « J'ai un an d'expérience professionnelle. »
\item \underline{\hspace{1.5cm}} = \zh{份} (\emph{fèn}) : classificateur des emplois et documents : \zh{一份工作}. \zh{个} serait une faute ici.
\end{enumerate}
Texte complété : 这是我的\zh{简历}。我2005年在杜阿拉\zh{出生}。我\zh{会}说法文和中文。我有一年的工作\zh{经验}。我想找一\zh{份}工作。

\emph{Zhè shì wǒ de jiǎnlì. Wǒ 2005 nián zài Dù'ālā chūshēng. Wǒ huì shuō Fǎwén hé Zhōngwén. Wǒ yǒu yì nián de gōngzuò jīngyàn. Wǒ xiǎng zhǎo yí fèn gōngzuò.} — « Voici mon CV. Je suis né en 2005 à Douala. Je sais parler français et chinois. J'ai un an d'expérience professionnelle. Je veux trouver un emploi. »
\end{exoresolu}

\courssub{Méthode 3 : Répondre aux questions de compréhension par des phrases complètes}

\begin{methode}[Rédiger une réponse de compréhension en chinois]
\begin{enumerate}
\item \textbf{Repérer le mot interrogatif} : \zh{谁} (qui), \zh{什么} (quoi), \zh{哪里}/\zh{哪儿} (où), \zh{什么时候} (quand), \zh{几}/\zh{多少} (combien), \zh{为什么} (pourquoi), \zh{还是} (choix entre deux options).
\item \textbf{Localiser la phrase du texte} qui contient la réponse et la souligner.
\item \textbf{Recopier la structure de la question} en remplaçant le mot interrogatif par l'information : \zh{他住在哪里？} → \zh{他住在雅温得。} Le mot interrogatif disparaît, l'ordre des mots ne change pas.
\item \textbf{Pour les questions en \zh{还是}} (est-ce A ou B ?), choisir une seule option : \zh{他是男的还是女的？} → \zh{他是男的。}
\item \textbf{Relire} : vérifier les tons, la ponctuation chinoise (\zh{。}), et que la phrase contient un verbe.
\end{enumerate}
\end{methode}

\begin{exoresolu}[Compréhension de texte : le CV de Li Ming]
\textbf{Énoncé.} Lis le texte de Li Ming (ci-dessus), puis réponds aux questions en chinois.

Questions : 1) \zh{李明是男的还是女的？} 2) \zh{他住在哪里？} 3) \zh{他会说几种语言？哪几种？} 4) \zh{他想找什么工作？}
\tcblower
\textbf{Solution détaillée.}
\begin{enumerate}
\item Question 1 : le texte commence par \zh{男} (\emph{nán}, « masculin »). Réponse : \zh{李明是男的。} (\emph{Lǐ Míng shì nán de.} — « Li Ming est un homme. »)
\item Question 2 : repérer \zh{住在雅温得}. Structure : \zh{住在} + lieu = « habiter à ». Réponse : \zh{他住在雅温得。} (\emph{Tā zhù zài Yǎwēndé.} — « Il habite à Yaoundé. »)
\item Question 3 : repérer \zh{会说中文、法文和英文}. La virgule chinoise \zh{、} énumère trois langues. Réponse : \zh{他会说三种语言：中文、法文和英文。} (\emph{Tā huì shuō sān zhǒng yǔyán : Zhōngwén, Fǎwén hé Yīngwén.} — « Il sait parler trois langues : le chinois, le français et l'anglais. ») \textbf{Attention au classificateur :} \zh{一种语言}, pas \zh{一个语言}.
\item Question 4 : repérer \zh{想找一份翻译的工作}. \zh{找} = « chercher », \zh{一份工作} = « un emploi » (classificateur \zh{份}). Réponse : \zh{他想找一份翻译的工作。} (\emph{Tā xiǎng zhǎo yí fèn fānyì de gōngzuò.} — « Il veut trouver un emploi de traducteur. »)
\end{enumerate}
\textbf{Conclusion méthodologique :} chaque réponse reprend la structure de la question et cite l'information du texte ; on répond par une phrase complète, jamais par un mot isolé.
\end{exoresolu}

\begin{exoresolu}[Questions de compréhension sur un second extrait (Wang Fang)]
\textbf{Énoncé.} Texte : \zh{王芳，女，2003年出生，是杜阿拉一所中学的学生。她2023年毕业，成绩很好。她会说三种语言，还会用电脑。她想当翻译。}

\emph{Wáng Fāng, nǚ, 2003 nián chūshēng, shì Dù'ālā yì suǒ zhōngxué de xuésheng. Tā 2023 nián bìyè, chéngjì hěn hǎo. Tā huì shuō sān zhǒng yǔyán, hái huì yòng diànnǎo. Tā xiǎng dāng fānyì.}

Questions : 1) \zh{王芳什么时候毕业？} 2) \zh{她会说几种语言？} 3) \zh{她想当什么？}
\tcblower
\textbf{Solution détaillée.}
\begin{enumerate}
\item Le mot interrogatif est \zh{什么时候} (« quand »). On cherche \zh{毕业} dans le texte : \zh{她2023年毕业}. Réponse complète : \zh{她2023年毕业。} (\emph{Tā 2023 nián bìyè.} — « Elle a obtenu son diplôme en 2023. ») Le temps se place avant le verbe, sans préposition « en ».
\item Le mot interrogatif est \zh{几} (« combien », petit nombre). Le texte dit \zh{三种语言}. Réponse : \zh{她会说三种语言。} (\emph{Tā huì shuō sān zhǒng yǔyán.} — « Elle sait parler trois langues. ») Classificateur \zh{种} pour les sortes/types.
\item Le mot interrogatif est \zh{什么}. Le texte dit \zh{她想当翻译} : \zh{当} + métier = « devenir, exercer comme ». Réponse : \zh{她想当翻译。} (\emph{Tā xiǎng dāng fānyì.} — « Elle veut devenir traductrice. »)
\end{enumerate}
\textbf{Conclusion :} la réponse reprend mot pour mot la phrase du texte, en supprimant le mot interrogatif ; c'est la stratégie la plus sûre à l'examen.
\end{exoresolu}

\courssub{Méthode 4 : Identifier les clés (radicaux) pour deviner le sens des caractères}

\begin{methode}[Utiliser les clés pour lire un mot inconnu]
\begin{enumerate}
\item \textbf{Repérer la clé} (partie significative du caractère) : \zh{讠} (parole) → \zh{说, 语, 译} ; \zh{亻} (personne) → \zh{住, 作, 们} ; \zh{女} (femme) → \zh{姓, 她, 妈} ; \zh{心} (cœur) → \zh{想, 爱, 忘} ; \zh{氵} (eau) → \zh{注, 洗, 河} ; \zh{⺮} (bambou) → \zh{简, 策, 箱}.
\item \textbf{Déduire le champ de sens} : \zh{译} contient \zh{讠} → lien avec la parole ; \zh{翻译} = « traduction ».
\item \textbf{Utiliser la clé pour distinguer les caractères proches (homophones ou quasi-homophones)} : \zh{作} (\zh{亻}, faire/activité) ≠ \zh{做} (faire/fabriquer, souvent manuel) ; \zh{住} (\zh{亻}, habiter/arrêter) ≠ \zh{注} (\zh{氵}, noter/verser/concentrer) ; \zh{已} (\zh{己} déjà fini, 3\textsuperscript{e} trait ne dépasse pas) ≠ \zh{己} (soi-même, 3\textsuperscript{e} trait dépasse à gauche) ≠ \zh{巳} (signe du serpent, trait central vertical).
\item \textbf{Appliquer l'ordre des traits} pour écrire : de haut en bas, de gauche à droite, l'horizontal avant le vertical, l'extérieur avant l'intérieur, on ferme la boîte en dernier.
\end{enumerate}
\end{methode}

\begin{exoresolu}[Clés et écriture : exercice d'application]
\textbf{Énoncé.} a) Quelle est la clé des caractères \zh{语, 译, 说} ? Que suggère-t-elle ? b) Décris l'ordre des traits de \zh{简} (dans \zh{简历}). c) Distingue \zh{住} et \zh{注} par la clé et le sens.
\tcblower
\textbf{Solution détaillée.}
\begin{enumerate}
\item[a)] \zh{语} (\emph{yǔ}, langue), \zh{译} (\emph{yì}, traduire), \zh{说} (\emph{shuō}, parler) contiennent tous la clé \zh{讠}, forme simplifiée de \zh{言}, « la parole ». Ces trois caractères appartiennent donc au champ sémantique du langage : on peut deviner qu'un caractère inconnu avec \zh{讠} a un rapport avec parler, lire ou les langues.
\item[b)] \zh{简} (\emph{jiǎn}, simple ; \zh{简历} = CV) : clé \zh{竹} → \zh{⺮} (bambou) en haut, car autrefois on écrivait sur des lamelles de bambou. Ordre d'écriture : d'abord la clé \zh{⺮} en haut (de gauche à droite, les deux éléments), puis la partie \zh{间} en dessous : le cadre \zh{门} (extérieur : vertical gauche, haut+droite, bas) avant \zh{日} (intérieur), conformément à la règle « l'extérieur avant l'intérieur ».
\item[c)] \zh{住} (\emph{zhù}, habiter) porte la clé \zh{亻} (personne) : une personne qui se fixe quelque part → « habiter, résider ». \zh{注} (\emph{zhù}, noter, verser, concentrer) porte la clé \zh{氵} (eau) : verser de l'eau → « verser, noter, enregistrer ». Même prononciation, même partie phonétique \zh{主}, mais sens différents : c'est la clé de gauche qui permet de choisir le bon caractère dans \zh{住在雅温得} (habiter à Yaoundé) vs \zh{注意} (faire attention).
\end{enumerate}
\textbf{Conclusion :} la clé donne le sens, la partie phonétique donne souvent le son : c'est l'outil n° 1 pour lire un texte sans dictionnaire.
\end{exoresolu}

\courssub{Méthode 5 : Rédiger son propre mini-CV en chinois}

\begin{methode}[Produire un mini-CV en 6 lignes types]
\begin{enumerate}
\item \textbf{Ligne 1 — Identité} : \zh{我叫……，我是男的/女的。} (\emph{Wǒ jiào..., wǒ shì nán de / nǚ de.})
\item \textbf{Ligne 2 — Naissance et lieu} : \zh{我……年……月……日出生，住在……。} Utiliser \zh{出生} pour « naître » et \zh{住在} + ville. \textbf{Piège :} pas de \zh{在} devant l'année avec \zh{出生} (on dit \zh{我2005年生}, pas \zh{我在2005年生}).
\item \textbf{Ligne 3 — Formation} : \zh{我在……中学学习，现在是毕业班的学生。} (\emph{Wǒ zài ... zhōngxué xuéxí, xiànzài shì bìyè bān de xuésheng.})
\item \textbf{Ligne 4 — Compétences} : \zh{我会说……和……，还会用电脑。} Enchaîner avec \zh{还} (« de plus »). Nuancer : \zh{会说一点儿中文} (« savoir parler un peu chinois »).
\item \textbf{Ligne 5 — Loisirs} : \zh{我喜欢……和……。}
\item \textbf{Ligne 6 — Projet professionnel} : \zh{我想找一份……的工作。} ou \zh{我想当……。} Relire : un verbe par phrase, temps avant le verbe, classificateur \zh{份} devant \zh{工作}.
\end{enumerate}
\end{methode}

\begin{exoresolu}[Production guidée : rédiger un mini-CV]
\textbf{Énoncé.} Rédige en chinois un mini-CV de 5 à 6 phrases pour Marie, née en 2005 à Buea, élève en Terminale A, qui parle anglais, français et un peu de chinois, aime la musique et le football, et veut devenir journaliste.
\tcblower
\textbf{Solution détaillée.}

On suit le plan en 6 lignes de la méthode :

\zh{我叫玛丽，我是女的。我2005年在布埃亚出生。我在一所中学学习，现在是毕业班的学生。我会说英文、法文，还会说一点儿中文。我喜欢音乐和足球。我想当记者。}

\emph{Wǒ jiào Mǎlì, wǒ shì nǚ de. Wǒ 2005 nián zài Bù'āiyà chūshēng. Wǒ zài yì suǒ zhōngxué xuéxí, xiànzài shì bìyè bān de xuésheng. Wǒ huì shuō Yīngwén, Fǎwén, hái huì shuō yìdiǎnr Zhōngwén. Wǒ xǐhuan yīnyuè hé zúqiú. Wǒ xiǎng dāng jìzhě.}

« Je m'appelle Marie, je suis une fille. Je suis née en 2005 à Buea. J'étudie dans un lycée, je suis maintenant élève en classe de Terminale. Je sais parler anglais, français, et je sais aussi parler un peu chinois. J'aime la musique et le football. Je veux devenir journaliste. »

\textbf{Justifications grammaticales :}
\begin{itemize}
\item \zh{我2005年出生} : la date se place entre le sujet et le verbe, sans préposition (calque du français « je suis née \emph{en} 2005 » à éviter : pas de \zh{在} devant l'année quand \zh{出生} suit).
\item \zh{还会说一点儿中文} : \zh{还} ajoute une compétence ; \zh{一点儿} (« un peu ») nuance le niveau, placé entre le verbe et le nom (ou après le verbe si pas de nom : \zh{会说一点儿}).
\item \zh{我想当记者} : \zh{当} + nom de métier, sans article, sans classificateur.
\item \zh{布埃亚} (\emph{Bù'āiyà}) : transcription phonétique standard de Buea.
\end{itemize}
\textbf{Conclusion :} un mini-CV correct tient en six structures apprises ; la réussite vient de l'exactitude (tons, ordre des mots, classificateurs), pas de la longueur.
\end{exoresolu}

\courssub{Élément socioculturel : Le CV en Chine et au Cameroun}

\begin{savaistu}[CV chinois vs CV camerounais : points de vigilance]
\begin{itemize}
\item \textbf{Photo d'identité :} En Chine, la photo (souvent fond bleu ou rouge, tenue formelle) est quasi obligatoire sur le CV. Au Cameroun, elle est fréquente mais pas systématique selon les secteurs.
\item \textbf{Informations personnelles :} Le CV chinois standard inclut souvent l'âge (\zh{年龄}), le lieu de naissance (\zh{籍贯} — origine familiale), l'état civil (\zh{婚姻状况}), le parti politique (\zh{政治面貌} : membre du PCC, de la Ligue de la Jeunesse, masses). Au Cameroun (RGPD/loi sur la protection des données), ces infos sont déconseillées, voire illégales à demander (âge, situation matrimoniale, religion).
\item \textbf{Écriture :} En Chine, le CV manuscrit (\zh{手写简历}) était jadis exigé pour juger la calligraphie et le sérieux ; aujourd'hui, le CV numérique (Word/PDF) domine, mais la calligraphie reste valorisée. Au Cameroun, seul le CV tapé est professionnel.
\item \textbf{Longueur :} En Chine, un jeune diplômé tient son CV sur \textbf{une page A4}. Au Cameroun, 1 à 2 pages sont acceptées.
\item \textbf{Références :} En Chine, on indique souvent \zh{联系人} (personne à contacter) et \zh{电话} en bas. Au Cameroun, on écrit « Références : sur demande ».
\end{itemize}
\emph{Conseil :} Pour un stage en entreprise chinoise au Cameroun, préparez un CV bilingue (français/chinois) sur une page, avec photo, sans infos privées sensibles, et soignez l'écriture des caractères si version manuscrite demandée.
\end{savaistu}

\courssub{Activité orale : Entretien d'embauche simulé}

\begin{experience}[Jeu de rôle : « Je recrute / Je postule »]
\textbf{Consigne :} Par binômes. L'un est le recruteur (HR d'une entreprise sino-camerounaise à Douala), l'autre est le candidat (élève de Terminale A).
\begin{enumerate}
\item \textbf{Préparation (5 min) :} Le candidat remplit oralement son mini-CV (Méthode 5). Le recruteur prépare 4 questions types : \zh{你会说中文吗？} / \zh{你有工作经验吗？} / \zh{你为什么想找这份工作？} / \zh{你的优点是什么？}
\item \textbf{Entretien (3 min) :} Le candidat présente son CV (\zh{这是我的简历，请看。}). Le recruteur pose ses questions. Le candidat répond par des phrases complètes.
\item \textbf{Échange :} Inverser les rôles.
\item \textbf{Retour classe :} 2 binômes jouent la scène devant la classe. L'enseignant note : tons, ordre des mots, classificateurs, fluidité.
\end{enumerate}
\textbf{Phrases utiles pour le recruteur :} \zh{请坐。请介绍一下你自己。你的中文水平怎么样？你会用Office软件吗？我们会通知你。} (\emph{Qǐng zuò. Qǐng jièshào yíxià nǐ zìjǐ. Nǐ de Zhōngwén shuǐpíng zěnmeyàng? Nǐ huì yòng Office ruǎnjiàn ma? Wǒmen huì tōngzhī nǐ.})
\end{experience}

\courssub{Exercices d'entraînement}

\begin{exercice}[Entraînement 1 : Vocabulaire et classificateurs]
Complète les phrases avec le bon classificateur : \zh{份 / 所 / 种 / 封 / 个 / 年}
\begin{enumerate}
\item 这是一\_\_\_\_\_\_简历。
\item 他在一\_\_\_\_\_\_中学学习。
\item 我会说三\_\_\_\_\_\_语言。
\item 我给公司写了一\_\_\_\_\_\_信。
\item 他有两\_\_\_\_\_\_工作经验。
\item 我2005\_\_\_\_\_\_出生。
\end{enumerate}
\end{exercice}

\begin{exercice}[Entraînement 2 : Order des mots / Mots-outils]
Remets les mots dans l'ordre pour former une phrase correcte.
\begin{enumerate}
\item 我 / 在 / 雅温得 / 住
\item 2023年 / 我 / 毕业 / 了
\item 想 / 找 / 工作 / 我 / 一份 / 翻译 / 的
\item 会 / 说 / 法文 / 和 / 中文 / 我
\item 喜欢 / 足球 / 我 / 和 / 音乐
\end{enumerate}
\end{exercice}

\begin{exercice}[Entraînement 3 : Compréhension écrite]
Lis le CV de Zhang Wei et réponds aux questions en chinois.
\zh{张伟，男，2004年10月1日出生，住在杜阿拉。2021年至2024年，他在雅温得一所中学学习。他会说中文、法文，还会用电脑办公。他喜欢看书和跑步。他想找一份导游的工作。}

Questions :
\begin{enumerate}
\item 张伟住在哪里？
\item 他在哪里学习的？几年？
\item 他会说几种语言？
\item 他会用电脑做什么？
\item 他想找什么工作？
\end{enumerate}
\end{exercice}

\begin{exercice}[Entraînement 4 : Expression écrite — Mon CV]
Rédige ton propre mini-CV en chinois (6 phrases minimum) en t'inspirant de la Méthode 5. Utilise ton vrai prénom (transcrit en chinois si possible), ta ville, ton école, tes langues, tes loisirs et ton projet professionnel (métier visé). Écris les caractères, le pinyin et la traduction française.
\end{exercice}

\courssub{Corrigés des exercices}

\begin{corrige}[Exercice 1 : Vocabulaire et classificateurs]
\begin{enumerate}
\item \zh{份} (\emph{fèn}) — document/CV.
\item \zh{所} (\emph{suǒ}) — école/institution.
\item \zh{种} (\emph{zhǒng}) — type/sorte (langues).
\item \zh{封} (\emph{fēng}) — lettre/email.
\item \zh{年} (\emph{nián}) — année (durée/expérience) — \emph{Note : on dit souvent 两年经验, le classificateur est l'unité de temps elle-même}.
\item \zh{年} (\emph{nián}) — pas de classificateur ajouté, l'année est le complément de temps direct.
\end{enumerate}
\end{corrige}

\begin{corrige}[Exercice 2 : Ordre des mots]
\begin{enumerate}
\item \zh{我住在雅温得。} (\emph{Wǒ zhù zài Yǎwēndé.})
\item \zh{我2023年毕业了。} (\emph{Wǒ 2023 nián bìyè le.}) — \zh{了} marque l'accomplissement/relevance actuelle.
\item \zh{我想找一份翻译的工作。} (\emph{Wǒ xiǎng zhǎo yí fèn fānyì de gōngzuò.})
\item \zh{我会说法文和中文。} (\emph{Wǒ huì shuō Fǎwén hé Zhōngwén.})
\item \zh{我喜欢足球和音乐。} (\emph{Wǒ xǐhuan zúqiú hé yīnyuè.})
\end{enumerate}
\end{corrige}

\begin{corrige}[Exercice 3 : Compréhension écrite]
\begin{enumerate}
\item \zh{他住在杜阿拉。} (\emph{Tā zhù zài Dù'ālā.})
\item \zh{他在雅温得一所中学学习，学了三年。} (\emph{Tā zài Yǎwēndé yì suǒ zhōngxué xuéxí, xué le sān nián.}) — 2021 à 2024 = 3 ans.
\item \zh{他会说两种语言：中文和法文。} (\emph{Tā huì shuō liǎng zhǒng yǔyán : Zhōngwén hé Fǎwén.})
\item \zh{他会用电脑办公。} (\emph{Tā huì yòng diànnǎo bàngōng.})
\item \zh{他想找一份导游的工作。} (\emph{Tā xiǎng zhǎo yí fèn dǎoyóu de gōngzuò.})
\end{enumerate}
\end{corrige}

\begin{corrige}[Exercice 4 : Expression écrite — Modèle type]
\textbf{Exemple pour un élève de Douala nommé Paul (保罗 \emph{Bǎoluó}), né en 2006, voulant être ingénieur (工程师 \emph{gōngchéngshī}) :}

\zh{我叫保罗，我是男的。我2006年在杜阿拉出生。我在一所中学学习，现在是毕业班的学生。我会说法文和英文，还会说一点儿中文。我喜欢篮球和看书。我想当工程师。}

\emph{Wǒ jiào Bǎoluó, wǒ shì nán de. Wǒ 2006 nián zài Dù'ālā chūshēng. Wǒ zài yì suǒ zhōngxué xuéxí, xiànzài shì bìyè bān de xuésheng. Wǒ huì shuō Fǎwén hé Yīngwén, hái huì shuō yìdiǎnr Zhōngwén. Wǒ xǐhuan lánqiú hé kànshū. Wǒ xiǎng dāng gōngchéngshī.}

\textbf{Critères de réussite :} 6 phrases + caractères corrects + pinyin avec tons + traduction + respect de l'ordre Sujet-Temps-Lieu-Verbe + classificateurs (\zh{所} école, \zh{份} travail) + \zh{会} compétence + \zh{当} métier.
\end{corrige}

\begin{attention}[Les 5 erreurs qui coûtent des points au Bac]
\begin{enumerate}
\item \textbf{Oublier ou mal placer les tons.} \zh{简历} s'écrit \emph{jiǎnlì} et non \emph{jianli} sans tons ; confondre \zh{买} (\emph{mǎi}, acheter) et \zh{卖} (\emph{mài}, vendre) change le sens de la phrase. Tout pinyin sans ton est sanctionné.
\item \textbf{Confondre les caractères proches (homographes / homophones).} \zh{住} (habiter, clé \zh{亻}) ≠ \zh{注} (noter/verser, clé \zh{氵}) ; \zh{作} (activité intellectuelle/artistique) ≠ \zh{做} (action manuelle/concrète) ; \zh{已} (déjà, fini, 3\textsuperscript{e} trait court) ≠ \zh{己} (soi, 3\textsuperscript{e} trait long à gauche) ≠ \zh{巳} (6\textsuperscript{e} branche terrestre). Relire chaque caractère en nommant sa \cle{clé}.
\item \textbf{Calquer l'ordre français (syntaxe).} On écrit \zh{我住在雅温得} (\emph{Wǒ zhù zài Yǎwēndé}), jamais \zh{我住雅温得在} : le complément de lieu introduit par \zh{在} se place \textbf{juste après le verbe}. De même, le complément de temps précède le verbe : \zh{我2023年毕业}, pas \zh{我毕业2023年}. Pas de préposition « en » devant l'année.
\item \textbf{Négliger les classificateurs (词).} \zh{一份工作} (pas \zh{一个工作}), \zh{一所学校}, \zh{三种语言}, \zh{一封信}, \zh{两年经验}. Un mauvais classificateur est une faute de grammaire comptée à l'examen.
\item \textbf{Mal employer les mots-outils (虚词).} \zh{的} relie un nom à son complément (\zh{我的简历}, \zh{翻译的工作}) ; \zh{了} marque l'action accomplie/le changement d'état (\zh{我学了三年中文}) et ne s'emploie pas pour un fait habituel sans durée ; \zh{过} marque l'expérience vécue (\zh{我去过中国}) ; \zh{会} (savoir faire, compétence apprise) ne se confond pas avec \zh{能} (pouvoir, possibilité physique/circonstancielle) : on dit \zh{我会说中文} (compétence CV), pas \zh{我能说中文} pour présenter ses compétences.
\end{enumerate}
\end{attention}

\newpage

\coursec{Exercices}

\courssub{Je vérifie mes connaissances}

\begin{exercice}[QCM : le vocabulaire du CV]
Entoure la bonne réponse.
\begin{enumerate}
\item \zh{简历} (jiǎnlì) signifie :
\begin{enumerate}
\item une lettre de motivation \item un curriculum vitae \item un diplôme
\end{enumerate}
\item \zh{国籍} (guójí) signifie :
\begin{enumerate}
\item la date de naissance \item la nationalité \item l'adresse
\end{enumerate}
\item \zh{毕业} (bìyè) signifie :
\begin{enumerate}
\item commencer ses études \item redoubler une classe \item terminer ses études, obtenir son diplôme
\end{enumerate}
\item \zh{申请} (shēnqǐng) signifie :
\begin{enumerate}
\item postuler, faire une demande \item refuser une offre \item téléphoner
\end{enumerate}
\end{enumerate}
\end{exercice}

\begin{exercice}[Vrai ou faux ?]
D'après le texte support « le CV de Li Ming », dis si chaque affirmation est vraie ou fausse et justifie ta réponse en citant la phrase du texte (en caractères).
\begin{enumerate}
\item Li Ming est né à Douala.
\item Li Ming a étudié le chinois à l'université.
\item Li Ming n'a aucune expérience professionnelle.
\item Li Ming parle trois langues.
\end{enumerate}
\end{exercice}

\begin{exercice}[Complète avec le bon mot]
Complète chaque phrase avec le mot du glossaire qui convient : \zh{经验} (jīngyàn), \zh{专业} (zhuānyè), \zh{爱好} (àihào), \zh{出生} (chūshēng), \zh{负责} (fùzé).
\begin{enumerate}
\item 我1998年\zh{在雅温得}\_\_\_\_\_\_。(Je suis né à Yaoundé en 1998.)
\item 我的\_\_\_\_\_\_是计算机科学。(Ma spécialité est l'informatique.)
\item 我有两年的工作\_\_\_\_\_\_。(J'ai deux ans d'expérience professionnelle.)
\item 我的\_\_\_\_\_\_是踢足球和听音乐。(Mes loisirs sont le football et la musique.)
\item 他\_\_\_\_\_\_公司的销售工作。(Il est responsable des ventes de l'entreprise.)
\end{enumerate}
\end{exercice}

\begin{exercice}[Associe caractère, pinyin et traduction]
Relie chaque caractère ou mot à son pinyin, puis à sa traduction française.
\begin{center}
\begin{tabularx}{\linewidth}{|X|X|X|}\hline
\rowcolor{popblueL} Caractères & Pinyin & Traduction \\ \hline
简历 & chūshēng & être responsable de \\ \hline
出生 & guójí & naître \\ \hline
国籍 & fùzé & curriculum vitae \\ \hline
负责 & shēnqǐng & nationalité \\ \hline
申请 & jiǎnlì & postuler \\ \hline
\end{tabularx}
\end{center}
\end{exercice}

\courssub{Je m'entraîne}

\begin{exercice}[Traduis en français]
Traduis les phrases suivantes en français.
\begin{enumerate}
\item \zh{李明2020年毕业于雅温得第二大学。}\\ \emph{Lǐ Míng 2020 nián bìyè yú Yǎwēndé Dì-èr Dàxué.}
\item \zh{他有一年的工作经验。}\\ \emph{Tā yǒu yì nián de gōngzuò jīngyàn.}
\item \zh{我想申请这个工作。}\\ \emph{Wǒ xiǎng shēnqǐng zhège gōngzuò.}
\item \zh{我的爱好是看书和打篮球。}\\ \emph{Wǒ de àihào shì kànshū hé dǎ lánqiú.}
\end{enumerate}
\end{exercice}

\begin{exercice}[Traduis en chinois]
Traduis les phrases suivantes en chinois (caractères et pinyin).
\begin{enumerate}
\item J'ai obtenu mon diplôme en 2024.
\item Je parle français, anglais et chinois.
\item Je voudrais postuler à cet emploi.
\item Je suis né(e) à Buea en 2006.
\end{enumerate}
\end{exercice}

\begin{exercice}[CV à trous]
Complète ce mini-CV avec les mots suivants : \zh{毕业} (bìyè), \zh{专业} (zhuānyè), \zh{经验} (jīngyàn), \zh{技能} (jìnéng), \zh{爱好} (àihào).
\begin{textebox}[Un CV à compléter]
姓名：王芳 \emph{(xìngmíng : Wáng Fāng)} — Nom : Wang Fang\\
国籍：喀麦隆 \emph{(guójí : Kāmàilóng)} — Nationalité : camerounaise\\
2024年\_\_\_\_\_\_于杜阿拉大学，\_\_\_\_\_\_是经济。\\
\emph{2024 nián \_\_\_\_\_\_ yú Dù'ālā Dàxué, \_\_\_\_\_\_ shì jīngjì.}\\
(Diplômée en 2024 de l'université de Douala, spécialité : économie.)\\
工作\_\_\_\_\_\_：两年。\\
\emph{gōngzuò \_\_\_\_\_\_ : liǎng nián.} — Expérience professionnelle : deux ans.\\
\_\_\_\_\_\_：电脑、法语、英语。\\
\emph{\_\_\_\_\_\_ : diànnǎo, Fǎyǔ, Yīngyǔ.} — Compétences : informatique, français, anglais.\\
\_\_\_\_\_\_：唱歌、旅游。\\
\emph{\_\_\_\_\_\_ : chànggē, lǚyóu.} — Loisirs : chanter, voyager.
\end{textebox}
\end{exercice}

\begin{exercice}[Remets les mots dans l'ordre]
Remets les mots dans le bon ordre pour faire une phrase correcte, puis traduis-la.
\begin{enumerate}
\item 于 / 毕业 / 我 / 大学 / 2023年 / 雅温得
\item 经验 / 工作 / 有 / 三年 / 他 / 的
\item 申请 / 想 / 这个 / 我 / 职位
\item 是 / 爱好 / 我的 / 足球 / 和 / 音乐
\end{enumerate}
\end{exercice}

\courssub{Je me prépare au Bac}

\begin{exercice}[Compréhension de l'écrit : le CV de Zhang Wei]
Lis le texte suivant, puis réponds aux questions en français.
\begin{textebox}[张伟的简历]
姓名：张伟 \emph{(xìngmíng : Zhāng Wěi)} — Nom : Zhang Wei\\
出生日期：1999年5月12日 \emph{(chūshēng rìqī : 1999 nián 5 yuè 12 rì)} — Date de naissance : le 12 mai 1999\\
国籍：中国 \emph{(guójí : Zhōngguó)} — Nationalité : chinoise\\
地址：雅温得，喀麦隆 \emph{(dìzhǐ : Yǎwēndé, Kāmàilóng)} — Adresse : Yaoundé, Cameroun\\
我2021年毕业于北京大学，专业是汉语国际教育。\\
\emph{Wǒ 2021 nián bìyè yú Běijīng Dàxué, zhuānyè shì Hànyǔ guójì jiàoyù.}\\
J'ai obtenu mon diplôme de l'université de Pékin en 2021, spécialité : enseignement international du chinois.\\
我有三年的工作经验。我在一所中学负责中文课。\\
\emph{Wǒ yǒu sān nián de gōngzuò jīngyàn. Wǒ zài yì suǒ zhōngxué fùzé Zhōngwén kè.}\\
J'ai trois ans d'expérience. Je suis responsable des cours de chinois dans un lycée.\\
我会说中文、英文和法文。我的爱好是看书和打乒乓球。\\
\emph{Wǒ huì shuō Zhōngwén, Yīngwén hé Fǎwén. Wǒ de àihào shì kànshū hé dǎ pīngpāngqiú.}\\
Je sais parler chinois, anglais et français. Mes loisirs sont la lecture et le tennis de table.
\end{textebox}
Questions de compréhension (réponds en français) :
\begin{enumerate}
\item Quel est le nom du candidat et quelle est sa nationalité ?
\item Quelle est sa date de naissance ?
\item Dans quelle université a-t-il étudié et quelle était sa spécialité ?
\item Quelle est son expérience professionnelle ?
\item Quelles langues parle-t-il ?
\end{enumerate}
Questions de langue :
\begin{enumerate}
\item Relève dans le texte la structure utilisée avec \zh{毕业于} et recopie la phrase complète.
\item Trouve dans le texte deux mots formés avec la clé « parole / langue » ou liés au domaine des études.
\end{enumerate}
\end{exercice}

\begin{exercice}[Thème et version]
\textbf{Thème.} Traduis en chinois (caractères et pinyin) :
\begin{enumerate}
\item Je suis né à Douala en 2005.
\item J'ai obtenu mon baccalauréat en 2024.
\item Je parle français, anglais et un peu de chinois.
\item Je voudrais postuler à cet emploi dans votre entreprise.
\end{enumerate}
\textbf{Version.} Traduis en français :
\begin{enumerate}
\item \zh{她2022年毕业于杜阿拉大学，专业是经济。}\\ \emph{Tā 2022 nián bìyè yú Dù'ālā Dàxué, zhuānyè shì jīngjì.}
\item \zh{她有两年的工作经验，会说法语和英语。}\\ \emph{Tā yǒu liǎng nián de gōngzuò jīngyàn, huì shuō Fǎyǔ hé Yīngyǔ.}
\end{enumerate}
\end{exercice}

\begin{exercice}[Expression écrite : mon CV]
Tu es élève de Terminale et tu veux postuler à un petit emploi d'été dans une entreprise de Yaoundé. Rédige ton CV en chinois.
Consignes :
\begin{itemize}
\item Écris entre 60 et 80 caractères chinois.
\item Indique : ton nom, ta date de naissance, ta nationalité, ta formation, tes langues, tes compétences et tes loisirs.
\item Utilise au moins cinq mots du glossaire de la leçon (\zh{简历}, \zh{出生}, \zh{国籍}, \zh{毕业}, \zh{专业}, \zh{经验}, \zh{技能}, \zh{爱好}, \zh{申请}…).
\item Termine par une phrase de candidature avec \zh{申请}.
\end{itemize}
\end{exercice}

\newpage

\coursec{Corrigés des exercices}

\begin{corrige}[Exercice 1 — QCM : le vocabulaire du CV]
\begin{enumerate}
\item \zh{简历} (jiǎnlì) = \textbf{b. un curriculum vitae}. \zh{简历} désigne le CV ; la lettre de motivation se dit \zh{求职信} (qiúzhíxìn) et le diplôme \zh{文凭} (wénpíng).
\item \zh{国籍} (guójí) = \textbf{b. la nationalité}. La date de naissance se dit \zh{出生日期} (chūshēng rìqī) et l'adresse \zh{地址} (dìzhǐ).
\item \zh{毕业} (bìyè) = \textbf{c. terminer ses études, obtenir son diplôme}. Attention : \zh{毕业} est un verbe séparable ; on dit \zh{毕业于 + école} (diplômé de…).
\item \zh{申请} (shēnqǐng) = \textbf{a. postuler, faire une demande}. Exemple : \zh{申请工作} (shēnqǐng gōngzuò), postuler à un emploi.
\end{enumerate}
\textbf{Point de vigilance :} ne pas confondre \zh{简历} (jiǎnlì, CV) et \zh{经历} (jīnglì, expérience vécue, parcours) : les deux mots se terminent par \zh{历} mais n'ont ni le même sens ni la même prononciation du premier caractère.
\end{corrige}

\begin{corrige}[Exercice 2 — Vrai ou faux ?]
\begin{enumerate}
\item \textbf{Vrai.} Le texte dit : \zh{我出生于杜阿拉。} \emph{(Wǒ chūshēng yú Dù'ālā.)} — « Je suis né à Douala. »
\item \textbf{Vrai.} Le texte dit : \zh{我毕业于雅温得第二大学，专业是汉语。} \emph{(Wǒ bìyè yú Yǎwēndé Dì-èr Dàxué, zhuānyè shì Hànyǔ.)} — « Diplômé de l'université de Yaoundé II, spécialité : chinois. »
\item \textbf{Faux.} Li Ming a une expérience professionnelle ; le texte dit : \zh{我有一年的工作经验。} \emph{(Wǒ yǒu yì nián de gōngzuò jīngyàn.)} — « J'ai un an d'expérience professionnelle. »
\item \textbf{Vrai.} Le texte dit : \zh{我会说法语、英语和汉语。} \emph{(Wǒ huì shuō Fǎyǔ, Yīngyǔ hé Hànyǔ.)} — « Je sais parler français, anglais et chinois », soit trois langues.
\end{enumerate}
\textbf{Méthode :} pour justifier, cite toujours la phrase exacte du texte en caractères, sans la modifier ; c'est la preuve attendue au Bac.
\end{corrige}

\begin{corrige}[Exercice 3 — Complète avec le bon mot]
\begin{enumerate}
\item 我1998年\zh{在雅温得出生}。 \emph{Wǒ 1998 nián zài Yǎwēndé chūshēng.} — Je suis né à Yaoundé en 1998. (\zh{出生} = naître ; le lieu est introduit par \zh{在}.)
\item 我的\zh{专业}是计算机科学。 \emph{Wǒ de zhuānyè shì jìsuànjī kēxué.} — Ma spécialité est l'informatique.
\item 我有两年的工作\zh{经验}。 \emph{Wǒ yǒu liǎng nián de gōngzuò jīngyàn.} — J'ai deux ans d'expérience professionnelle. (Collocation fixe : \zh{工作经验}.)
\item 我的\zh{爱好}是踢足球和听音乐。 \emph{Wǒ de àihào shì tī zúqiú hé tīng yīnyuè.} — Mes loisirs sont le football et la musique.
\item 他\zh{负责}公司的销售工作。 \emph{Tā fùzé gōngsī de xiāoshòu gōngzuò.} — Il est responsable des ventes de l'entreprise.
\end{enumerate}
\textbf{Point de vigilance :} \zh{爱好} (àihào) se prononce avec le quatrième ton sur \zh{好} quand il signifie « loisir » ; \zh{好} au troisième ton (hǎo) signifie « bon ».
\end{corrige}

\begin{corrige}[Exercice 4 — Associe caractère, pinyin et traduction]
\begin{center}
\begin{tabularx}{\linewidth}{|X|X|X|}\hline
\rowcolor{popblueL} Caractères & Pinyin & Traduction \\ \hline
简历 & jiǎnlì & curriculum vitae \\ \hline
出生 & chūshēng & naître \\ \hline
国籍 & guójí & nationalité \\ \hline
负责 & fùzé & être responsable de \\ \hline
申请 & shēnqǐng & postuler \\ \hline
\end{tabularx}
\end{center}
\textbf{Astuce mémoire :} \zh{简} contient la clé du bambou \zh{⺮} (les premiers documents étaient écrits sur des lamelles de bambou) ; \zh{籍} (registre) contient la même clé. \zh{申} dans \zh{申请} évoque la demande formelle.
\end{corrige}

\begin{corrige}[Exercice 5 — Traduis en français]
\begin{enumerate}
\item \zh{李明2020年毕业于雅温得第二大学。} — « Li Ming a obtenu son diplôme de l'université de Yaoundé II en 2020. » Structure : sujet + date + \zh{毕业于} + établissement. En français, la date se place en fin de phrase ; en chinois, elle se place avant le verbe.
\item \zh{他有一年的工作经验。} — « Il a un an d'expérience professionnelle. » \zh{一年} devient \zh{一年} avec \zh{一} prononcé au quatrième ton devant un deuxième ton (yì nián).
\item \zh{我想申请这个工作。} — « Je voudrais postuler à cet emploi. » \zh{想} (xiǎng) exprime ici le souhait poli ; \zh{这个} (zhège) = ce/cette + classificateur \zh{个}.
\item \zh{我的爱好是看书和打篮球。} — « Mes loisirs sont la lecture et le basket-ball. » Structure : \zh{我的爱好是 + A + 和 + B}.
\end{enumerate}
\end{corrige}

\begin{corrige}[Exercice 6 — Traduis en chinois]
\begin{enumerate}
\item \zh{我2024年毕业了。} \emph{Wǒ 2024 nián bìyè le.} — J'ai obtenu mon diplôme en 2024. (\zh{了} marque l'accomplissement ; la date précède le verbe.)
\item \zh{我会说法语、英语和汉语。} \emph{Wǒ huì shuō Fǎyǔ, Yīngyǔ hé Hànyǔ.} — Je parle français, anglais et chinois. (Énumération avec la virgule chinoise \zh{、} ; \zh{和} avant le dernier élément.)
\item \zh{我想申请这个工作。} \emph{Wǒ xiǎng shēnqǐng zhège gōngzuò.} — Je voudrais postuler à cet emploi.
\item \zh{我2006年出生在布埃亚。} ou \zh{我2006年在布埃亚出生。} \emph{Wǒ 2006 nián zài Bù'āiyà chūshēng.} — Je suis né(e) à Buea en 2006. (Le lieu introduit par \zh{在} se place avant le verbe \zh{出生}.)
\end{enumerate}
\textbf{Point de vigilance :} en chinois, l'ordre est toujours Sujet + (date) + (lieu avec \zh{在}) + verbe. Ne jamais placer le lieu ou la date après le verbe comme en français.
\end{corrige}

\begin{corrige}[Exercice 7 — CV à trous]
\begin{textebox}[Corrigé du CV de Wang Fang]
姓名：王芳 \emph{(xìngmíng : Wáng Fāng)} — Nom : Wang Fang\\
国籍：喀麦隆 \emph{(guójí : Kāmàilóng)} — Nationalité : camerounaise\\
2024年\zh{毕业}于杜阿拉大学，\zh{专业}是经济。\\
\emph{2024 nián bìyè yú Dù'ālā Dàxué, zhuānyè shì jīngjì.}\\
Diplômée en 2024 de l'université de Douala, spécialité : économie.\\
工作\zh{经验}：两年。\\
\emph{gōngzuò jīngyàn : liǎng nián.} — Expérience professionnelle : deux ans.\\
\zh{技能}：电脑、法语、英语。\\
\emph{jìnéng : diànnǎo, Fǎyǔ, Yīngyǔ.} — Compétences : informatique, français, anglais.\\
\zh{爱好}：唱歌、旅游。\\
\emph{àihào : chànggē, lǚyóu.} — Loisirs : chanter, voyager.
\end{textebox}
\textbf{Vérification :} \zh{毕业于 + université} (structure fixe) ; \zh{工作经验} (collocation) ; \zh{技能} précède la liste des compétences ; \zh{爱好} précède la liste des loisirs.
\end{corrige}

\begin{corrige}[Exercice 8 — Remets les mots dans l'ordre]
\begin{enumerate}
\item \zh{我2023年毕业于雅温得大学。} \emph{Wǒ 2023 nián bìyè yú Yǎwēndé Dàxué.} — « J'ai obtenu mon diplôme de l'université de Yaoundé en 2023. » Ordre : sujet + date + \zh{毕业于} + lieu d'études.
\item \zh{他有三年的工作经验。} \emph{Tā yǒu sān nián de gōngzuò jīngyàn.} — « Il a trois ans d'expérience professionnelle. » La durée \zh{三年} précède \zh{的} + nom.
\item \zh{我想申请这个职位。} \emph{Wǒ xiǎng shēnqǐng zhège zhíwèi.} — « Je voudrais postuler à ce poste. »
\item \zh{我的爱好是足球和音乐。} \emph{Wǒ de àihào shì zúqiú hé yīnyuè.} — « Mes loisirs sont le football et la musique. »
\end{enumerate}
\textbf{Point de vigilance :} \zh{毕业于} est inséparable dans cet ordre ; on ne dit jamais \zh{毕业雅温得大学于}. La date se place avant le verbe, jamais à la fin.
\end{corrige}

\begin{corrige}[Exercice 9 — Compréhension de l'écrit : le CV de Zhang Wei]
\textbf{Questions de compréhension :}
\begin{enumerate}
\item Le candidat s'appelle Zhang Wei (\zh{张伟}) et il est de nationalité chinoise (\zh{国籍：中国}).
\item Il est né le 12 mai 1999 (\zh{出生日期：1999年5月12日}).
\item Il a étudié à l'université de Pékin (\zh{北京大学}), avec pour spécialité l'enseignement international du chinois (\zh{专业是汉语国际教育}).
\item Il a trois ans d'expérience professionnelle : il est responsable des cours de chinois dans un lycée (\zh{我有三年的工作经验。我在一所中学负责中文课。}).
\item Il parle chinois, anglais et français (\zh{我会说中文、英文和法文。}).
\end{enumerate}
\textbf{Questions de langue :}
\begin{enumerate}
\item Structure avec \zh{毕业于} : \zh{我2021年毕业于北京大学，专业是汉语国际教育。} Schéma : sujet + date + \zh{毕业于} + établissement.
\item Mots liés aux études ou à la langue : \zh{汉语国际教育} (enseignement international du chinois), \zh{大学} (université), \zh{中学} (lycée/collège), \zh{课} (cours — contient la clé de la parole \zh{讠}). Deux réponses suffisent.
\end{enumerate}
\textbf{Barème indicatif (sur 20) :} compréhension 5 $\times$ 2 = 10 points ; questions de langue 2 $\times$ 3 = 6 points ; qualité de la rédaction en français 4 points.
\textbf{Points de vigilance :} répondre en phrases complètes ; ne pas confondre \zh{出生日期} (date de naissance) et \zh{国籍} (nationalité) ; \zh{一所中学} utilise le classificateur \zh{所} propre aux établissements scolaires.
\end{corrige}

\begin{corrige}[Exercice 10 — Thème et version]
\textbf{Thème :}
\begin{enumerate}
\item \zh{我2005年出生在杜阿拉。} \emph{Wǒ 2005 nián chūshēng zài Dù'ālā.} (ou \zh{我2005年在杜阿拉出生。} \emph{Wǒ 2005 nián zài Dù'ālā chūshēng.})
\item \zh{我2024年获得了高中毕业证书。} \emph{Wǒ 2024 nián huòdé le gāozhōng bìyè zhèngshū.} — Formulation simple acceptée : \zh{我2024年高中毕业了。} \emph{Wǒ 2024 nián gāozhōng bìyè le.}
\item \zh{我会说法语、英语和一点儿汉语。} \emph{Wǒ huì shuō Fǎyǔ, Yīngyǔ hé yìdiǎnr Hànyǔ.}
\item \zh{我想申请贵公司的这个工作。} \emph{Wǒ xiǎng shēnqǐng guì gōngsī de zhège gōngzuò.} (\zh{贵} = « votre » honorifique, très apprécié dans une candidature.)
\end{enumerate}
\textbf{Version :}
\begin{enumerate}
\item \zh{她2022年毕业于杜阿拉大学，专业是经济。} — « Elle a obtenu son diplôme de l'université de Douala en 2022, avec une spécialité en économie. »
\item \zh{她有两年的工作经验，会说法语和英语。} — « Elle a deux ans d'expérience professionnelle et sait parler français et anglais. »
\end{enumerate}
\textbf{Barème indicatif (sur 20) :} thème 4 $\times$ 2,5 = 10 points ; version 2 $\times$ 4 = 8 points ; exactitude des tons du pinyin 2 points.
\textbf{Points de vigilance :} date avant le verbe ; \zh{一点儿} pour « un peu de » ; \zh{了} pour marquer l'accomplissement du diplôme ; en version, rendre \zh{会} par « savoir » (capacité apprise), pas par « pouvoir ».
\end{corrige}

\begin{corrige}[Exercice 11 — Expression écrite : mon CV]
\textbf{Proposition de rédaction modèle :}
\begin{textebox}[Modèle de CV — 我的简历]
\zh{姓名：玛丽} \emph{xìngmíng : Mǎlì} — Nom : Marie\\
\zh{出生日期：2006年3月5日} \emph{chūshēng rìqī : 2006 nián 3 yuè 5 rì} — Née le 5 mars 2006\\
\zh{国籍：喀麦隆} \emph{guójí : Kāmàilóng} — Nationalité : camerounaise\\
\zh{我是雅温得一所中学的高中生，明年毕业。}\\
\emph{Wǒ shì Yǎwēndé yì suǒ zhōngxué de gāozhōngshēng, míngnián bìyè.}\\
Je suis élève de Terminale dans un lycée de Yaoundé, diplômée l'an prochain.\\
\zh{我会说法语、英语和一点儿汉语。我的技能是电脑。我的爱好是踢足球和听音乐。}\\
\emph{Wǒ huì shuō Fǎyǔ, Yīngyǔ hé yìdiǎnr Hànyǔ. Wǒ de jìnéng shì diànnǎo. Wǒ de àihào shì tī zúqiú hé tīng yīnyuè.}\\
Je parle français, anglais et un peu de chinois. Ma compétence : l'informatique. Mes loisirs : le football et la musique.\\
\zh{我想申请贵公司的暑期工作。}\\
\emph{Wǒ xiǎng shēnqǐng guì gōngsī de shǔqī gōngzuò.}\\
Je voudrais postuler à un emploi d'été dans votre entreprise.
\end{textebox}
\textbf{Barème indicatif (sur 20) :} respect des consignes et des éléments demandés (nom, naissance, nationalité, formation, langues, compétences, loisirs) 8 points ; emploi d'au moins cinq mots du glossaire 4 points ; correction grammaticale et ordre des mots 4 points ; exactitude des caractères et du pinyin 2 points ; phrase finale avec \zh{申请} 2 points.
\textbf{Points de vigilance :}
\begin{itemize}
\item La date de naissance suit l'ordre année + mois + jour : \zh{2006年3月5日}.
\item Le lieu avec \zh{在} se place avant le verbe : \zh{在雅温得出生}.
\item Ne pas oublier \zh{的} dans \zh{我的爱好} et \zh{我的技能}.
\item \zh{一点儿} (un peu de) se place avant le nom : \zh{一点儿汉语}.
\item Terminer impérativement par une phrase de candidature contenant \zh{申请}.
\end{itemize}
\end{corrige}

\newpage

\coursec{Fiche bilan}

\courssub{Fiche Bilan — Leçon 38 : Rédaction d'un CV}

Cette fiche bilan synthétise l'essentiel de la leçon : le lexique des rubriques d'un CV chinois, les structures grammaticales et collocations indispensables, ainsi qu'une méthode de rédaction pas à pas. Elle sert de support de révision avant l'examen.

\begin{popfigure}
\begin{tikzpicture}[
  mindmap,
  grow cyclic,
  every node/.style={concept, circular drop shadow, font=\small\bfseries, text=white, execute at begin node=\strut},
  root concept/.append style={concept color=popblue, font=\large\bfseries, minimum size=2.8cm},
  level 1/.append style={level distance=4.2cm, sibling angle=60, minimum size=2.2cm},
  text width=2.2cm, align=center
]
\node[root concept, concept color=popblue, align=center]{Rédaction\\d'un CV\\ \zh{简历} (jiǎnlì)}
  child[concept color=poporange] { node[concept, align=center]{Vocabulaire\\clé\\ \zh{核心词汇}} }
  child[concept color=popgreen] { node[concept, align=center]{Structure\\type\\ \zh{标准结构}} }
  child[concept color=poppurple] { node[concept, align=center]{Grammaire\\et collocations\\ \zh{语法搭配}} }
  child[concept color=poppink] { node[concept, align=center]{Caractères\\à maîtriser\\ \zh{重点汉字}} }
  child[concept color=popgold] { node[concept, align=center]{Différences\\Cameroun/Chine\\ \zh{中喀差异}} }
  child[concept color=popblue] { node[concept, align=center]{Conseils\\de rédaction\\ \zh{写作技巧}} };
\end{tikzpicture}
\legende{Carte mentale récapitulative de la leçon 38 : les six piliers pour réussir son CV en chinois.}
\end{popfigure}

\courssub{Lexique clé — 核心词汇}

\begin{center}
\begin{tabularx}{\linewidth}{|X|X|X|X|}
\hline \rowcolor{popblueL}
\textbf{Caractères} & \textbf{Pinyin} & \textbf{Nature} & \textbf{Traduction} \\ \hline
简历 & jiǎnlì & n. & CV, curriculum vitæ \\ \hline
个人信息 & gèrén xìnxī & n. & informations personnelles \\ \hline
求职意向 & qiúzhí yìxiàng & n. & poste visé, objectif professionnel \\ \hline
教育背景 & jiàoyù bèijǐng & n. & formation, parcours scolaire \\ \hline
工作经历 & gōngzuò jīnglì & n. & expérience professionnelle \\ \hline
实习经历 & shíxí jīnglì & n. & expérience de stage \\ \hline
技能特长 & jìnéng tècháng & n. & compétences et atouts \\ \hline
证书 & zhèngshū & n. & certificat, diplôme \\ \hline
熟练掌握 & shúliàn zhǎngwò & v. & maîtriser couramment \\ \hline
负责 & fùzé & v. & être responsable de, gérer \\ \hline
实习 & shíxí & v./n. & faire un stage ; stage \\ \hline
应届毕业生 & yìngjiè bìyèshēng & n. & jeune diplômé, étudiant en dernière année \\ \hline
自我评价 & zìwǒ píngjià & n. & auto-évaluation, présentation de soi \\ \hline
联系电话 & liánxì diànhuà & n. & téléphone de contact \\ \hline
\end{tabularx}
\end{center}

\begin{attention}[Faux amis et pièges de vocabulaire]
\begin{itemize}
\item \zh{经历} (jīnglì, « expérience vécue ») et \zh{经验} (jīngyàn, « expérience acquise, savoir-faire ») ne sont pas interchangeables : on dit \zh{工作经历} (expérience professionnelle) mais \zh{工作经验} (savoir-faire professionnel). Dans un CV, la rubrique standard est \zh{工作经历}.
\item \zh{简历} (jiǎnlì) se dit uniquement du CV ; ne pas confondre avec \zh{历史} (lìshǐ, « histoire »).
\item Classificateurs : \zh{一份简历} (yí fèn jiǎnlì, « un CV »), \zh{一段经历} (yí duàn jīnglì, « une expérience »), \zh{一份工作} (yí fèn gōngzuò, « un emploi »).
\end{itemize}
\end{attention}

\courssub{Structures grammaticales et collocations types pour le CV}

\begin{center}
\begin{tabularx}{\linewidth}{|X|X|X|}
\hline \rowcolor{popblueL}
\textbf{Structure / Fonction} & \textbf{Exemple (caractères + pinyin)} & \textbf{Traduction} \\ \hline
\zh{毕业于 + école} (formation) & \zh{我2023年毕业于雅温得第一大学。} (Wǒ èr líng èr sān nián bìyè yú Yǎwēndé Dì-yī Dàxué.) & J'ai été diplômé de l'Université de Yaoundé I en 2023. \\ \hline
\zh{主修 + spécialité} & \zh{我主修计算机科学。} (Wǒ zhǔxiū jìsuànjī kēxué.) & Ma spécialité est l'informatique. \\ \hline
\zh{在 + lieu/entreprise + 工作/实习过} (expérience) & \zh{我在一家中国公司实习过。} (Wǒ zài yì jiā Zhōngguó gōngsī shíxí guò.) & J'ai fait un stage dans une entreprise chinoise. \\ \hline
\zh{负责 + tâche} (responsabilité) & \zh{我负责市场调研工作。} (Wǒ fùzé shìchǎng diàoyán gōngzuò.) & J'étais responsable des études de marché. \\ \hline
\zh{熟练使用 / 精通 + outil/langue} (compétence) & \zh{我熟练使用电脑，精通英语。} (Wǒ shúliàn shǐyòng diànnǎo, jīngtōng Yīngyǔ.) & J'utilise couramment l'ordinateur et je maîtrise l'anglais. \\ \hline
\zh{获得 + certificat/prix} (certification) & \zh{我获得了HSK四级证书。} (Wǒ huòdé le HSK sì jí zhèngshū.) & J'ai obtenu le certificat HSK niveau 4. \\ \hline
\zh{具有 + qualité} (atout) & \zh{我具有团队合作精神。} (Wǒ jùyǒu tuánduì hézuò jīngshén.) & J'ai l'esprit d'équipe. \\ \hline
\end{tabularx}
\end{center}

\begin{propriete}[Rappels grammaticaux utiles pour le CV]
\begin{itemize}
\item \zh{毕业于} : la préposition \zh{于} (yú) introduit le complément de lieu ou d'origine après le verbe ; structure : Sujet + (année) + \zh{毕业于} + école. L'année se lit chiffre par chiffre : \zh{2023年} = \emph{èr líng èr sān nián}.
\item \zh{过} (guo) après le verbe marque l'expérience vécue : \zh{实习过} « avoir déjà fait un stage ». C'est la marque idéale pour la rubrique « expérience ».
\item \zh{了} (le) après \zh{获得} marque l'accomplissement d'une action ponctuelle : \zh{获得了证书} « avoir obtenu le diplôme ».
\item Échelle des niveaux de maîtrise : \zh{了解} (liǎojiě, notions) < \zh{掌握} (zhǎngwò, bonne maîtrise) < \zh{熟练} (shúliàn, courant) < \zh{精通} (jīngtōng, expert).
\end{itemize}
\end{propriete}

\begin{methode}[Rédiger son mini-CV en quatre étapes]
\begin{enumerate}
\item \textbf{Analyser l'offre} : repérer les mots-clés de l'annonce (compétences, diplômes exigés) et les réutiliser dans les rubriques \zh{求职意向} et \zh{技能特长}.
\item \textbf{Remplir les six rubriques obligatoires} dans l'ordre : \zh{个人信息} (nom, téléphone, courriel ; la photo est courante en Chine) $\to$ \zh{求职意向} $\to$ \zh{教育背景} (ordre antéchronologique : du plus récent au plus ancien) $\to$ \zh{工作/实习经历} (avec des verbes d'action : \zh{组织} organiser, \zh{完成} accomplir, \zh{负责} gérer, \zh{协助} assister) $\to$ \zh{技能证书} $\to$ \zh{自我评价} (trois lignes maximum).
\item \textbf{Vérifier le vocabulaire technique} : utiliser le glossaire de la leçon et, si besoin, un dictionnaire pour les termes spécifiques (ex. \zh{市场营销} shìchǎng yíngxiāo, « marketing » ; \zh{财务管理} cáiwù guǎnlǐ, « gestion financière »).
\item \textbf{Relire et mettre en page} : une page maximum, présentation claire et alignée ; vérifier chaque caractère, la ponctuation chinoise pleine chasse （，。；：） et les tons des mots-clés pour la présentation orale.
\end{enumerate}
\end{methode}

\begin{savaistu}[Le CV en Chine et au Cameroun]
En Chine, le CV (\zh{简历}) comporte très souvent une photo d'identité, la date de naissance et parfois la situation familiale, ce qui est de moins en moins courant dans les CV francophones. L'ordre antéchronologique (la formation ou l'expérience la plus récente en premier) est la norme dans les deux pays. Au Cameroun, de nombreuses entreprises chinoises (BTP, télécommunications, commerce) recrutent des jeunes diplômés : savoir rédiger un CV en chinois, même bref, est un atout réel pour décrocher un stage ou un emploi à Douala ou à Yaoundé.
\end{savaistu}

\begin{aretenir}[Checklist avant le Bac — Leçon 38]
$\square$ Je connais par cœur le vocabulaire des six rubriques standard d'un CV chinois (de \zh{个人信息} à \zh{自我评价}). \\
$\square$ Je sais écrire correctement les dix caractères clés : \zh{简、历、教、育、工、作、经、验、技、能} (ordre des traits respecté : haut $\to$ bas, gauche $\to$ droite, horizontal avant vertical). \\
$\square$ Je maîtrise les structures \zh{毕业于 + école} et \zh{主修 + spécialité} pour la rubrique « Formation ». \\
$\square$ J'utilise des verbes d'action précis pour décrire mes expériences (\zh{负责、组织、完成、协助、提高}) et non des phrases vagues. \\
$\square$ Je sais exprimer mes niveaux de langue et d'outil : \zh{精通} (expert) > \zh{熟练} (courant) > \zh{了解} (notions). \\
$\square$ Je connais les différences culturelles majeures : photo fréquente en Chine, mention de l'âge possible, format antéchronologique strict. \\
$\square$ Je suis capable de présenter mon CV à l'oral en une minute (\zh{自我介绍} structuré : nom $\to$ formation $\to$ expérience clé $\to$ atout principal $\to$ motivation). \\
$\square$ Je fais attention aux caractères visuellement proches : \zh{经} (jīng, expérience) et \zh{径} (jìng, chemin) ; \zh{历} (lì, calendrier, CV) et \zh{厉} (lì, sévère). \\
$\square$ Je vérifie les classificateurs (\zh{一份简历}, \zh{一段经历}, \zh{一份工作}) et la ponctuation chinoise pleine chasse. \\
$\square$ Je sais adapter mon CV à une annonce chinoise (mots-clés, format, photo) pour un stage ou un emploi en Chine ou dans une entreprise chinoise au Cameroun.
\end{aretenir}

\findecours

\end{document}
